\input{preamble}

% OK, start here.
%
\begin{document}

\title{Complexes dualisants}


\maketitle

\phantomsection
\label{section-phantom}

\tableofcontents

\section{Introduction}
\label{section-introduction}

\noindent
Dans ce chapitre, nous étudions les complexes dualisants en algèbre commutative.
On pourra consulter \cite{RD}.

\medskip\noindent
Nous commençons par les surjections et injections essentielles,
les couvertures projectives, les enveloppes injectives et la dualité
sur les anneaux artiniens, puis étudions les enveloppes injectives
des corps résiduels, ce qui conduit rapidement à une démonstration
de la dualité de Matlis. Voir les sections \ref{section-essential},
\ref{section-injective-modules}, \ref{section-projective-cover}, \ref{section-injective-hull}, \ref{section-artinian} et \ref{section-hull-residue-field},
ainsi que la proposition \ref{proposition-matlis}.

\medskip\noindent
Viennent ensuite trois sections consacrées à la cohomologie locale
dans une grande généralité ; voir les sections \ref{section-bad-local-cohomology},
\ref{section-local-cohomology} et \ref{section-local-cohomology-noetherian}.
Nous en appliquons certains résultats à l'étude de la profondeur
dans la section \ref{section-depth}. Une autre application montre comment,
étant donné un idéal de type fini $I$ d'un anneau $A$,
les objets « $I$-complets » et « de $I$-torsion »
de la catégorie dérivée de $A$ se correspondent par une équivalence ;
voir la section \ref{section-torsion-and-complete}.
Pour approfondir la cohomologie locale, notamment le théorème de finitude
(qui repose sur la dualité locale -- voir ci-dessous), on se reportera à
Cohomologie locale, section \ref{local-cohomology-section-introduction}.

\medskip\noindent
L'essentiel de ce chapitre est consacré à la dualité pour un homomorphisme
d'anneaux et aux complexes dualisants. Voir les sections \ref{section-trivial},
\ref{section-base-change-trivial-duality}, \ref{section-dualizing}, \ref{section-dualizing-local}, \ref{section-dimension-function}, \ref{section-local-duality}, \ref{section-dualizing-module},
\ref{section-CM}, \ref{section-gorenstein}, \ref{section-ubiquity-dualizing} et \ref{section-formal-fibres}.
La définition fondamentale est celle d'un complexe dualisant
$\omega_A^\bullet$ sur un anneau noethérien $A$ : c'est un objet
$\omega_A^\bullet \in D^{+}(A)$ dont les modules de cohomologie
$H^i(\omega_A^\bullet)$ sont des $A$-modules de type fini, qui est de dimension injective finie,
et tel que le morphisme
$$
A \longrightarrow R\Hom_A(\omega_A^\bullet, \omega_A^\bullet)
$$
soit un quasi-isomorphisme. Après avoir établi quelques propriétés
élémentaires des complexes dualisants, nous montrons qu'un complexe
dualisant définit une fonction de dimension. Nous démontrons ensuite
le théorème de dualité locale de Grothendieck. Après une brève étude
des modules dualisants et des anneaux de Cohen-Macaulay, nous introduisons
les anneaux de Gorenstein et montrons que beaucoup d'anneaux noethériens
usuels possèdent des complexes dualisants. Dans une dernière section,
nous appliquons ces résultats pour établir une bonne théorie des anneaux
locaux noethériens dont les fibres formelles sont de Gorenstein ou
localement des intersections complètes.

\medskip\noindent
Dans les dernières sections, nous décrivons une construction algébrique
des « foncteurs image inverse extraordinaire » utilisés en géométrie
algébrique, par exemple dans l'ouvrage \cite{RD}. Ce sujet est poursuivi
dans le chapitre sur la dualité pour les schémas. Voir
Dualité pour les schémas, section \ref{duality-section-introduction}.







\section{Surjections et injections essentielles}
\label{section-essential}

\noindent
Nous travaillerons surtout dans des catégories de modules, mais autant
donner la définition dans le cas général.

\begin{definition}
\label{definition-essential}
Soit $\mathcal{A}$ une catégorie abélienne.
\begin{enumerate}
\item Une injection $A \subset B$ de $\mathcal{A}$ est dite {\it essentielle},
ou l'on dit que $B$ est une {\it extension essentielle de} $A$,
si tout sous-objet non nul $B' \subset B$ a une intersection non nulle avec $A$.
\item Une surjection $f : A \to B$ de $\mathcal{A}$ est dite {\it essentielle}
si, pour tout sous-objet propre $A' \subset A$, on a $f(A') \not = B$.
\end{enumerate}
\end{definition}

\noindent
Voici quelques lemmes sur cette notion.

\begin{lemma}
\label{lemma-essential}
Soit $\mathcal{A}$ une catégorie abélienne.
\begin{enumerate}
\item Si $A \subset B$ et $B \subset C$ sont des extensions essentielles,
alors $A \subset C$ est une extension essentielle.
\item Si $A \subset B$ est une extension essentielle et $C \subset B$
un sous-objet, alors $A \cap C \subset C$ est une extension essentielle.
\item Si $A \to B$ et $B \to C$ sont des surjections essentielles,
alors $A \to C$ est une surjection essentielle.
\item Étant donné une surjection essentielle $f : A \to B$ et une surjection
$A \to C$ de noyau $K$, le morphisme $C \to B/f(K)$ est une surjection essentielle.
\end{enumerate}
\end{lemma}

\begin{proof}
Omise.
\end{proof}

\begin{lemma}
\label{lemma-union-essential-extensions}
Soit $R$ un anneau. Soit $M$ un $R$-module. Soit $E = \colim E_i$
une limite inductive filtrante de $R$-modules. Supposons donné
un système compatible d'injections essentielles $M \to E_i$ de $R$-modules.
Alors $M \to E$ est une injection essentielle.
\end{lemma}

\begin{proof}
Cela résulte immédiatement des définitions et de l'exactitude
des limites inductives filtrantes (Algèbre, lemme \ref{algebra-lemma-directed-colimit-exact}).
\end{proof}

\begin{lemma}
\label{lemma-essential-extension}
Soit $R$ un anneau. Soient $M \subset N$ des $R$-modules.
Les conditions suivantes sont équivalentes :
\begin{enumerate}
\item $M \subset N$ est une extension essentielle ;
\item pour tout $x \in N$ non nul, il existe $f \in R$ tel que $fx \in M$
et $fx \not = 0$.
\end{enumerate}
\end{lemma}

\begin{proof}
Supposons (1) et soit $x \in N$ un élément non nul. D'après (1), on a
$Rx \cap M \not = 0$, ce qui entraîne (2).

\medskip\noindent
Supposons (2). Soit $N' \subset N$ un sous-module non nul. Choisissons $x \in N'$
non nul. D'après (2), il existe $f \in R$ tel que $fx \in M$ et $fx \not = 0$.
Ainsi $N' \cap M \not = 0$.
\end{proof}




\section{Modules injectifs}
\label{section-injective-modules}

\noindent
Quelques résultats sur les modules injectifs sur des anneaux.

\begin{lemma}
\label{lemma-product-injectives}
Soit $R$ un anneau. Tout produit de $R$-modules injectifs est injectif.
\end{lemma}

\begin{proof}
C'est un cas particulier de Homologie, lemme \ref{homology-lemma-product-injectives}.
\end{proof}

\begin{lemma}
\label{lemma-injective-flat}
Soit $R \to S$ un homomorphisme plat d'anneaux. Si $E$ est un $S$-module injectif,
alors $E$ est injectif en tant que $R$-module.
\end{lemma}

\begin{proof}
Cela résulte de l'égalité $\Hom_R(M, E) = \Hom_S(M \otimes_R S, E)$
(Algèbre, lemme \ref{algebra-lemma-adjoint-tensor-restrict})
et de l'exactitude du produit tensoriel avec $S$.
\end{proof}

\begin{lemma}
\label{lemma-injective-epimorphism}
Soit $R \to S$ un épimorphisme d'anneaux. Soit $E$ un $S$-module.
Si $E$ est injectif en tant que $R$-module, alors $E$ est un $S$-module injectif.
\end{lemma}

\begin{proof}
Cela résulte de l'égalité $\Hom_R(N, E) = \Hom_S(N, E)$ pour tout $S$-module $N$ ;
voir Algèbre, lemme \ref{algebra-lemma-epimorphism-modules}.
\end{proof}

\begin{lemma}
\label{lemma-hom-injective}
Soit $R \to S$ un homomorphisme d'anneaux. Si $E$ est un $R$-module injectif,
alors $\Hom_R(S, E)$ est un $S$-module injectif.
\end{lemma}

\begin{proof}
Cela résulte de l'égalité $\Hom_S(N, \Hom_R(S, E)) = \Hom_R(N, E)$ donnée par
Algèbre, lemme \ref{algebra-lemma-adjoint-hom-restrict}.
\end{proof}

\begin{lemma}
\label{lemma-essential-extensions-in-injective}
Soit $R$ un anneau. Soit $I$ un $R$-module injectif. Soit $E \subset I$
un sous-module. Les conditions suivantes sont équivalentes :
\begin{enumerate}
\item $E$ est injectif ;
\item pour tout $E \subset E' \subset I$ tel que $E \subset E'$ soit essentielle,
on a $E = E'$.
\end{enumerate}
En particulier, un $R$-module est injectif si et seulement si toute
extension essentielle est triviale.
\end{lemma}

\begin{proof}
La dernière assertion résulte de la première et du fait que la catégorie
des $R$-modules possède suffisamment d'injectifs
(Compléments d'algèbre, section \ref{more-algebra-section-injectives-modules}).

\medskip\noindent
Supposons (1). Soit $E \subset E' \subset I$ comme dans (2).
Alors le morphisme $\text{id}_E : E \to E$ se prolonge en un morphisme $\alpha : E' \to E$.
Le noyau de $\alpha$ est nécessairement nul, car son intersection avec
$E$ est nulle et $E'$ est une extension essentielle. Ainsi $E = E'$.

\medskip\noindent
Supposons (2). Soient $M \subset N$ des $R$-modules et $\varphi : M \to E$
un homomorphisme de $R$-modules. Pour prouver (1), il faut montrer que
$\varphi$ se prolonge en un morphisme $N \to E$. Considérons l'ensemble $\mathcal{S}$
des couples $(M', \varphi')$ où $M \subset M' \subset N$ et où $\varphi' : M' \to E$
est un homomorphisme de $R$-modules coïncidant avec $\varphi$ sur $M$.
Définissons un ordre sur $\mathcal{S}$ par la règle $(M', \varphi') \leq (M'', \varphi'')$
si et seulement si $M' \subset M''$ et $\varphi''|_{M'} = \varphi'$.
Il est clair que l'on peut prendre le maximum d'une partie totalement
ordonnée de $\mathcal{S}$. Le lemme de Zorn permet donc de supposer que $(M, \varphi)$
est un élément maximal.

\medskip\noindent
Choisissons un prolongement $\psi : N \to I$ de $\varphi$ composé avec l'inclusion $E \to I$.
C'est possible puisque $I$ est injectif.
Si $\psi(N) \subset E$, alors $\psi$ est le prolongement cherché.
Si $\psi(N)$ n'est pas contenu dans $E$, alors, d'après (2), l'inclusion
$E \subset E + \psi(N)$ n'est pas essentielle. On peut donc trouver un sous-module non nul
$K \subset E + \psi(N)$ dont l'intersection avec $E$ est $0$.
Cela signifie que $M' = \psi^{-1}(E + K)$ contient strictement $M$.
On peut donc prolonger $\varphi$ à $M'$ au moyen de
$$
M' \xrightarrow{\psi|_{M'}} E + K \to (E + K)/K = E
$$
Cela contredit la maximalité de $(M, \varphi)$.
\end{proof}

\begin{example}
\label{example-reduced-ring-injective}
Soit $R$ un anneau réduit. Soit $\mathfrak p \subset R$ un idéal premier minimal ;
alors $K = R_\mathfrak p$ est un corps
(Algèbre, lemme \ref{algebra-lemma-minimal-prime-reduced-ring}).
Le module $K$ est un $R$-module injectif. En effet, on a
$\Hom_R(M, K) = \Hom_K(M_\mathfrak p, K)$ pour tout $R$-module $M$.
La localisation étant un foncteur exact, ainsi que le passage au dual
sur les $K$-espaces vectoriels, on en déduit que $\Hom_R(-, K)$
est un foncteur exact, c'est-à-dire que $K$ est un $R$-module injectif.
\end{example}

\begin{lemma}
\label{lemma-sum-injective-modules}
Soit $R$ un anneau noethérien. Toute somme directe de modules injectifs
est injective.
\end{lemma}

\begin{proof}
Soit $E_i$ une famille de modules injectifs indexée par un ensemble $I$.
Posons $E = \bigoplus E_i$. Pour montrer que $E$ est injectif, utilisons
Injectifs, lemme \ref{injectives-lemma-criterion-baer}.
Soit donc $\varphi : J \to E$ un homomorphisme de modules d'un idéal de $R$
dans $E$. Puisque $J$ est un $R$-module de type fini
(car $R$ est noethérien), il existe un nombre fini d'éléments $i_1, \ldots, i_r \in I$
tels que $\varphi$ prenne ses valeurs dans $\bigoplus_{j = 1, \ldots, r} E_{i_j}$.
On peut alors prolonger $\varphi$ à valeurs dans $\bigoplus_{j = 1, \ldots, r} E_{i_j}$
en utilisant l'injectivité des modules $E_{i_j}$.
\end{proof}

\begin{lemma}
\label{lemma-localization-injective-modules}
Soit $R$ un anneau noethérien. Soit $S \subset R$ une partie multiplicative.
Si $E$ est un $R$-module injectif, alors $S^{-1}E$ est un $S^{-1}R$-module injectif.
\end{lemma}

\begin{proof}
Puisque $R \to S^{-1}R$ est un épimorphisme d'anneaux, il suffit de montrer que
$S^{-1}E$ est injectif en tant que $R$-module ; voir le lemme \ref{lemma-injective-epimorphism}.
Pour cela, utilisons Injectifs, lemme \ref{injectives-lemma-criterion-baer}.
Soit donc $I \subset R$ un idéal et soit $\varphi : I \to S^{-1} E$ un homomorphisme de $R$-modules.
Puisque $I$ est un $R$-module de présentation finie
(car $R$ est noethérien), il existe $f \in S$ et un homomorphisme
de $R$-modules $I \to E$ tels que $f\varphi$ soit le composé $I \to E \to S^{-1}E$
(Algèbre, lemme \ref{algebra-lemma-hom-from-finitely-presented}).
On peut alors prolonger $I \to E$ en un homomorphisme $R \to E$.
Le composé
$$
R \to E \to S^{-1}E \xrightarrow{f^{-1}} S^{-1}E
$$
est le prolongement cherché de $\varphi$ à $R$.
\end{proof}

\begin{lemma}
\label{lemma-injective-module-divide}
Soit $R$ un anneau noethérien. Soit $I$ un $R$-module injectif.
\begin{enumerate}
\item Soit $f \in R$. Alors $E = \bigcup I[f^n] = I[f^\infty]$
est un sous-module injectif de $I$.
\item Soit $J \subset R$ un idéal. Le sous-module de torsion
par les puissances de $J$, noté $I[J^\infty]$, est un sous-module injectif de $I$.
\end{enumerate}
\end{lemma}

\begin{proof}
Utilisons le lemme \ref{lemma-essential-extensions-in-injective} pour prouver (1).
Supposons que $E \subset E' \subset I$ et que $E'$ soit une extension essentielle de $E$.
Montrons que $E' = E$. Sinon, il existe $x \in E'$ tel que $x \not \in E$.
Posons $J = \{ a \in R \mid ax \in E\}$. Puisque $R$ est noethérien, on peut écrire $J = (g_1, \ldots, g_t)$
pour certains $g_i \in R$. Par définition, $E$ est l'ensemble des éléments
de $I$ annulés par des puissances de $f$ ; on peut donc choisir
des entiers $n_i$ tels que $f^{n_i}g_ix = 0$.
Posons $n = \mathrm{max}\{ n_i \}$. Alors $x' = f^n x$ est un élément de $E'$
qui n'appartient pas à $E$ et qui est annulé par $J$.
Posons $J' = \{ a \in R \mid ax' \in E \}$, de sorte que $J \subset J'$. Réciproquement, on a $a \in J'$
si et seulement si $ax' \in E$, ou encore si et seulement si $f^m a x' = 0$
pour un certain $m \geq 0$. Mais alors $f^m a x' = f^{m + n} a x$ entraîne $ax \in E$, c'est-à-dire $a \in J$.
Ainsi $J = J'$. On a donc $J = J' = \text{Ann}(x')$, d'où $Rx' \cap E  = 0$.
Par conséquent, $E'$ n'est pas une extension essentielle de $E$,
ce qui est contradictoire.

\medskip\noindent
Pour prouver (2), écrivons $J = (f_1, \ldots, f_t)$. Alors $I[J^\infty]$ est égal à
$$
(\ldots((I[f_1^\infty])[f_2^\infty])\ldots)[f_t^\infty]
$$
et le résultat découle de (1) par récurrence.
\end{proof}

\begin{lemma}
\label{lemma-injective-dimension-over-polynomial-ring}
Soit $A$ un anneau noethérien. Soit $E$ un $A$-module injectif.
Alors $E \otimes_A A[x]$ est d'amplitude injective $[0, 1]$
en tant qu'objet de $D(A[x])$. En particulier, $E \otimes_A A[x]$
est de dimension injective finie en tant que $A[x]$-module.
\end{lemma}

\begin{proof}
Écrivons $E[x] = E \otimes_A A[x]$. Considérons la suite exacte courte de $A[x]$-modules
$$
0 \to E[x] \to \Hom_A(A[x], E[x]) \to \Hom_A(A[x], E[x]) \to 0
$$
où le premier morphisme envoie $p \in E[x]$ sur $f \mapsto fp$ et le second envoie
$\varphi$ sur $f \mapsto \varphi(xf) - x\varphi(f)$.
Le second morphisme est surjectif, car $\Hom_A(A[x], E[x]) = \prod_{n \geq 0} E[x]$ en tant que groupe abélien,
et le morphisme envoie $(e_n)$ sur $(e_{n + 1} - xe_n)$, ce qui définit une surjection.
En tant que $A$-module, on a $E[x] \cong \bigoplus_{n \geq 0} E$,
qui est injectif d'après le lemme \ref{lemma-sum-injective-modules}.
Le $A[x]$-module $\Hom_A(A[x], E[x])$ est donc injectif d'après le lemme \ref{lemma-hom-injective},
ce qui achève la démonstration.
\end{proof}



\section{Couvertures projectives}
\label{section-projective-cover}

\noindent
Dans cette section, nous étudions brièvement les couvertures projectives.

\begin{definition}
\label{definition-projective-cover}
Soit $R$ un anneau. Une surjection $P \to M$ de $R$-modules est appelée
{\it couverture projective}, ou parfois {\it enveloppe projective},
si $P$ est un $R$-module projectif et si $P \to M$ est une surjection essentielle.
\end{definition}

\noindent
Les couvertures projectives n'existent pas toujours. Par exemple, si $k$
est un corps et $R = k[x]$ l'anneau de polynômes sur $k$, le module $M = R/(x)$
n'a pas de couverture projective. En effet, pour toute surjection $f : P \to M$
avec $P$ projectif sur $R$, le sous-module propre $(x - 1)P$ se surjecte
sur $M$. Ainsi $f$ n'est pas essentielle.

\begin{lemma}
\label{lemma-projective-cover-unique}
Soit $R$ un anneau et soit $M$ un $R$-module. Si une couverture projective
de $M$ existe, elle est unique à isomorphisme près.
\end{lemma}

\begin{proof}
Soient $P \to M$ et $P' \to M$ des couvertures projectives. Puisque $P$ est
un $R$-module projectif et que $P' \to M$ est surjective, il existe
un homomorphisme de $R$-modules $\alpha : P \to P'$ compatible avec les morphismes vers $M$.
Puisque $P' \to M$ est essentielle, $\alpha$ est surjectif.
Comme $P'$ est un $R$-module projectif, on peut choisir une décomposition
en somme directe $P = \Ker(\alpha) \oplus P'$. Puisque $P' \to M$ est surjective
et que $P \to M$ est essentielle, on en déduit que $\Ker(\alpha)$ est nul, comme voulu.
\end{proof}

\noindent
Voici un exemple où les couvertures projectives existent.

\begin{lemma}
\label{lemma-projective-covers-local}
Soit $(R, \mathfrak m, \kappa)$ un anneau local. Tout $R$-module de type fini possède
une couverture projective.
\end{lemma}

\begin{proof}
Soit $M$ un $R$-module de type fini. Posons $r = \dim_\kappa(M/\mathfrak m M)$.
Choisissons $x_1, \ldots, x_r \in M$ dont les images forment une base de $M/\mathfrak m M$.
Considérons le morphisme $f : R^{\oplus r} \to M$. D'après le lemme de Nakayama, il est
surjectif (Algèbre, lemme \ref{algebra-lemma-NAK}). Si $N \subset R^{\oplus r}$ est un sous-module propre,
alors $N/\mathfrak m N \to \kappa^{\oplus r}$ n'est pas surjectif (encore par le lemme de Nakayama),
donc $N/\mathfrak m N \to M/\mathfrak m M$ n'est pas surjectif.
Ainsi $f$ est une surjection essentielle.
\end{proof}







\section{Enveloppes injectives}
\label{section-injective-hull}

\noindent
Dans cette section, nous étudions brièvement les enveloppes injectives.

\begin{definition}
\label{definition-injective-hull}
Soit $R$ un anneau. Une injection $M \to I$ de $R$-modules est appelée
{\it enveloppe injective} si $I$ est un $R$-module injectif
et si $M \to I$ est une injection essentielle.
\end{definition}

\noindent
Les enveloppes injectives existent toujours.

\begin{lemma}
\label{lemma-injective-hull}
Soit $R$ un anneau. Tout $R$-module possède une enveloppe injective.
\end{lemma}

\begin{proof}
Soit $M$ un $R$-module. D'après Compléments d'algèbre, section \ref{more-algebra-section-injectives-modules},
la catégorie des $R$-modules possède suffisamment d'injectifs.
Choisissons une injection $M \to I$ où $I$ est un $R$-module injectif.
Considérons l'ensemble $\mathcal{S}$ des sous-modules $M \subset E \subset I$
tels que $E$ soit une extension essentielle de $M$.
Ordonnons $\mathcal{S}$ par inclusion. Si $\{E_\alpha\}$ est une partie totalement
ordonnée de $\mathcal{S}$, alors $\bigcup E_\alpha$ est aussi une extension essentielle
de $M$ (lemme \ref{lemma-union-essential-extensions}).
Le lemme de Zorn fournit donc un élément maximal $E \in \mathcal{S}$.
Affirmons que $M \subset E$ est une enveloppe injective, c'est-à-dire
que $E$ est un $R$-module injectif. Cela résulte du lemme \ref{lemma-essential-extensions-in-injective}.
\end{proof}

\begin{lemma}
\label{lemma-injective-hull-unique}
Soit $R$ un anneau. Soient $M$, $N$ des $R$-modules et soient $M \to E$
et $N \to E'$ des enveloppes injectives. Alors :
\begin{enumerate}
\item pour tout homomorphisme de $R$-modules $\varphi : M \to N$, il existe un
homomorphisme de $R$-modules $\psi : E \to E'$ tel que le diagramme
$$
\xymatrix{
M \ar[r] \ar[d]_\varphi & E \ar[d]^\psi \\
N \ar[r] & E'
}
$$
soit commutatif ;
\item si $\varphi$ est injectif, alors $\psi$ est injectif ;
\item si $\varphi$ est une injection essentielle, alors $\psi$ est un isomorphisme ;
\item si $\varphi$ est un isomorphisme, alors $\psi$ est un isomorphisme ;
\item si $M \to I$ est un plongement de $M$ dans un $R$-module injectif,
alors il existe un isomorphisme $I \cong E \oplus I'$ compatible avec les plongements de $M$.
\end{enumerate}
En particulier, l'enveloppe injective $E$ de $M$ est unique à isomorphisme près.
\end{lemma}

\begin{proof}
L'assertion (1) résulte de ce que $E'$ est un $R$-module injectif.
L'assertion (2) résulte de ce que $\Ker(\psi) \cap M = 0$
et que $E$ est une extension essentielle de $M$.
Supposons que $\varphi$ soit une injection essentielle. Alors
$E \cong \psi(E) \subset E'$ d'après (2), d'où $E' = \psi(E) \oplus E''$ puisque $E$ est injectif.
Puisque $E'$ est une extension essentielle de $M$
(lemme \ref{lemma-essential}), on obtient $E'' = 0$.
L'assertion (4) est un cas particulier de (3).
Supposons $M \to I$ comme dans (5).
Choisissons un morphisme $\alpha : E \to I$ prolongeant le morphisme $M \to I$.
Le même raisonnement montre que $\alpha$ est injectif.
Comme précédemment, $\alpha(E)$ est donc facteur direct de $I$.
Cela prouve (5).
\end{proof}

\begin{example}
\label{example-injective-hull-domain}
Soit $R$ un anneau intègre de corps des fractions $K$.
Alors $R \subset K$ est une enveloppe injective de $R$.
En effet, l'exemple \ref{example-reduced-ring-injective} montre que $K$ est un $R$-module injectif,
et le lemme \ref{lemma-essential-extension} montre que $R \subset K$ est une extension essentielle.
\end{example}

\begin{definition}
\label{definition-indecomposable}
Un objet $X$ d'une catégorie additive est dit {\it indécomposable}
s'il est non nul et si $X = Y \oplus Z$ entraîne que $Y = 0$ ou $Z = 0$.
\end{definition}

\begin{lemma}
\label{lemma-indecomposable-injective}
Soit $R$ un anneau. Soit $E$ un $R$-module injectif indécomposable.
Alors :
\begin{enumerate}
\item $E$ est l'enveloppe injective de tout sous-module non nul de $E$ ;
\item l'intersection de deux sous-modules non nuls de $E$ est non nulle ;
\item $\text{End}_R(E)$ est un anneau local non commutatif dont l'idéal maximal
est formé des $\varphi : E \to E$ de noyau non nul ;
\item l'ensemble des diviseurs de zéro sur $E$ est un idéal premier $\mathfrak p$
de $R$, et $E$ est un $R_\mathfrak p$-module injectif.
\end{enumerate}
\end{lemma}

\begin{proof}
L'assertion (1) résulte du lemme \ref{lemma-injective-hull-unique}.
L'assertion (2) découle de (1) et de la définition des enveloppes injectives.

\medskip\noindent
Preuve de (3). Posons $A = \text{End}_R(E)$ et $I = \{\varphi \in A \mid \Ker(\varphi) \not = 0\}$.
L'énoncé signifie que $I$ est un idéal bilatère et que tout
$\varphi \in A$ tel que $\varphi \not \in I$ est inversible.
Supposons que $\varphi$ et $\psi$ ne soient pas injectifs.
Alors $\Ker(\varphi) \cap \Ker(\psi)$ est non nul d'après (2).
Ainsi $\varphi + \psi \in I$. Il en résulte que $I$ est un idéal bilatère.
Si $\varphi \in A$ et $\varphi \not \in I$, alors $E \cong \varphi(E) \subset E$ est un sous-module injectif,
donc $E = \varphi(E)$ puisque $E$ est indécomposable.

\medskip\noindent
Preuve de (4). Considérons l'homomorphisme d'anneaux $R \to A$ et soit $\mathfrak p \subset R$
l'image réciproque de l'idéal maximal $I$. Il est alors clair
que $\mathfrak p$ est un idéal premier et que $R \to A$ se prolonge en $R_\mathfrak p \to A$.
Ainsi $E$ est un $R_\mathfrak p$-module.
Le lemme \ref{lemma-injective-epimorphism} montre que $E$ est injectif en tant que $R_\mathfrak p$-module.
\end{proof}

\begin{lemma}
\label{lemma-injective-hull-indecomposable}
Soit $\mathfrak p \subset R$ un idéal premier d'un anneau $R$.
Soit $E$ l'enveloppe injective de $R/\mathfrak p$. Alors :
\begin{enumerate}
\item $E$ est indécomposable ;
\item $E$ est l'enveloppe injective de $\kappa(\mathfrak p)$ ;
\item $E$ est l'enveloppe injective de $\kappa(\mathfrak p)$ sur l'anneau $R_\mathfrak p$.
\end{enumerate}
\end{lemma}

\begin{proof}
D'après le lemme \ref{lemma-essential-extension}, l'inclusion $R/\mathfrak p \subset \kappa(\mathfrak p)$ est une extension essentielle.
Le lemme \ref{lemma-injective-hull-unique} donne alors (2).
Pour $f \in R$ tel que $f \not \in \mathfrak p$, le morphisme $f : \kappa(\mathfrak p) \to \kappa(\mathfrak p)$ est un isomorphisme ;
le morphisme $f : E \to E$ est donc un isomorphisme, voir le lemme \ref{lemma-injective-hull-unique}.
Ainsi $E$ est un $R_\mathfrak p$-module. Il est injectif en tant que
$R_\mathfrak p$-module d'après le lemme \ref{lemma-injective-epimorphism}.
Enfin, soit $E' \subset E$ un sous-$R$-module injectif non nul.
Alors $J = (R/\mathfrak p) \cap E'$ est non nul. Quitte à remplacer $E'$ par un sous-module,
on peut supposer que $E'$ est l'enveloppe injective de $J$
(voir par exemple le lemme \ref{lemma-injective-hull-unique}).
Observons que $R/\mathfrak p$ est une extension essentielle de $J$,
par exemple d'après le lemme \ref{lemma-essential-extension}. Ainsi $E' \to E$
est un isomorphisme d'après le lemme \ref{lemma-injective-hull-unique} (3).
Par conséquent, $E$ est indécomposable.
\end{proof}

\begin{lemma}
\label{lemma-indecomposable-injective-noetherian}
Soit $R$ un anneau noethérien. Soit $E$ un $R$-module injectif
indécomposable. Il existe alors un idéal premier $\mathfrak p$ de $R$
tel que $E$ soit l'enveloppe injective de $\kappa(\mathfrak p)$.
\end{lemma}

\begin{proof}
Soit $\mathfrak p$ l'idéal premier obtenu dans le lemme \ref{lemma-indecomposable-injective}.
Écrivons $\mathfrak p = (f_1, \ldots, f_r)$.
Choisissons un élément non nul $x \in \bigcap \Ker(f_i : E \to E)$, voir le lemme \ref{lemma-indecomposable-injective}.
Alors $(R_\mathfrak p)x$ est un module isomorphe à $\kappa(\mathfrak p)$ contenu dans $E$.
On conclut par le lemme \ref{lemma-indecomposable-injective}.
\end{proof}

\begin{proposition}[Structure des modules injectifs sur les anneaux noethériens]
\label{proposition-structure-injectives-noetherian}
Soit $R$ un anneau noethérien.
Tout module injectif est somme directe de modules injectifs indécomposables.
Tout module injectif indécomposable est l'enveloppe injective
du corps résiduel en un idéal premier.
\end{proposition}

\begin{proof}
La seconde assertion est le lemme \ref{lemma-indecomposable-injective-noetherian}.
Pour la première, soit $I$ un $R$-module injectif.
Construisons par récurrence transfinie des $I_\alpha \subset I$
pour des ordinaux $\alpha$, sommes directes de $R$-modules injectifs
indécomposables $E_{\beta + 1}$ pour $\beta < \alpha$.
Pour $\alpha = 0$, posons $I_0 = 0$. Supposons donné un ordinal $\alpha$
pour lequel $I_\alpha$ a été construit. Alors $I_\alpha$ est un
$R$-module injectif d'après le lemme \ref{lemma-sum-injective-modules}.
Ainsi $I \cong I_\alpha \oplus I'$. Si $I' = 0$, c'est terminé.
Sinon, $I'$ possède un idéal premier associé d'après
Algèbre, lemme \ref{algebra-lemma-ass-zero}.
Le module $I'$ contient donc une copie de $R/\mathfrak p$ pour un idéal premier $\mathfrak p$.
Ainsi $I'$ contient un sous-module indécomposable $E$ d'après
les lemmes \ref{lemma-injective-hull-unique} et \ref{lemma-injective-hull-indecomposable}. Posons $I_{\alpha + 1} = I_\alpha \oplus E_\alpha$.
Si $\alpha$ est un ordinal limite et si $I_\beta$ a été construit
pour $\beta < \alpha$, posons $I_\alpha = \bigcup_{\beta < \alpha} I_\beta$.
Observons que $I_\alpha = \bigoplus_{\beta < \alpha} E_{\beta + 1}$.
Cela achève la démonstration.
\end{proof}



\section{Dualité sur les anneaux locaux artiniens}
\label{section-artinian}

\noindent
Soit $(R, \mathfrak m, \kappa)$ un anneau local artinien.
Rappelons qu'alors $R$ est noethérien et que $R$ est de longueur finie
en tant que $R$-module. De plus, un $R$-module est de type fini
si et seulement s'il est de longueur finie. Nous utiliserons ces faits
sans autre mention dans cette section. Pour plus de détails, voir
Algèbre, sections \ref{algebra-section-length} et \ref{algebra-section-artinian},
ainsi que Algèbre, proposition \ref{algebra-proposition-dimension-zero-ring}.

\begin{lemma}
\label{lemma-finite}
Soit $(R, \mathfrak m, \kappa)$ un anneau local artinien.
Soit $E$ une enveloppe injective de $\kappa$. Pour tout
$R$-module de type fini $M$, on a
$$
\text{longueur}_R(M) = \text{longueur}_R(\Hom_R(M, E))
$$
En particulier, l'enveloppe injective $E$ de $\kappa$ est un $R$-module de type fini.
\end{lemma}

\begin{proof}
Puisque $E$ est une extension essentielle de $\kappa$, on a
$\kappa = E[\mathfrak m]$, où $E[\mathfrak m]$ désigne la $\mathfrak m$-torsion de $E$
(avec les notations de Compléments d'algèbre, section \ref{more-algebra-section-torsion}).
Ainsi $\Hom_R(\kappa, E) \cong \kappa$ et l'égalité des longueurs est vraie pour $M = \kappa$.
Démontrons l'égalité affichée dans le lemme par récurrence sur
la longueur de $M$. Si $M$ est non nul, il existe une surjection
$M \to \kappa$ de noyau $M'$. Puisque le foncteur $M \mapsto \Hom_R(M, E)$
est exact, on obtient une suite exacte courte
$$
0 \to \Hom_R(\kappa, E) \to \Hom_R(M, E) \to \Hom_R(M', E) \to 0.
$$
L'additivité de la longueur pour cette suite et pour la suite $0 \to M' \to M \to \kappa \to 0$,
ainsi que l'égalité pour $M'$ (hypothèse de récurrence) et pour $\kappa$,
donnent l'égalité pour $M$.
La dernière assertion du lemme résulte de l'égalité $E = \Hom_R(R, E)$.
\end{proof}

\begin{lemma}
\label{lemma-evaluate}
Soit $(R, \mathfrak m, \kappa)$ un anneau local artinien.
Soit $E$ une enveloppe injective de $\kappa$.
Pour tout $R$-module de type fini $M$, le morphisme d'évaluation
$$
M \longrightarrow \Hom_R(\Hom_R(M, E), E)
$$
est un isomorphisme. En particulier, $R = \Hom_R(E, E)$.
\end{lemma}

\begin{proof}
Observons que la flèche affichée est injective. En effet, si $x \in M$
est un élément non nul, il existe un morphisme non nul $Rx \to \kappa$ que l'on peut
prolonger en un morphisme $\varphi : M \to E$ ne s'annulant pas sur $x$.
Puisque la source et le but de la flèche ont même longueur d'après
le lemme \ref{lemma-finite}, c'est un isomorphisme.
La dernière assertion s'obtient en prenant $M = R$.
\end{proof}

\noindent
Pour énoncer le lemme suivant, notons $\text{Mod}^{fg}_R$ la catégorie des
$R$-modules de type fini sur un anneau $R$.

\begin{lemma}
\label{lemma-duality}
Soit $(R, \mathfrak m, \kappa)$ un anneau local artinien.
Soit $E$ une enveloppe injective de $\kappa$.
Le foncteur $D(-) = \Hom_R(-, E)$ induit une anti-équivalence exacte
$\text{Mod}^{fg}_R \to \text{Mod}^{fg}_R$, et $D \circ D \cong \text{id}$.
\end{lemma}

\begin{proof}
Nous avons vu que $D \circ D = \text{id}$ sur $\text{Mod}^{fg}_R$ dans le lemme \ref{lemma-evaluate}.
Il en résulte immédiatement que $D$ est une anti-équivalence.
\end{proof}

\begin{lemma}
\label{lemma-duality-torsion-cotorsion}
Gardons les hypothèses et les notations du lemme \ref{lemma-duality}.
Soient $I \subset R$ un idéal et $M$ un $R$-module de type fini.
Alors
$$
D(M[I]) = D(M)/ID(M) \quad\text{et}\quad D(M/IM) = D(M)[I]
$$
\end{lemma}

\begin{proof}
Écrivons $I = (f_1, \ldots, f_t)$. Considérons le morphisme
$$
M^{\oplus t} \xrightarrow{f_1, \ldots, f_t} M
$$
de conoyau $M/IM$. En appliquant le foncteur exact $D$,
on en déduit que $D(M/IM)$ est $D(M)[I]$. L'autre cas se démontre de la même manière.
\end{proof}



\section{Enveloppe injective du corps résiduel}
\label{section-hull-residue-field}

\noindent
Dans cette section, la plupart des résultats concernent les anneaux locaux noethériens.

\begin{lemma}
\label{lemma-quotient}
Soit $R \to S$ un homomorphisme surjectif d'anneaux locaux, de noyau $I$.
Soit $E$ l'enveloppe injective du corps résiduel de $R$ sur $R$.
Alors $E[I]$ est l'enveloppe injective du corps résiduel de $S$ sur $S$.
\end{lemma}

\begin{proof}
Observons que $E[I] = \Hom_R(S, E)$ puisque $S = R/I$. Ainsi $E[I]$ est un $S$-module injectif
d'après le lemme \ref{lemma-hom-injective}. Puisque $E$ est une extension essentielle
de $\kappa = R/\mathfrak m_R$, il en résulte que $E[I]$ est aussi une extension essentielle
de $\kappa$. D'où le résultat.
\end{proof}

\begin{lemma}
\label{lemma-torsion-submodule-sum-injective-hulls}
Soit $(R, \mathfrak m, \kappa)$ un anneau local.
Soit $E$ l'enveloppe injective de $\kappa$.
Soit $M$ un module de torsion par les puissances de $\mathfrak m$ sur $R$,
tel que $n = \dim_\kappa(M[\mathfrak m]) < \infty$.
Alors $M$ est isomorphe à un sous-module de $E^{\oplus n}$.
\end{lemma}

\begin{proof}
Observons que $E^{\oplus n}$ est l'enveloppe injective de $\kappa^{\oplus n} = M[\mathfrak m]$.
Il existe donc un homomorphisme de $R$-modules
$M \to E^{\oplus n}$ qui est injectif sur $M[\mathfrak m]$.
Puisque $M$ est de torsion par les puissances de $\mathfrak m$,
l'inclusion $M[\mathfrak m] \subset M$ est une extension essentielle
(par exemple d'après le lemme \ref{lemma-essential-extension}).
On en déduit que le noyau de $M \to E^{\oplus n}$ est nul.
\end{proof}

\begin{lemma}
\label{lemma-union-artinian}
Soit $(R, \mathfrak m, \kappa)$ un anneau local noethérien.
Soit $E$ une enveloppe injective de $\kappa$ sur $R$.
Soit $E_n$ une enveloppe injective de $\kappa$ sur $R/\mathfrak m^n$.
Alors $E = \bigcup E_n$ et $E_n = E[\mathfrak m^n]$.
\end{lemma}

\begin{proof}
On a $E_n = E[\mathfrak m^n]$ d'après le lemme \ref{lemma-quotient}.
On a $E = \bigcup E_n$ parce que $\bigcup E_n = E[\mathfrak m^\infty]$ est un sous-$R$-module injectif
qui contient $\kappa$ ; voir le lemme \ref{lemma-injective-module-divide}.
\end{proof}

\noindent
Le lemme suivant montre que l'enveloppe injective du corps résiduel
d'un anneau local noethérien ne dépend que du complété.

\begin{lemma}
\label{lemma-compare}
Soit $R \to S$ un homomorphisme local plat d'anneaux locaux noethériens
tel que $R/\mathfrak m_R \cong S/\mathfrak m_R S$.
Alors l'enveloppe injective du corps résiduel de $R$
est l'enveloppe injective du corps résiduel de $S$.
\end{lemma}

\begin{proof}
Notons que $\mathfrak m_RS = \mathfrak m_S$, puisque le quotient par le premier idéal est un corps.
Posons $\kappa = R/\mathfrak m_R = S/\mathfrak m_S$.
Soit $E_R$ l'enveloppe injective de $\kappa$ sur $R$.
Soit $E_S$ l'enveloppe injective de $\kappa$ sur $S$.
Observons que $E_S$ est un $R$-module injectif d'après le lemme \ref{lemma-injective-flat}.
Choisissons un prolongement $E_R \to E_S$ de l'identification des corps résiduels.
Ce morphisme est un isomorphisme d'après le lemme \ref{lemma-union-artinian},
car $R \to S$ induit un isomorphisme $R/\mathfrak m_R^n \to S/\mathfrak m_S^n$ pour tout $n$.
\end{proof}

\begin{lemma}
\label{lemma-endos}
Soit $(R, \mathfrak m, \kappa)$ un anneau local noethérien.
Soit $E$ une enveloppe injective de $\kappa$ sur $R$.
Alors $\Hom_R(E, E)$ est canoniquement isomorphe au complété de $R$.
\end{lemma}

\begin{proof}
Écrivons $E = \bigcup E_n$ avec $E_n = E[\mathfrak m^n]$ comme dans le lemme \ref{lemma-union-artinian}.
Tout endomorphisme de $E$ préserve cette filtration. Ainsi
$$
\Hom_R(E, E) = \lim \Hom_R(E_n, E_n)
$$
Le lemme en résulte, puisque
$\Hom_R(E_n, E_n) = \Hom_{R/\mathfrak m^n}(E_n, E_n) = R/\mathfrak m^n$
d'après le lemme \ref{lemma-evaluate}.
\end{proof}

\begin{lemma}
\label{lemma-injective-hull-has-dcc}
Soit $(R, \mathfrak m, \kappa)$ un anneau local noethérien.
Soit $E$ une enveloppe injective de $\kappa$ sur $R$.
Alors $E$ satisfait à la condition de chaîne descendante.
\end{lemma}

\begin{proof}
Si $E \supset M_1 \supset M_2 \supset \ldots$ est une suite de sous-modules, alors
$$
\Hom_R(E, E) \to \Hom_R(M_1, E) \to \Hom_R(M_2, E) \to \ldots
$$
est une suite de surjections. D'après le lemme \ref{lemma-endos},
chacun de ces modules est un module sur le complété $R^\wedge = \Hom_R(E, E)$.
Puisque $R^\wedge$ est noethérien (Algèbre, lemme \ref{algebra-lemma-completion-Noetherian-Noetherian}),
la suite stationne : $\Hom_R(M_n, E) = \Hom_R(M_{n + 1}, E) = \ldots$.
Puisque $E$ est injectif, cela ne peut se produire que si $\Hom_R(M_n/M_{n + 1}, E)$ est nul.
Or, si $M_n/M_{n + 1}$ est non nul, il contient un élément non nul annulé par $\mathfrak m$,
car $E$ est de torsion par les puissances de $\mathfrak m$ d'après le lemme \ref{lemma-union-artinian}.
Dans ce cas, $M_n/M_{n + 1}$ admet un morphisme non nul vers $E$,
ce qui contredit l'annulation supposée. Cela achève la démonstration.
\end{proof}

\begin{lemma}
\label{lemma-describe-categories}
Soit $(R, \mathfrak m, \kappa)$ un anneau local noethérien.
Soit $E$ une enveloppe injective de $\kappa$.
\begin{enumerate}
\item Pour un $R$-module $M$, les conditions suivantes sont équivalentes :
\begin{enumerate}
\item $M$ satisfait à la condition de chaîne ascendante ;
\item $M$ est un $R$-module de type fini ;
\item il existe $n, m$ et une suite exacte
$R^{\oplus m} \to R^{\oplus n} \to M \to 0$.
\end{enumerate}
\item Pour un $R$-module $M$, les conditions suivantes sont équivalentes :
\begin{enumerate}
\item $M$ satisfait à la condition de chaîne descendante ;
\item $M$ est de torsion par les puissances de $\mathfrak m$ et
$\dim_\kappa(M[\mathfrak m]) < \infty$ ;
\item il existe $n, m$ et une suite exacte
$0 \to M \to E^{\oplus n} \to E^{\oplus m}$.
\end{enumerate}
\end{enumerate}
\end{lemma}

\begin{proof}
Nous omettons la démonstration de (1).

\medskip\noindent
Soit $M$ un $R$-module satisfaisant à la condition de chaîne descendante.
Soit $x \in M$. Alors $\mathfrak m^n x$ est une chaîne descendante de sous-modules,
donc elle stationne. Ainsi $\mathfrak m^nx = \mathfrak m^{n + 1}x$ pour un certain $n$.
Le lemme de Nakayama (Algèbre, lemme \ref{algebra-lemma-NAK}) entraîne alors $\mathfrak m^n x = 0$,
c'est-à-dire que $x$ est de torsion par les puissances de $\mathfrak m$.
Puisque $M[\mathfrak m]$ est un espace vectoriel sur $\kappa$, il doit être de dimension
finie pour satisfaire à la condition de chaîne descendante.

\medskip\noindent
Supposons que $M$ soit de torsion par les puissances de $\mathfrak m$
et que son sous-module de $\mathfrak m$-torsion $M[\mathfrak m]$ soit de dimension finie.
D'après le lemme \ref{lemma-torsion-submodule-sum-injective-hulls}, on voit que $M$ est un sous-module
de $E^{\oplus n}$ pour un certain $n$.
Considérons le quotient $N = E^{\oplus n}/M$. D'après le lemme \ref{lemma-injective-hull-has-dcc},
le module $E$ satisfait à la condition de chaîne descendante,
donc il en est de même de $E^{\oplus n}$ et de $N$.
Ainsi $N$ satisfait à (2)(a), ce qui entraîne que $N$ satisfait
à (2)(b) d'après le deuxième paragraphe de la démonstration.
En appliquant à nouveau le lemme \ref{lemma-torsion-submodule-sum-injective-hulls}, on voit donc que $N$
est un sous-module de $E^{\oplus m}$ pour un certain $m$.
On a ainsi une suite exacte courte
$0 \to M \to E^{\oplus n} \to E^{\oplus m}$.

\medskip\noindent
Supposons donnée une suite exacte courte
$0 \to M \to E^{\oplus n} \to E^{\oplus m}$.
Puisque $E$ satisfait à la condition de chaîne descendante
d'après le lemme \ref{lemma-injective-hull-has-dcc}, il en est de même de $M$.
\end{proof}

\begin{proposition}[Dualité de Matlis]
\label{proposition-matlis}
Soit $(R, \mathfrak m, \kappa)$ un anneau local noethérien complet.
Soit $E$ l'enveloppe injective de $\kappa$ sur $R$.
Le foncteur $D(-) = \Hom_R(-, E)$ est une anti-équivalence
$$
\left\{
\begin{matrix}
R\text{-modules satisfaisant la} \\
\text{condition de chaîne descendante}
\end{matrix}
\right\}
\longleftrightarrow
\left\{
\begin{matrix}
R\text{-modules satisfaisant la} \\
\text{condition de chaîne ascendante}
\end{matrix}
\right\}
$$
et l'on a $D \circ D = \text{id}$ de chaque côté de l'équivalence.
\end{proposition}

\begin{proof}
D'après le lemme \ref{lemma-endos}, on a $R = \Hom_R(E, E) = D(E)$.
Bien entendu, on a $E = \Hom_R(R, E) = D(R)$. Puisque $E$ est injectif,
le foncteur $D$ est exact. Le résultat découle alors immédiatement
de la description des catégories donnée dans le lemme \ref{lemma-describe-categories}.
\end{proof}

\begin{remark}
\label{remark-matlis}
Soit $(R, \mathfrak m, \kappa)$ un anneau local noethérien.
Soit $E$ une enveloppe injective de $\kappa$ sur $R$.
Voici un complément à la dualité de Matlis : si $N$ est un module
de torsion par les puissances de $\mathfrak m$ et si $M = \Hom_R(N, E)$ est un module de type fini
sur le complété de $R$, alors $N$ satisfait à la condition de chaîne
descendante. En effet, pour tous sous-modules $N'' \subset N' \subset N$ tels que $N'' \not = N'$,
il existe un plongement $\kappa \subset N''/N'$, donc un morphisme non nul $N' \to E$
annulant $N''$, que l'on peut prolonger en un morphisme $N \to E$ annulant $N''$.
Ainsi $N \supset N' \mapsto M' = \Hom_R(N/N', E) \subset M$ est une application préservant les inclusions des sous-modules
de $N$ vers les sous-modules de $M$, d'où la conclusion.
\end{remark}




















\section{Dérivation du foncteur de torsion}
\label{section-bad-local-cohomology}

\noindent
Soient $A$ un anneau et $I \subset A$ un idéal de type fini
(si $I$ n'est pas de type fini, il faudrait peut-être utiliser
une autre définition). Posons $Z = V(I) \subset \Spec(A)$. Rappelons que la catégorie $I^\infty\text{-torsion}$
des modules de torsion par les puissances de $I$ ne dépend que
du fermé $Z$ et non du choix de l'idéal de type fini $I$ tel que $Z = V(I)$ ;
voir Compléments d'algèbre, lemme \ref{more-algebra-lemma-local-cohomology-closed}.
Dans cette section, nous considérerons le foncteur
$$
H^0_{I} : \text{Mod}_A \longrightarrow I^\infty\text{-torsion},\quad
M \longmapsto M[I^\infty] = \bigcup M[I^n]
$$
qui associe à $M$ son sous-module de torsion par les puissances de $I$.

\medskip\noindent
Soient $A$ un anneau et $I$ un idéal de type fini.
Notons que $I^\infty\text{-torsion}$ est une catégorie abélienne de Grothendieck
(les sommes directes existent, les limites inductives filtrantes sont
exactes et $\bigoplus A/I^n$ est un générateur d'après
Compléments d'algèbre, lemme \ref{more-algebra-lemma-I-power-torsion-presentation}).
La catégorie dérivée $D(I^\infty\text{-torsion})$ existe donc ; voir Injectifs, remarque \ref{injectives-remark-existence-D}.
Notre foncteur $H^0_I$ est exact à gauche et possède un prolongement dérivé
que nous noterons
$$
R\Gamma_I : D(A) \longrightarrow D(I^\infty\text{-torsion}).
$$
{\bf Attention :} ce foncteur ne mérite le nom de cohomologie
locale que si l'anneau $A$ est noethérien.
Les foncteurs $H^0_I$, $R\Gamma_I$ et les satellites $H^p_I$
ne dépendent que du fermé $Z \subset \Spec(A)$ et non du choix de l'idéal
de type fini $I$ tel que $V(I) = Z$.
Nous tenons cependant à utiliser l'indice $I$ pour les foncteurs
ci-dessus, car la notation $R\Gamma_Z$ sera réservée à un autre foncteur,
voir (\ref{equation-local-cohomology}), qui ne coïncide avec le foncteur $R\Gamma_I$
(à notre connaissance) que si $A$ est noethérien
(voir le lemme \ref{lemma-local-cohomology-noetherian}).

\begin{lemma}
\label{lemma-adjoint}
Soient $A$ un anneau et $I \subset A$ un idéal de type fini.
Le foncteur $R\Gamma_I$ est adjoint à droite du foncteur
$D(I^\infty\text{-torsion}) \to D(A)$.
\end{lemma}

\begin{proof}
Cela résulte du fait que le passage au sous-module de torsion
par les puissances de $I$ est adjoint à droite du foncteur d'inclusion
$I^\infty\text{-torsion} \to \text{Mod}_A$. Voir Catégories dérivées, lemme \ref{derived-lemma-derived-adjoint-functors}.
\end{proof}

\begin{lemma}
\label{lemma-local-cohomology-ext}
Soient $A$ un anneau et $I \subset A$ un idéal de type fini.
Pour tout objet $K$ de $D(A)$, on a
$$
R\Gamma_I(K) = \text{hocolim}\ R\Hom_A(A/I^n, K)
$$
dans $D(A)$ et
$$
R^q\Gamma_I(K) = \colim_n \Ext_A^q(A/I^n, K)
$$
en tant que modules, pour tout $q \in \mathbf{Z}$.
\end{lemma}

\begin{proof}
Soit $J^\bullet$ un complexe K-injectif représentant $K$. Alors
$$
R\Gamma_I(K) = J^\bullet[I^\infty] = \colim J^\bullet[I^n] =
\colim \Hom_A(A/I^n, J^\bullet)
$$
où la première égalité est la définition de $R\Gamma_I(K)$.
Catégories dérivées, lemme \ref{derived-lemma-colim-hocolim}, donne la première égalité
affichée dans l'énoncé. La seconde en résulte, puisque $H^q(\Hom_A(A/I^n, J^\bullet)) = \Ext^q_A(A/I^n, K)$
et puisque les limites inductives filtrantes sont exactes
dans la catégorie des groupes abéliens.
\end{proof}

\begin{lemma}
\label{lemma-bad-local-cohomology-vanishes}
Soient $A$ un anneau et $I \subset A$ un idéal de type fini.
Soit $K^\bullet$ un complexe de $A$-modules tel que
$f : K^\bullet \to K^\bullet$ soit un isomorphisme pour un certain $f \in I$,
c'est-à-dire que $K^\bullet$ soit un complexe de $A_f$-modules. Alors
$R\Gamma_I(K^\bullet) = 0$.
\end{lemma}

\begin{proof}
En effet, les modules de cohomologie de $R\Gamma_I(K^\bullet)$ sont alors de torsion
par les puissances de $f$ et $f$ y agit par automorphismes.
Ces modules de cohomologie sont donc nuls, et l'objet est donc nul.
\end{proof}

\noindent
Soient $A$ un anneau et $I \subset A$ un idéal de type fini.
D'après Compléments d'algèbre, lemme \ref{more-algebra-lemma-I-power-torsion}, la catégorie des modules
de torsion par les puissances de $I$ est une sous-catégorie de Serre
de la catégorie de tous les $A$-modules. On a donc un foncteur
\begin{equation}
\label{equation-compare-torsion}
D(I^\infty\text{-torsion}) \longrightarrow D_{I^\infty\text{-torsion}}(A)
\end{equation}
où le membre de droite est la sous-catégorie pleine de $D(A)$
formée des objets dont les cohomologies sont des modules de torsion
par les puissances de $I$. Le foncteur et la catégorie sont tous deux
étudiés dans Catégories dérivées, section \ref{derived-section-triangulated-sub}.

\begin{lemma}
\label{lemma-not-equal}
Soient $A$ un anneau et $I$ un idéal de type fini.
Soient $M$ et $N$ des modules de torsion par les puissances de $I$.
\begin{enumerate}
\item $\Hom_{D(A)}(M, N) = \Hom_{D({I^\infty\text{-torsion}})}(M, N)$ ;
\item $\Ext^1_{D(A)}(M, N) =
\Ext^1_{D({I^\infty\text{-torsion}})}(M, N)$ ;
\item $\Ext^2_{D({I^\infty\text{-torsion}})}(M, N) \to
\Ext^2_{D(A)}(M, N)$ n'est pas surjectif en général ;
\item (\ref{equation-compare-torsion}) n'est pas une équivalence en général.
\end{enumerate}
\end{lemma}

\begin{proof}
Les assertions (1) et (2) résultent immédiatement de ce que les modules
de torsion par les puissances de $I$ forment une sous-catégorie de Serre
de $\text{Mod}_A$. L'assertion (4) découle de (3).

\medskip\noindent
Pour (3), soit $A$ un anneau possédant un élément $f \in A$ tel que
$A[f]$ contienne un élément non nul $x$ annulé par $f$ et que
$A$ contienne des éléments $x_n$ tels que $f^nx_n = x$.
Un tel anneau $A$ existe, car on peut prendre
$$
A = \mathbf{Z}[f, x, x_n]/(fx, f^nx_n - x)
$$
Étant donné $A$, posons $I = (f)$. Alors la suite exacte
$$
0 \to A[f] \to A \xrightarrow{f} A \to A/fA \to 0
$$
définit un élément de $\Ext^2_A(A/fA, A[f])$. Affirmons que cet élément
ne provient pas d'un élément de
$\Ext^2_{D(f^\infty\text{-torsion})}(A/fA, A[f])$.
En effet, sinon, il existerait une suite exacte
$$
0 \to A[f] \to M \to N \to A/fA \to 0
$$
où $M$ et $N$ sont des modules de torsion par les puissances de $f$,
définissant la même classe de $2$-extension.
Puisque $A \to A$ est un complexe de modules libres et que les classes
de $2$-extensions sont les mêmes, on pourrait trouver un morphisme
$$
\xymatrix{
0 \ar[r] &
A[f] \ar[r] \ar[d] &
A \ar[r] \ar[d]_\varphi &
A \ar[r] \ar[d]_\psi &
A/fA \ar[r] \ar[d] & 0 \\
0 \ar[r] &
A[f] \ar[r] &
M \ar[r] &
N \ar[r] &
A/fA \ar[r] & 0
}
$$
(nous omettons certains détails). On pourrait alors remplacer $M$
par l'image de $\varphi$ et $N$ par l'image de $\psi$.
Alors $M$ serait un module cyclique, donc $f^n M = 0$ pour un certain $n$.
En considérant $\varphi(x_{n + 1})$, on obtient une contradiction avec le fait que
$f^{n + 1}x_n = x$ est non nul dans $A[f]$.
\end{proof}









\section{Cohomologie locale}
\label{section-local-cohomology}

\noindent
Soient $A$ un anneau et $I \subset A$ un idéal de type fini.
Posons $Z = V(I) \subset \Spec(A)$. Nous construirons un foncteur
\begin{equation}
\label{equation-local-cohomology}
R\Gamma_Z : D(A) \longrightarrow D_{I^\infty\text{-torsion}}(A).
\end{equation}
adjoint à droite du foncteur d'inclusion. Pour les notations,
voir la section \ref{section-bad-local-cohomology}. Les modules de cohomologie de $R\Gamma_Z(K)$
sont les {\it groupes de cohomologie locale de $K$ relativement à $Z$}.
D'après le lemme \ref{lemma-not-equal}, ce foncteur ne sera en général {\bf pas}
égal à $R\Gamma_I( - )$, même en les considérant comme foncteurs à valeurs dans $D(A)$.
Dans la section \ref{section-local-cohomology-noetherian}, nous montrerons que, si $A$ est noethérien,
les deux coïncident.

\medskip\noindent
Nous poursuivrons l'étude de la cohomologie locale dans le chapitre
qui lui est consacré ; voir Cohomologie locale, section \ref{local-cohomology-section-introduction}.
Nous y montrerons notamment que $R\Gamma_Z$ calcule la cohomologie
à support dans $Z$ du complexe associé de faisceaux quasi-cohérents
sur $\Spec(A)$. Voir Cohomologie locale, lemme \ref{local-cohomology-lemma-local-cohomology-is-local-cohomology}.


\begin{lemma}
\label{lemma-local-cohomology-adjoint}
Soient $A$ un anneau et $I \subset A$ un idéal de type fini.
Il existe un adjoint à droite $R\Gamma_Z$ (\ref{equation-local-cohomology})
du foncteur d'inclusion $D_{I^\infty\text{-torsion}}(A) \to D(A)$.
En fait, si $I$ est engendré par $f_1, \ldots, f_r \in A$, on a
$$
R\Gamma_Z(K) =
(A \to \prod\nolimits_{i_0} A_{f_{i_0}} \to
\prod\nolimits_{i_0 < i_1} A_{f_{i_0}f_{i_1}}
\to \ldots \to A_{f_1\ldots f_r}) \otimes_A^\mathbf{L} K
$$
fonctoriellement en $K \in D(A)$.
\end{lemma}

\begin{proof}
Soit $I = (f_1, \ldots, f_r)$ un idéal.
Soit $K^\bullet$ un complexe de $A$-modules.
Il existe un morphisme canonique de complexes
$$
(A \to \prod\nolimits_{i_0} A_{f_{i_0}} \to
\prod\nolimits_{i_0 < i_1} A_{f_{i_0}f_{i_1}} \to
\ldots \to A_{f_1\ldots f_r}) \longrightarrow A.
$$
du complexe de {\v C}ech augmenté vers $A$.
En tensorisant par $K^\bullet$ et en prenant le complexe total associé,
on obtient un morphisme
$$
\text{Tot}\left(
K^\bullet \otimes_A
(A \to \prod\nolimits_{i_0} A_{f_{i_0}} \to
\prod\nolimits_{i_0 < i_1} A_{f_{i_0}f_{i_1}} \to
\ldots \to A_{f_1\ldots f_r})\right)
\longrightarrow
K^\bullet
$$
dans $D(A)$. Affirmons que les modules de cohomologie du complexe
de gauche sont de torsion par les puissances de $I$,
c'est-à-dire que le membre de gauche est un objet de $D_{I^\infty\text{-torsion}}(A)$.
En effet, on a
$$
(A \to \prod\nolimits_{i_0} A_{f_{i_0}} \to
\prod\nolimits_{i_0 < i_1} A_{f_{i_0}f_{i_1}} \to
\ldots \to A_{f_1\ldots f_r}) = \colim K(A, f_1^n, \ldots, f_r^n)
$$
d'après Compléments d'algèbre, lemme \ref{more-algebra-lemma-extended-alternating-Cech-is-colimit-koszul}.
De plus, la multiplication par $f_i^n$ sur le complexe
$K(A, f_1^n, \ldots, f_r^n)$ est homotope à zéro d'après Compléments d'algèbre, lemme \ref{more-algebra-lemma-homotopy-koszul}.
Puisque
$$
H^q\left( LHS \right) =
\colim H^q(\text{Tot}(K^\bullet \otimes_A K(A, f_1^n, \ldots, f_r^n)))
$$
on obtient l'assertion. D'autre part, si $K^\bullet$ est un objet de $D_{I^\infty\text{-torsion}}(A)$,
les complexes $K^\bullet \otimes_A A_{f_{i_0} \ldots f_{i_p}}$ ont une cohomologie nulle.
Dans ce cas, le morphisme $LHS \to K^\bullet$ est donc un isomorphisme dans $D(A)$.
La construction
$$
R\Gamma_Z(K^\bullet) =
\text{Tot}\left(
K^\bullet \otimes_A
(A \to \prod\nolimits_{i_0} A_{f_{i_0}} \to
\prod\nolimits_{i_0 < i_1} A_{f_{i_0}f_{i_1}} \to
\ldots \to A_{f_1\ldots f_r})\right)
$$
est fonctorielle en $K^\bullet$ et définit un foncteur exact
$D(A) \to D_{I^\infty\text{-torsion}}(A)$ entre catégories triangulées.
L'existence de la transformation naturelle $R\Gamma_Z \to \text{id}$ donnée ci-dessus,
et le fait qu'elle soit un isomorphisme sur les $K^\bullet$ appartenant
à la sous-catégorie, entraînent formellement que $R\Gamma_Z$
est l'adjoint à droite cherché.
\end{proof}

\begin{lemma}
\label{lemma-local-cohomology-and-restriction}
Soient $A \to B$ un homomorphisme d'anneaux et $I \subset A$ un idéal de type fini.
Posons $J = IB$, puis $Z = V(I)$ et $Y = V(J)$. Alors
$$
R\Gamma_Z(M_A) = R\Gamma_Y(M)_A
$$
fonctoriellement en $M \in D(B)$. Ici $(-)_A$ désigne les restrictions
$D(B) \to D(A)$ et $D_{J^\infty\text{-torsion}}(B) \to D_{I^\infty\text{-torsion}}(A)$.
\end{lemma}

\begin{proof}
Cela résulte de l'unicité des adjoints, puisque $R\Gamma_Z((-)_A)$ et $R\Gamma_Y(-)_A$
sont tous deux adjoints à droite du foncteur $D_{I^\infty\text{-torsion}}(A) \to D(B)$,
$K \mapsto K \otimes_A^\mathbf{L} B$.
On peut aussi utiliser la description de $R\Gamma_Z$ et de $R\Gamma_Y$
par les complexes alternés de {\v C}ech
(lemme \ref{lemma-local-cohomology-adjoint}).
En effet, si $I = (f_1, \ldots, f_r)$, alors $J$ est engendré par les images
$g_1, \ldots, g_r \in B$ de $f_1, \ldots, f_r$.
L'énoncé résulte alors de l'existence d'un isomorphisme canonique
\begin{align*}
& M_A \otimes_A (A \to \prod\nolimits_{i_0} A_{f_{i_0}} \to
\prod\nolimits_{i_0 < i_1} A_{f_{i_0}f_{i_1}}
\to \ldots \to A_{f_1\ldots f_r}) \\
& = 
M \otimes_B (B \to \prod\nolimits_{i_0} B_{g_{i_0}} \to
\prod\nolimits_{i_0 < i_1} B_{g_{i_0}g_{i_1}}
\to \ldots \to B_{g_1\ldots g_r})
\end{align*}
pour tout $B$-module $M$.
\end{proof}

\begin{lemma}
\label{lemma-torsion-change-rings}
Soient $A \to B$ un homomorphisme d'anneaux et $I \subset A$ un idéal de type fini.
Posons $J = IB$. Soient $Z = V(I)$ et $Y = V(J)$. Alors
$$
R\Gamma_Z(K) \otimes_A^\mathbf{L} B = R\Gamma_Y(K \otimes_A^\mathbf{L} B)
$$
fonctoriellement en $K \in D(A)$.
\end{lemma}

\begin{proof}
Écrivons $I = (f_1, \ldots, f_r)$. Alors $J$ est engendré par les images
$g_1, \ldots, g_r \in B$ de $f_1, \ldots, f_r$. On a donc
$$
(A \to \prod A_{f_{i_0}} \to \ldots \to A_{f_1\ldots f_r}) \otimes_A B =
(B \to \prod B_{g_{i_0}} \to \ldots \to B_{g_1\ldots g_r})
$$
en tant que complexes de $B$-modules. Représentons $K$
par un complexe K-plat $K^\bullet$ de $A$-modules.
Puisque les complexes totaux associés à
$$
K^\bullet \otimes_A
(A \to \prod A_{f_{i_0}} \to \ldots \to A_{f_1\ldots f_r}) \otimes_A B
$$
et à
$$
K^\bullet \otimes_A B \otimes_B
(B \to \prod B_{g_{i_0}} \to \ldots \to B_{g_1\ldots g_r})
$$
représentent les membres de gauche et de droite de la formule
affichée dans le lemme (voir le lemme \ref{lemma-local-cohomology-adjoint}), on conclut.
\end{proof}

\begin{lemma}
\label{lemma-local-cohomology-vanishes}
Soient $A$ un anneau et $I \subset A$ un idéal de type fini.
Soit $K^\bullet$ un complexe de $A$-modules tel que
$f : K^\bullet \to K^\bullet$ soit un isomorphisme pour un certain $f \in I$,
c'est-à-dire que $K^\bullet$ soit un complexe de $A_f$-modules. Alors
$R\Gamma_Z(K^\bullet) = 0$.
\end{lemma}

\begin{proof}
En effet, les modules de cohomologie de $R\Gamma_Z(K^\bullet)$ sont alors de torsion
par les puissances de $f$ et $f$ y agit par automorphismes.
Ces modules de cohomologie sont donc nuls, et l'objet est donc nul.
\end{proof}

\begin{lemma}
\label{lemma-torsion-tensor-product}
Soient $A$ un anneau et $I \subset A$ un idéal de type fini.
Pour $K, L \in D(A)$, on a
$$
R\Gamma_Z(K \otimes_A^\mathbf{L} L) =
K \otimes_A^\mathbf{L} R\Gamma_Z(L) =
R\Gamma_Z(K) \otimes_A^\mathbf{L} L =
R\Gamma_Z(K) \otimes_A^\mathbf{L} R\Gamma_Z(L)
$$
Si $K$ ou $L$ appartient à $D_{I^\infty\text{-torsion}}(A)$,
il en est de même de $K \otimes_A^\mathbf{L} L$.
\end{lemma}

\begin{proof}
D'après le lemme \ref{lemma-local-cohomology-adjoint}, on sait que $R\Gamma_Z$ est donné par $C \otimes^\mathbf{L} -$
pour un certain $C \in D(A)$. Ainsi, pour $K, L \in D(A)$ quelconques, on a
$$
R\Gamma_Z(K \otimes_A^\mathbf{L} L) =
K \otimes^\mathbf{L} L \otimes_A^\mathbf{L} C =
K \otimes_A^\mathbf{L} R\Gamma_Z(L)
$$
Les autres égalités en découlent formellement.
Cela entraîne aussi la dernière assertion du lemme.
\end{proof}

\begin{lemma}
\label{lemma-local-cohomology-ss}
Soient $A$ un anneau et $I, J \subset A$ des idéaux de type fini.
Posons $Z = V(I)$ et $Y = V(J)$. Alors $Z \cap Y = V(I + J)$
et $R\Gamma_Y \circ R\Gamma_Z = R\Gamma_{Y \cap Z}$ en tant que foncteurs
$D(A) \to D_{(I + J)^\infty\text{-torsion}}(A)$. Pour $K \in D^+(A)$,
il existe une suite spectrale
$$
E_2^{p, q} = H^p_Y(H^q_Z(K)) \Rightarrow H^{p + q}_{Y \cap Z}(K)
$$
comme dans Catégories dérivées, lemme \ref{derived-lemma-grothendieck-spectral-sequence}.
\end{lemma}

\begin{proof}
Le lemme comporte un léger abus de notation, car, à proprement parler,
on ne peut pas composer $R\Gamma_Y$ et $R\Gamma_Z$.
L'énoncé signifie simplement que l'on compose $R\Gamma_Z$ avec l'inclusion $D_{I^\infty\text{-torsion}}(A) \to D(A)$,
puis avec $R\Gamma_Y$. L'égalité $R\Gamma_Y \circ R\Gamma_Z = R\Gamma_{Y \cap Z}$ résulte alors de ce que
$$
D_{I^\infty\text{-torsion}}(A) \to D(A) \xrightarrow{R\Gamma_Y}
D_{(I + J)^\infty\text{-torsion}}(A)
$$
est adjoint à droite de l'inclusion
$D_{(I + J)^\infty\text{-torsion}}(A) \to D_{I^\infty\text{-torsion}}(A)$.
On peut aussi démontrer la formule à l'aide du lemme \ref{lemma-local-cohomology-adjoint}
et du fait que le produit tensoriel des complexes de
{\v C}ech augmentés sur $f_1, \ldots, f_r$ et sur $g_1, \ldots, g_m$
est le complexe de {\v C}ech augmenté sur $f_1, \ldots, f_n. g_1, \ldots, g_m$.
La dernière assertion en résulte, avec le lemme cité.
\end{proof}

\noindent
Le lemme suivant est l'analogue de Compléments d'algèbre, lemme \ref{more-algebra-lemma-restriction-derived-complete-equivalence},
pour les complexes dont les cohomologies sont de torsion.

\begin{lemma}
\label{lemma-torsion-flat-change-rings}
Soient $A \to B$ un homomorphisme plat d'anneaux et $I \subset A$ un idéal de type fini
tel que $A/I = B/IB$. Le changement de base et la restriction induisent alors
des équivalences quasi-inverses
$D_{I^\infty\text{-torsion}}(A) = D_{(IB)^\infty\text{-torsion}}(B)$.
\end{lemma}

\begin{proof}
Plus précisément, les foncteurs sont $K \mapsto K \otimes_A^\mathbf{L} B$
pour $K$ dans $D_{I^\infty\text{-torsion}}(A)$ et $M \mapsto M_A$
pour $M$ dans $D_{(IB)^\infty\text{-torsion}}(B)$. Cela fonctionne parce que $H^i(K \otimes_A^\mathbf{L} B) = H^i(K) \otimes_A B = H^i(K)$.
La première égalité résulte de la platitude de $A \to B$ et la seconde
de Compléments d'algèbre, lemme \ref{more-algebra-lemma-neighbourhood-isomorphism}.
\end{proof}

\noindent
Le lemme suivant a été démontré pour les $\Hom$ et les $\Ext^1$ de modules
dans Compléments d'algèbre, lemmes \ref{more-algebra-lemma-neighbourhood-equivalence} et \ref{more-algebra-lemma-neighbourhood-extensions}.

\begin{lemma}
\label{lemma-neighbourhood-extensions}
Soient $A \to B$ un homomorphisme plat d'anneaux et $I \subset A$
un idéal de type fini tel que $A/I \to B/IB$ soit un isomorphisme.
Pour $K \in D_{I^\infty\text{-torsion}}(A)$ et $L \in D(A)$, le morphisme
$$
R\Hom_A(K, L) \longrightarrow R\Hom_B(K \otimes_A B, L \otimes_A B)
$$
est un quasi-isomorphisme. En particulier, si $M$, $N$
sont des $A$-modules et si $M$ est de torsion par les puissances de $I$,
le morphisme canonique
$$
\Ext^i_A(M, N)
\longrightarrow
\Ext^i_B(M \otimes_A B, N \otimes_A B)
$$
est un isomorphisme pour tout $i$.
\end{lemma}

\begin{proof}
Soient $Z = V(I) \subset \Spec(A)$ et $Y = V(IB) \subset \Spec(B)$.
Puisque les modules de cohomologie de $K$ sont de torsion par les puissances
de $I$, le morphisme canonique $R\Gamma_Z(L) \to L$ induit un isomorphisme
$$
R\Hom_A(K, R\Gamma_Z(L)) \to R\Hom_A(K, L)
$$
dans $D(A)$. De même, les modules de cohomologie de $K \otimes_A B$
sont de torsion par les puissances de $IB$ et l'on a un isomorphisme
$$
R\Hom_B(K \otimes_A B, R\Gamma_Y(L \otimes_A B)) \to 
R\Hom_B(K \otimes_A B, L \otimes_A B)
$$
dans $D(B)$.
D'après le lemme \ref{lemma-torsion-change-rings}, on a $R\Gamma_Z(L) \otimes_A B = R\Gamma_Y(L \otimes_A B)$.
Il suffit donc de montrer que le morphisme
$$
R\Hom_A(K, R\Gamma_Z(L)) \to R\Hom_B(K \otimes_A B, R\Gamma_Z(L) \otimes_A B)
$$
est un quasi-isomorphisme, ce qui résulte du lemme \ref{lemma-torsion-flat-change-rings}.
\end{proof}




\section{Cohomologie locale pour les anneaux noethériens}
\label{section-local-cohomology-noetherian}

\noindent
Soient $A$ un anneau et $I \subset A$ un idéal de type fini.
Posons $Z = V(I) \subset \Spec(A)$. Rappelons que (\ref{equation-compare-torsion}) est le foncteur
$$
D(I^\infty\text{-torsion}) \to D_{I^\infty\text{-torsion}}(A)
$$
En fait, il existe une transformation naturelle de foncteurs
\begin{equation}
\label{equation-compare-torsion-functors}
(\ref{equation-compare-torsion}) \circ R\Gamma_I(-)
\longrightarrow
R\Gamma_Z(-)
\end{equation}
En effet, étant donné un complexe de $A$-modules $K^\bullet$,
le morphisme canonique $R\Gamma_I(K^\bullet) \to K^\bullet$ dans $D(A)$ se factorise
(de manière unique) par $R\Gamma_Z(K^\bullet)$, puisque $R\Gamma_I(K^\bullet)$ a des modules
de cohomologie de torsion par les puissances de $I$
(voir le lemme \ref{lemma-adjoint}).
En général, ce morphisme n'est pas un isomorphisme
(nous l'avons vu dans le lemme \ref{lemma-not-equal}).

\begin{lemma}
\label{lemma-local-cohomology-noetherian}
Soient $A$ un anneau noethérien et $I \subset A$ un idéal.
\begin{enumerate}
\item Le morphisme d'adjonction $R\Gamma_I(K) \to K$ est un isomorphisme
pour $K \in D_{I^\infty\text{-torsion}}(A)$ ;
\item le foncteur
(\ref{equation-compare-torsion})
$D(I^\infty\text{-torsion}) \to D_{I^\infty\text{-torsion}}(A)$
est une équivalence ;
\item la transformation de foncteurs
(\ref{equation-compare-torsion-functors}) est un isomorphisme,
autrement dit $R\Gamma_I(K) = R\Gamma_Z(K)$ pour $K \in D(A)$.
\end{enumerate}
\end{lemma}

\begin{proof}
Un argument formel, que nous omettons, montre qu'il suffit de prouver (1).

\medskip\noindent
Soit $M$ un module de torsion par les puissances de $I$ sur $A$.
Choisissons un plongement $M \to J$ dans un $A$-module injectif.
Alors $J[I^\infty]$ est un $A$-module injectif, voir le lemme \ref{lemma-injective-module-divide},
et l'on obtient un plongement $M \to J[I^\infty]$.
Ainsi tout module de torsion par les puissances de $I$
possède une résolution injective $M \to J^\bullet$ dont les $J^n$
sont aussi de torsion par les puissances de $I$.
Il en résulte que $R\Gamma_I(M) = M$ (ce n'est pas une trivialité,
et ce n'est pas vrai en général si $A$ n'est pas noethérien).
Supposons ensuite que $K \in D_{I^\infty\text{-torsion}}^+(A)$. La suite spectrale
$$
R^q\Gamma_I(H^p(K)) \Rightarrow R^{p + q}\Gamma_I(K)
$$
(Catégories dérivées, lemme \ref{derived-lemma-two-ss-complex-functor}) converge,
et nous avons vu ci-dessus que seuls les termes avec $q = 0$
peuvent être non nuls. Ainsi $R\Gamma_I(K) \to K$ est un isomorphisme.

\medskip\noindent
Supposons que $K$ soit un objet quelconque de $D_{I^\infty\text{-torsion}}(A)$.
On a
$$
R^q\Gamma_I(K) = \colim \Ext^q_A(A/I^n, K)
$$
par le lemme \ref{lemma-local-cohomology-ext}. Prenons $f_1, \ldots, f_r \in A$ engendrant $I$.
Soit $K_n^\bullet = K(A, f_1^n, \ldots, f_r^n)$ le complexe de Koszul placé en degrés $-r, \ldots, 0$.
Puisque les pro-objets $\{A/I^n\}$ et $\{K_n^\bullet\}$ dans $D(A)$ sont les mêmes
d'après Compléments d'algèbre, lemme \ref{more-algebra-lemma-sequence-Koszul-complexes}, on voit que
$$
R^q\Gamma_I(K) = \colim \Ext^q_A(K_n^\bullet, K)
$$
Choisissons un complexe quelconque $K^\bullet$ de $A$-modules représentant $K$.
Puisque $K_n^\bullet$ est un complexe fini de modules libres de type fini,
on voit que
$$
\Ext^q_A(K_n, K) =
H^q(\text{Tot}((K_n^\bullet)^\vee \otimes_A K^\bullet))
$$
où $(K_n^\bullet)^\vee$ est le dual du complexe $K_n^\bullet$.
Voir Compléments d'algèbre, lemme \ref{more-algebra-lemma-RHom-out-of-projective}.
Puisque $(K_n^\bullet)^\vee$ est un complexe de $A$-modules libres de type fini
placés en degrés $0, \ldots, r$, les termes du complexe $\text{Tot}((K_n^\bullet)^\vee \otimes_A K^\bullet)$
sont les mêmes que ceux du complexe $\text{Tot}((K_n^\bullet)^\vee \otimes_A \tau_{\geq q - r - 2} K^\bullet)$
en degrés $q - 1$ et supérieurs. On voit donc que
$$
\Ext^q_A(K_n, K) = \text{Ext}^q_A(K_n, \tau_{\geq q - r - 2}K)
$$
pour tout $n$. Il en résulte que
$$
R^q\Gamma_I(K) = R^q\Gamma_I(\tau_{\geq q - r - 2}K) =
H^q(\tau_{\geq q - r - 2}K) = H^q(K)
$$
Ainsi le morphisme $R\Gamma_I(K) \to K$ est un isomorphisme.
\end{proof}

\begin{lemma}
\label{lemma-compute-local-cohomology-noetherian}
Soient $A$ un anneau noethérien et $I = (f_1, \ldots, f_r)$ un idéal de $A$.
Posons $Z = V(I) \subset \Spec(A)$. On a des isomorphismes canoniques
$$
R\Gamma_I(A) \to
(A \to \prod\nolimits_{i_0} A_{f_{i_0}} \to
\prod\nolimits_{i_0 < i_1} A_{f_{i_0}f_{i_1}} \to
\ldots \to A_{f_1\ldots f_r}) \to R\Gamma_Z(A)
$$
dans $D(A)$. Si $M$ est un $A$-module, on a de même
$$
R\Gamma_I(M) \cong
(M \to \prod\nolimits_{i_0} M_{f_{i_0}} \to
\prod\nolimits_{i_0 < i_1} M_{f_{i_0}f_{i_1}} \to
\ldots \to M_{f_1\ldots f_r}) \cong R\Gamma_Z(M)
$$
dans $D(A)$.
\end{lemma}

\begin{proof}
Cela résulte du lemme \ref{lemma-local-cohomology-noetherian} et du calcul du foncteur $R\Gamma_Z$
donné dans le lemme \ref{lemma-local-cohomology-adjoint}.
\end{proof}

\begin{lemma}
\label{lemma-local-cohomology-change-rings}
Si $A \to B$ est un homomorphisme d'anneaux noethériens et $I \subset A$
un idéal, alors on a dans $D(B)$
$$
R\Gamma_I(A) \otimes_A^\mathbf{L} B =
R\Gamma_Z(A) \otimes_A^\mathbf{L} B =
R\Gamma_Y(B) = R\Gamma_{IB}(B)
$$
où $Y = V(IB) \subset \Spec(B)$.
\end{lemma}

\begin{proof}
Il suffit de combiner les lemmes \ref{lemma-compute-local-cohomology-noetherian} et \ref{lemma-torsion-change-rings}.
\end{proof}






\section{Profondeur}
\label{section-depth}

\noindent
Dans cette section, nous revenons sur la notion de profondeur introduite
dans Algèbre, section \ref{algebra-section-depth}.

\begin{lemma}
\label{lemma-depth}
Soient $A$ un anneau noethérien, $I \subset A$ un idéal
et $M$ un $A$-module de type fini tel que $IM \not = M$.
Alors les entiers suivants sont égaux :
\begin{enumerate}
\item $\text{profondeur}_I(M)$ ;
\item le plus petit entier $i$ tel que $\Ext_A^i(A/I, M)$ soit non nul ;
\item le plus petit entier $i$ tel que $H^i_I(M)$ soit non nul.
\end{enumerate}
De plus, on a $\Ext^i_A(N, M) = 0$ pour $i < \text{profondeur}_I(M)$
et pour tout $A$-module de type fini $N$ annulé par une puissance de $I$.
\end{lemma}

\begin{proof}
Montrons par récurrence sur $\text{profondeur}_I(M)$ que (1) et (2) sont égaux,
ce qui est permis par Algèbre, lemme \ref{algebra-lemma-depth-finite-noetherian}.

\medskip\noindent
Cas initial. Si $\text{profondeur}_I(M) = 0$, alors $I$ est contenu dans la réunion
des idéaux premiers associés de $M$
(Algèbre, lemme \ref{algebra-lemma-ass-zero-divisors}).
Par évitement des idéaux premiers (Algèbre, lemme \ref{algebra-lemma-silly}),
on voit que $I \subset \mathfrak p$ pour un idéal premier associé $\mathfrak p$.
Ainsi $\Hom_A(A/I, M)$ est non nul. L'égalité est donc vraie dans ce cas.

\medskip\noindent
Supposons que $\text{profondeur}_I(M) > 0$. Soit $f \in I$ un non-diviseur de zéro
sur $M$ tel que $\text{profondeur}_I(M/fM) = \text{profondeur}_I(M) - 1$.
Considérons la suite exacte courte
$$
0 \to M \to M \to M/fM \to 0
$$
et la suite exacte longue associée pour $\Ext^*_A(A/I, -)$.
Notons que $\Ext^i_A(A/I, M)$ est un $A/I$-module de type fini
(Algèbre, lemmes \ref{algebra-lemma-ext-noetherian} et \ref{algebra-lemma-annihilate-ext}). On obtient donc
$$
\Hom_A(A/I, M/fM) = \Ext^1_A(A/I, M)
$$
et des suites exactes courtes
$$
0 \to \Ext^i_A(A/I, M) \to \text{Ext}^i_A(A/I, M/fM) \to
\Ext^{i + 1}_A(A/I, M) \to 0
$$
D'où l'égalité de (1) et (2) par récurrence.

\medskip\noindent
Observons que $\text{profondeur}_I(M) = \text{profondeur}_{I^n}(M)$ pour tout $n \geq 1$,
par exemple d'après Algèbre, lemme \ref{algebra-lemma-regular-sequence-powers}.
L'égalité de (1) et (2) donne donc $\Ext^i_A(A/I^n, M) = 0$ pour tous $n$ et $i < \text{profondeur}_I(M)$.
Soit $N$ un $A$-module de type fini annulé par une puissance de $I$.
On peut alors choisir une suite exacte courte
$$
0 \to N' \to (A/I^n)^{\oplus m} \to N \to 0
$$
pour certains $n, m \geq 0$. Alors
$\Hom_A(N, M) \subset \Hom_A((A/I^n)^{\oplus m}, M)$
et
$\Ext^i_A(N, M) \subset \text{Ext}^{i - 1}_A(N', M)$
pour $i < \text{profondeur}_I(M)$. Un simple raisonnement par récurrence
établit donc la dernière assertion du lemme.

\medskip\noindent
Enfin, prouvons que (3) est égal à (1) et (2).
On a $H^p_I(M) = \colim \Ext^p_A(A/I^n, M)$ d'après le lemme \ref{lemma-local-cohomology-ext}.
On voit donc que $H^i_I(M) = 0$ pour $i < \text{profondeur}_I(M)$.
Pour $i = \text{profondeur}_I(M)$, l'annulation de $\Ext_A^{i - 1}(I/I^n, M)$ montre que le morphisme
$\Ext_A^i(A/I, M) \to H_I^i(M)$ est injectif, ce qui prouve la non-annulation au degré voulu.
\end{proof}

\begin{lemma}
\label{lemma-depth-in-ses}
Soit $A$ un anneau noethérien. Soit $0 \to N' \to N \to N'' \to 0$
une suite exacte courte de $A$-modules de type fini.
Soit $I \subset A$ un idéal.
\begin{enumerate}
\item
$\text{profondeur}_I(N) \geq \min\{\text{profondeur}_I(N'), \text{profondeur}_I(N'')\}$
\item
$\text{profondeur}_I(N'') \geq \min\{\text{profondeur}_I(N), \text{profondeur}_I(N') - 1\}$
\item
$\text{profondeur}_I(N') \geq \min\{\text{profondeur}_I(N), \text{profondeur}_I(N'') + 1\}$
\end{enumerate}
\end{lemma}

\begin{proof}
Supposons $IN \not = N$, $IN' \not = N'$ et $IN'' \not = N''$.
On peut alors utiliser la caractérisation de la profondeur
par les groupes Ext $\Ext^i(A/I, N)$, voir le lemme \ref{lemma-depth},
ainsi que la suite exacte longue de cohomologie
$$
\begin{matrix}
0
\to \Hom_A(A/I, N')
\to \Hom_A(A/I, N)
\to \Hom_A(A/I, N'')
\\
\phantom{0\ }
\to \Ext^1_A(A/I, N')
\to \Ext^1_A(A/I, N)
\to \Ext^1_A(A/I, N'')
\to \ldots
\end{matrix}
$$
donnée par Algèbre, lemme \ref{algebra-lemma-long-exact-seq-ext}.
Cet argument fonctionne aussi si $IN = N$,
car alors $\Ext^i_A(A/I, N) = 0$ pour tout $i$.
De même dans le cas où $IN' \not = N'$ et/ou $IN'' \not = N''$.
\end{proof}

\begin{lemma}
\label{lemma-depth-drops-by-one}
Soient $A$ un anneau noethérien, $I \subset A$ un idéal et $M$
un $A$-module de type fini tel que $IM \not = M$.
\begin{enumerate}
\item Si $x \in I$ est un non-diviseur de zéro sur $M$, alors
$\text{profondeur}_I(M/xM) = \text{profondeur}_I(M) - 1$.
\item Toute suite $M$-régulière $x_1, \ldots, x_r$ dans $I$
peut être prolongée en une suite $M$-régulière dans $I$
de longueur $\text{profondeur}_I(M)$.
\end{enumerate}
\end{lemma}

\begin{proof}
L'assertion (2) est une conséquence formelle de (1).
Soit $x \in I$ comme dans (1). La suite exacte courte $0 \to M \to M \to M/xM \to 0$
et le lemme \ref{lemma-depth-in-ses} donnent $\text{profondeur}_I(M/xM) \geq \text{profondeur}_I(M) - 1$.
D'autre part, si $x_1, \ldots, x_r \in I$ est une suite régulière pour $M/xM$,
alors $x, x_1, \ldots, x_r$ est une suite régulière pour $M$. D'où (1).
\end{proof}

\begin{lemma}
\label{lemma-depth-CM}
Soit $R$ un anneau local noethérien. Si $M$ est un
$R$-module de Cohen-Macaulay de type fini et $I \subset R$ un idéal non trivial, alors
$$
\text{profondeur}_I(M) = \dim(\text{Supp}(M)) - \dim(\text{Supp}(M/IM)).
$$
\end{lemma}

\begin{proof}
Nous allons le démontrer par récurrence sur $\text{profondeur}_I(M)$.

\medskip\noindent
Si $\text{profondeur}_I(M) = 0$, alors $I$ est contenu dans l'un des idéaux premiers
associés $\mathfrak p$ de $M$ (Algèbre, lemme \ref{algebra-lemma-ideal-nonzerodivisor}).
Alors $\mathfrak p \in \text{Supp}(M/IM)$, d'où
$\dim(\text{Supp}(M/IM)) \geq \dim(R/\mathfrak p) = \dim(\text{Supp}(M))$
où l'égalité résulte d'Algèbre, lemme \ref{algebra-lemma-CM-ass-minimal-support}.
Le lemme est donc vrai dans ce cas.

\medskip\noindent
Si $\text{profondeur}_I(M) > 0$, choisissons $x \in I$ qui soit un non-diviseur de zéro sur $M$.
Notons que $(M/xM)/I(M/xM) = M/IM$. D'autre part, on a
$\text{profondeur}_I(M/xM) = \text{profondeur}_I(M) - 1$
d'après le lemme \ref{lemma-depth-drops-by-one}, et $\dim(\text{Supp}(M/xM)) = \dim(\text{Supp}(M)) -  1$
d'après Algèbre, lemme \ref{algebra-lemma-one-equation-module}.
Le résultat découle donc de l'hypothèse de récurrence.
\end{proof}

\begin{lemma}
\label{lemma-depth-flat-CM}
Soit $R \to S$ un homomorphisme local plat d'anneaux locaux noethériens.
Notons $\mathfrak m \subset R$ l'idéal maximal.
Soit $I \subset S$ un idéal.
Si $S/\mathfrak mS$ est de Cohen-Macaulay, alors
$$
\text{profondeur}_I(S) \geq \dim(S/\mathfrak mS) - \dim(S/\mathfrak mS + I)
$$
\end{lemma}

\begin{proof}
D'après Algèbre, lemme \ref{algebra-lemma-grothendieck-regular-sequence}, toute suite dans $S$
dont l'image est une suite régulière dans $S/\mathfrak mS$
est une suite régulière dans $S$.
Il suffit donc de prouver le lemme lorsque $R$ est un corps.
C'est un cas particulier du lemme \ref{lemma-depth-CM}.
\end{proof}

\begin{lemma}
\label{lemma-divide-by-torsion}
Soient $A$ un anneau et $I \subset A$ un idéal de type fini.
Soit $M$ un $A$-module. Soit $Z = V(I)$.
Alors $H^0_I(M) = H^0_Z(M)$. Notons $N$ cette valeur commune
et posons $M' = M/N$. Alors :
\begin{enumerate}
\item $H^0_I(M') = 0$, $H^p_I(M) = H^p_I(M')$ et $H^p_I(N) = 0$
pour tout $p > 0$ ;
\item $H^0_Z(M') = 0$, $H^p_Z(M) = H^p_Z(M')$ et $H^p_Z(N) = 0$
pour tout $p > 0$.
\end{enumerate}
\end{lemma}

\begin{proof}
Par définition, $H^0_I(M) = M[I^\infty]$ est de torsion par les puissances de $I$.
D'après le lemme \ref{lemma-local-cohomology-adjoint}, on voit que
$$
H^0_Z(M) = \Ker(M \longrightarrow M_{f_1} \times \ldots \times M_{f_r})
$$
si $I = (f_1, \ldots, f_r)$. Ainsi $H^0_I(M) \subset H^0_Z(M)$ et, réciproquement, si $x \in H^0_Z(M)$,
cet élément est annulé par un $f_i^{e_i}$ pour un certain $e_i \geq 1$,
donc par une puissance de $I$.
Cela prouve la première égalité ; de plus, $N$ est de torsion
par les puissances de $I$.
D'après le lemme \ref{lemma-adjoint}, on voit que $R\Gamma_I(N) = N$.
D'après le lemme \ref{lemma-local-cohomology-adjoint}, on voit que $R\Gamma_Z(N) = N$.
Cela prouve l'annulation de $H^p_I(N)$ et de $H^p_Z(N)$ en degrés strictement
positifs dans (1) et (2).
L'annulation de $H^0_I(M')$ et de $H^0_Z(M')$ résulte des remarques précédentes
et du fait que $M'$ est sans torsion par les puissances de $I$
d'après Compléments d'algèbre, lemme \ref{more-algebra-lemma-divide-by-torsion}.
L'égalité des cohomologies supérieures de $M$ et de $M'$
découle immédiatement de la suite exacte longue de cohomologie.
\end{proof}









\section{Modules de torsion et modules complets}
\label{section-torsion-and-complete}

\noindent
Soient $A$ un anneau et $I$ un idéal de type fini.
On peut alors considérer la catégorie dérivée $D_{I^\infty\text{-torsion}}(A)$
des complexes dont les modules de cohomologie sont de torsion
par les puissances de $I$ (section \ref{section-local-cohomology})
et la catégorie dérivée $D_{comp}(A, I)$ des complexes complets au sens dérivé
(Compléments d'algèbre, section \ref{more-algebra-section-derived-completion}).
Dans cette section, nous montrons que ces catégories sont équivalentes.
On trouvera un énoncé plus général dans \cite{Dwyer-Greenlees}.

\begin{lemma}
\label{lemma-complete-and-local}
\begin{slogan}
Les résultats de cette nature sont parfois appelés dualité de Greenlees-May.
\end{slogan}
Soient $A$ un anneau et $I$ un idéal de type fini.
Soit $R\Gamma_Z$ comme dans le lemme \ref{lemma-local-cohomology-adjoint}.
Notons ${\ }^\wedge$ la complétion dérivée, comme dans
Compléments d'algèbre, lemme \ref{more-algebra-lemma-derived-completion}.
Pour un objet $K$ de $D(A)$, on a
$$
R\Gamma_Z(K^\wedge) = R\Gamma_Z(K)
\quad\text{et}\quad
(R\Gamma_Z(K))^\wedge = K^\wedge
$$
dans $D(A)$.
\end{lemma}

\begin{proof}
Choisissons $f_1, \ldots, f_r \in A$ engendrant $I$. Rappelons que
$$
K^\wedge = R\Hom_A\left((A \to \prod A_{f_{i_0}}
\to \prod A_{f_{i_0i_1}} \to \ldots \to A_{f_1 \ldots f_r}), K\right)
$$
d'après Compléments d'algèbre, lemme \ref{more-algebra-lemma-derived-completion}.
Le cône $C = \text{Cone}(K \to K^\wedge)$ est donc donné par
$$
R\Hom_A\left((\prod A_{f_{i_0}}
\to \prod A_{f_{i_0i_1}} \to \ldots \to A_{f_1 \ldots f_r}), K\right)
$$
qui peut être représenté par un complexe muni d'une filtration finie
dont les quotients successifs sont isomorphes aux
$$
R\Hom_A(A_{f_{i_0} \ldots f_{i_p}}, K), \quad p > 0
$$
Ces complexes s'annulent lorsqu'on leur applique $R\Gamma_Z$,
voir le lemme \ref{lemma-local-cohomology-vanishes}. En appliquant $R\Gamma_Z$
au triangle distingué $K \to K^\wedge \to C \to K[1]$,
on voit que la première formule du lemme est correcte.

\medskip\noindent
Rappelons que
$$
R\Gamma_Z(K) =
K \otimes^\mathbf{L} (A \to \prod A_{f_{i_0}}
\to \prod A_{f_{i_0i_1}} \to \ldots \to A_{f_1 \ldots f_r})
$$
d'après le lemme \ref{lemma-local-cohomology-adjoint}.
Le cône $C = \text{Cone}(R\Gamma_Z(K) \to K)$ peut donc être représenté par un complexe muni
d'une filtration finie dont les quotients successifs sont isomorphes aux
$$
K \otimes_A A_{f_{i_0} \ldots f_{i_p}}, \quad p > 0
$$
Ces complexes s'annulent lorsqu'on leur applique ${\ }^\wedge$,
voir Compléments d'algèbre, lemme \ref{more-algebra-lemma-derived-completion-vanishes}.
En appliquant la complétion dérivée au triangle distingué
$R\Gamma_Z(K) \to K \to C \to R\Gamma_Z(K)[1]$
on voit que la seconde formule du lemme est correcte.
\end{proof}

\noindent
Le résultat suivant est un cas particulier d'un phénomène très général
concernant les sous-catégories admissibles d'une catégorie triangulée.

\begin{proposition}
\label{proposition-torsion-complete}
\begin{reference}
C'est un cas particulier de \cite[théorème 1.1]{Porta-Liran-Yekutieli}.
\end{reference}
Soient $A$ un anneau et $I \subset A$ un idéal de type fini.
Les foncteurs $R\Gamma_Z$ et ${\ }^\wedge$ définissent des équivalences
quasi-inverses de catégories
$$
D_{I^\infty\text{-torsion}}(A) \leftrightarrow D_{comp}(A, I)
$$
\end{proposition}

\begin{proof}
Cela résulte immédiatement du lemme \ref{lemma-complete-and-local}.
\end{proof}

\noindent
Le complément suivant à la proposition précédente précise
la correspondance sur les morphismes.

\begin{lemma}
\label{lemma-compare-RHom}
Gardons les notations du lemme \ref{lemma-complete-and-local}.
Pour des objets $K, L$ de $D(A)$, il existe un isomorphisme canonique
$$
R\Hom_A(K^\wedge, L^\wedge) \longrightarrow R\Hom_A(R\Gamma_Z(K), R\Gamma_Z(L))
$$
dans $D(A)$.
\end{lemma}

\begin{proof}
Écrivons $I = (f_1, \ldots, f_r)$. Notons $C = (A \to \prod A_{f_i} \to \ldots \to A_{f_1 \ldots f_r})$ le complexe alterné de
{\v C}ech. La complétion dérivée est alors donnée par
$R\Hom_A(C, -)$ (Compléments d'algèbre, lemme \ref{more-algebra-lemma-derived-completion})
et la cohomologie locale par $C \otimes^\mathbf{L} -$ (lemme \ref{lemma-local-cohomology-adjoint}).
En combinant l'isomorphisme
$$
R\Hom_A(K \otimes^\mathbf{L} C, L \otimes^\mathbf{L} C) =
R\Hom_A(K, R\Hom_A(C,  L \otimes^\mathbf{L} C))
$$
(Compléments d'algèbre, lemme \ref{more-algebra-lemma-internal-hom})
et le morphisme
$$
L \to R\Hom_A(C,  L \otimes^\mathbf{L} C)
$$
(Compléments d'algèbre, lemme \ref{more-algebra-lemma-internal-hom-diagonal}), on obtient un morphisme
$$
\gamma :
R\Hom_A(K, L)
\longrightarrow
R\Hom_A(K \otimes^\mathbf{L} C, L \otimes^\mathbf{L} C)
$$
D'autre part, le membre de droite est complet au sens dérivé,
car il est égal à
$$
R\Hom_A(C, R\Hom_A(K, L \otimes^\mathbf{L} C)).
$$
Ainsi $\gamma$ se factorise par le complété dérivé de $R\Hom_A(K, L)$
en vertu de la propriété universelle de la complétion dérivée.
Or, la complétion dérivée passe à l'intérieur du $R\Hom_A$
d'après Compléments d'algèbre, lemme \ref{more-algebra-lemma-completion-RHom},
ce qui donne le morphisme cherché.

\medskip\noindent
Pour montrer que le morphisme du lemme est un isomorphisme,
on peut supposer que $K$ et $L$ sont complets au sens dérivé,
c'est-à-dire que $K = K^\wedge$ et $L = L^\wedge$. Il s'agit alors du morphisme
$$
\gamma : R\Hom_A(K, L) \longrightarrow R\Hom_A(R\Gamma_Z(K), R\Gamma_Z(L))
$$
D'après la proposition \ref{proposition-torsion-complete}, les groupes de cohomologie
des membres de gauche et de droite coïncident.
Autrement dit, il faut vérifier que le morphisme $\gamma$
envoie un morphisme $\alpha : K \to L$ dans $D(A)$ sur le morphisme
$R\Gamma_Z(\alpha) : R\Gamma_Z(K) \to R\Gamma_Z(L)$.
Nous omettons cette vérification (indication : noter que $R\Gamma_Z(\alpha)$
est simplement le morphisme $\alpha \otimes \text{id}_C :
K \otimes^\mathbf{L} C
\to
L \otimes^\mathbf{L} C$, ce qui est presque identique
à la construction du morphisme dans Compléments d'algèbre, lemme \ref{more-algebra-lemma-internal-hom-diagonal}).
\end{proof}

\begin{lemma}
\label{lemma-completion-local}
Soient $I$ et $J$ des idéaux d'un anneau noethérien $A$.
Soit $M$ un $A$-module de type fini. Posons $Z =V(J)$.
Considérons le complété $I$-adique dérivé $R\Gamma_Z(M)^\wedge$
de la cohomologie locale. Alors :
\begin{enumerate}
\item on a $R\Gamma_Z(M)^\wedge = R\lim R\Gamma_Z(M/I^nM)$ ;
\item on a des suites exactes courtes
$$
0 \to R^1\lim H^{i - 1}_Z(M/I^nM) \to H^i(R\Gamma_Z(M)^\wedge) \to
\lim H^i_Z(M/I^nM) \to 0
$$
\end{enumerate}
En particulier, $R\Gamma_Z(M)^\wedge$ a une cohomologie nulle en degrés négatifs.
\end{lemma}

\begin{proof}
Supposons que $J = (g_1, \ldots, g_m)$.
Alors $R\Gamma_Z(M)$ est calculé par le complexe
$$
M \to \prod M_{g_{j_0}} \to \prod M_{g_{j_0}g_{j_1}} \to
\ldots \to M_{g_1g_2\ldots g_m}
$$
d'après le lemme \ref{lemma-local-cohomology-adjoint}.
D'après Compléments d'algèbre, lemme \ref{more-algebra-lemma-when-derived-completion-is-completion},
le complété $I$-adique dérivé de ce complexe
est donné par le complexe
$$
\lim M/I^nM \to \prod \lim (M/I^nM)_{g_{j_0}} \to
\ldots \to \lim (M/I^nM)_{g_1g_2\ldots g_m}
$$
des complétés usuels. Puisque $R\Gamma_Z(M/I^nM)$ est calculé par le complexe
$ M/I^nM \to \prod (M/I^nM)_{g_{j_0}} \to
\ldots \to (M/I^nM)_{g_1g_2\ldots g_m}$ et que les morphismes de transition entre ces complexes
sont surjectifs, on en déduit (1) par
Compléments d'algèbre, lemme \ref{more-algebra-lemma-compute-Rlim-modules}.
L'assertion (2) découle alors de Compléments d'algèbre, lemme \ref{more-algebra-lemma-break-long-exact-sequence-modules}.
\end{proof}

\begin{lemma}
\label{lemma-completion-local-H0}
Sous les notations et hypothèses du lemme \ref{lemma-completion-local},
supposons $A$ complet pour la topologie $I$-adique. Alors
$$
H^0(R\Gamma_Z(M)^\wedge) = \colim H^0_{V(J')}(M)
$$
où la limite inductive filtrante porte sur les $J' \subset J$ tels que
$V(J') \cap V(I) = V(J) \cap V(I)$.
\end{lemma}

\begin{proof}
Puisque $M$ est un $A$-module de type fini, $M$
est complet pour la topologie $I$-adique.
La démonstration du lemme \ref{lemma-completion-local} montre que
$$
H^0(R\Gamma_Z(M)^\wedge) =
\Ker(M^\wedge \to \prod M_{g_j}^\wedge) =
\Ker(M \to \prod M_{g_j}^\wedge)
$$
où le membre de droite fait intervenir la complétion $I$-adique usuelle.
Le noyau $K_j$ de $M_{g_j} \to M_{g_j}^\wedge$ est $\bigcap I^n M_{g_j}$.
D'après Algèbre, lemme \ref{algebra-lemma-intersection-powers-ideal-module},
pour tout $\mathfrak p \in V(IA_{g_j})$, il existe $f \in A_{g_j}$ tel que $f \not \in \mathfrak p$ et $(K_j)_f = 0$.

\medskip\noindent
Soit $s \in H^0(R\Gamma_Z(M)^\wedge)$.
D'après ce qui précède, on peut considérer $s$ comme un élément de $M$.
L'intersection du support $Z'$ de $s$ avec $D(g_j)$
est disjointe de $D(g_j) \cap V(I)$ d'après les arguments ci-dessus.
Ainsi $Z'$ est un fermé de $\Spec(A)$ tel que $Z' \cap V(I) \subset V(J)$.
On a donc $Z' \cup V(J) = V(J')$ pour un idéal $J' \subset J$ tel que $V(J') \cap V(I) \subset V(J)$,
et l'on a $s \in H^0_{V(J')}(M)$.
Réciproquement, tout $s \in H^0_{V(J')}(M)$ avec $J' \subset J$ et $V(J') \cap V(I) \subset V(J)$
a une image nulle dans $M_{g_j}^\wedge$ pour tout $j$.
Cela prouve le lemme.
\end{proof}









\section{Dualité triviale pour un homomorphisme d'anneaux}
\label{section-trivial}

\noindent
Soit $A \to B$ un homomorphisme d'anneaux. Considérons le foncteur
$$
\Hom_A(B, -) : \text{Mod}_A \longrightarrow \text{Mod}_B,\quad
M \longmapsto \Hom_A(B, M)
$$
Ce foncteur est exact à gauche et possède un prolongement dérivé
$R\Hom(B, -) : D(A) \to D(B)$.

\begin{lemma}
\label{lemma-right-adjoint}
Soit $A \to B$ un homomorphisme d'anneaux. Le foncteur $R\Hom(B, -)$
construit ci-dessus est adjoint à droite du foncteur de restriction
$D(B) \to D(A)$.
\end{lemma}

\begin{proof}
Cela résulte de ce que la restriction et $\Hom_A(B, -)$ sont des foncteurs adjoints
d'après Algèbre, lemme \ref{algebra-lemma-adjoint-hom-restrict}.
Voir Catégories dérivées, lemme \ref{derived-lemma-derived-adjoint-functors}.
\end{proof}

\begin{lemma}
\label{lemma-composition-right-adjoints}
Soient $A \to B \to C$ des homomorphismes d'anneaux. Alors
$R\Hom(C, -) \circ R\Hom(B, -) : D(A) \to D(C)$
est le foncteur $R\Hom(C, -) : D(A) \to D(C)$.
\end{lemma}

\begin{proof}
Cela résulte de l'unicité des adjoints à droite et du lemme \ref{lemma-right-adjoint}.
\end{proof}

\begin{lemma}
\label{lemma-RHom-ext}
Soit $\varphi : A \to B$ un homomorphisme d'anneaux. Pour $K$ dans $D(A)$, on a
$$
\varphi_*R\Hom(B, K) = R\Hom_A(B, K)
$$
où $\varphi_* : D(B) \to D(A)$ est la restriction. En particulier,
$R^q\Hom(B, K) = \Ext_A^q(B, K)$.
\end{lemma}

\begin{proof}
Choisissons un complexe K-injectif $I^\bullet$ représentant $K$.
Alors $R\Hom(B, K)$ est représenté par le complexe $\Hom_A(B, I^\bullet)$ de $B$-modules.
Puisque ce complexe, considéré comme complexe de $A$-modules,
représente $R\Hom_A(B, K)$, le lemme est vrai.
\end{proof}

\noindent
Soit $A$ un anneau noethérien. Nous noterons
$$
D_{\textit{Coh}}(A) \subset D(A)
$$
la sous-catégorie pleine formée des objets $K$ de $D(A)$
dont les modules de cohomologie sont tous des $A$-modules de type fini.
Cela a un sens d'après Catégories dérivées, section \ref{derived-section-triangulated-sub},
car, puisque $A$ est noethérien, la sous-catégorie des
$A$-modules de type fini est une sous-catégorie de Serre de $\text{Mod}_A$.

\begin{lemma}
\label{lemma-exact-support-coherent}
Avec les notations ci-dessus, supposons que $A \to B$ soit un homomorphisme
fini d'anneaux noethériens. Alors $R\Hom(B, -)$ envoie $D^+_{\textit{Coh}}(A)$ dans $D^+_{\textit{Coh}}(B)$.
\end{lemma}

\begin{proof}
Il faut montrer que, si $K \in D^+(A)$ a des modules de cohomologie de type fini,
le complexe $R\Hom(B, K)$ a lui aussi des modules de cohomologie de type fini.
Cela résulte par exemple du lemme \ref{lemma-RHom-ext}, si l'on montre que
les modules Ext $\Ext^i_A(B, K)$ sont des $A$-modules de type fini.
Puisque $K$ est borné inférieurement, on a une suite spectrale convergente
$$
\Ext^p_A(B, H^q(K)) \Rightarrow \text{Ext}^{p + q}_A(B, K)
$$
Cela achève la démonstration, car les modules $\Ext^p_A(B, H^q(K))$
sont de type fini d'après Algèbre, lemme \ref{algebra-lemma-ext-noetherian}.
\end{proof}

\begin{remark}
\label{remark-exact-support}
Soient $A$ un anneau et $I \subset A$ un idéal. Posons $B = A/I$.
Dans ce cas, le foncteur $\Hom_A(B, -)$ est égal au foncteur
$$
\text{Mod}_A \longrightarrow \text{Mod}_B,\quad M \longmapsto M[I]
$$
qui associe à $M$ son sous-module de $I$-torsion.
\end{remark}

\begin{situation}
\label{situation-resolution}
Soit $R \to A$ un homomorphisme d'anneaux.
Donnons une autre construction de $R\Hom(A, -)$, qui nous sera utile
dans la suite de ce chapitre.
Supposons donnée une algèbre différentielle graduée $(E, d)$
sur $R$ et un quasi-isomorphisme $E \to A$, où l'on considère $A$
comme une algèbre différentielle graduée sur $R$ à différentielle nulle.
On a alors des diagrammes commutatifs
$$
\vcenter{
\xymatrix{
D(E, \text{d}) \ar[rd] & &  D(A) \ar[ll] \ar[ld] \\
& D(R)
}
}
\quad\text{et}\quad
\vcenter{
\xymatrix{
D(E, \text{d}) \ar[rr]_{- \otimes_E^\mathbf{L} A} & &  D(A) \\
& D(R) \ar[lu]^{- \otimes_R^\mathbf{L} E} \ar[ru]_{- \otimes_R^\mathbf{L} A}
}
}
$$
où les flèches horizontales sont des équivalences de catégories
(Algèbre différentielle graduée, lemme \ref{dga-lemma-qis-equivalence}).
La commutativité du premier diagramme est claire.
Le second est commutatif parce que le premier l'est et que nos foncteurs
en sont les adjoints à gauche
(Algèbre différentielle graduée, exemple \ref{dga-example-map-hom-tensor}),
ou encore parce que l'on a $E \otimes^\mathbf{L}_E A = E \otimes_E A$ et que l'on peut appliquer
Algèbre différentielle graduée, lemme \ref{dga-lemma-compose-tensor-functors-general}.
\end{situation}

\begin{lemma}
\label{lemma-RHom-dga}
Dans la situation \ref{situation-resolution}, le foncteur $R\Hom(A, -)$ est égal au composé de
$R\Hom(E, -) : D(R) \to D(E, \text{d})$
et de l'équivalence $- \otimes^\mathbf{L}_E A : D(E, \text{d}) \to D(A)$.
\end{lemma}

\begin{proof}
Cela résulte de ce que $R\Hom(E, -)$ est adjoint à droite
de $- \otimes^\mathbf{L}_R E$ ; voir Algèbre différentielle graduée, lemme \ref{dga-lemma-tensor-hom-adjoint}.
Ce foncteur joue donc le même rôle que le foncteur $R\Hom(A, -)$
pour l'homomorphisme $R \to A$ (lemme \ref{lemma-right-adjoint}).
Ces foncteurs doivent donc se correspondre par l'équivalence
$- \otimes^\mathbf{L}_E A : D(E, \text{d}) \to D(A)$.
\end{proof}

\begin{lemma}
\label{lemma-RHom-is-tensor}
Dans la situation \ref{situation-resolution}, supposons que
\begin{enumerate}
\item $E$, considéré comme objet de $D(R)$, soit compact ;
\item $N = \Hom^\bullet_R(E^\bullet, R)$ calcule $R\Hom(E, R)$.
\end{enumerate}
Alors $R\Hom(E, -) : D(R) \to D(E)$ est isomorphe à
$K \mapsto K \otimes_R^\mathbf{L} N$.
\end{lemma}

\begin{proof}
C'est un cas particulier d'Algèbre différentielle graduée, lemme \ref{dga-lemma-RHom-is-tensor}.
\end{proof}

\begin{lemma}
\label{lemma-RHom-is-tensor-special}
Dans la situation \ref{situation-resolution}, supposons que $A$ soit un $R$-module parfait.
Alors
$$
R\Hom(A, -) : D(R) \to D(A)
$$
est donné par $K \mapsto K \otimes_R^\mathbf{L} M$,
où $M = R\Hom(A, R) \in D(A)$.
\end{lemma}

\begin{proof}
Appliquons Algèbre à puissances divisées, lemme \ref{dpa-lemma-tate-resoluton-pseudo-coherent-ring-map},
pour choisir une résolution de Tate $(E, \text{d})$ de $A$ sur $R$.
Notons que $E^i = 0$ pour $i > 0$, que $E^0 = R[x_1, \ldots, x_n]$ est une algèbre de polynômes
et que $E^i$ est un $E^0$-module libre de type fini pour $i < 0$.
Il en résulte que $E$, considéré comme complexe de $R$-modules,
est un complexe borné supérieurement de $R$-modules libres.
Vérifions les hypothèses du lemme \ref{lemma-RHom-is-tensor}.
La première est satisfaite puisque $A$ est parfait
(donc compact d'après Compléments d'algèbre, proposition \ref{more-algebra-proposition-perfect-is-compact}),
et la seconde d'après Compléments d'algèbre, lemme \ref{more-algebra-lemma-RHom-out-of-projective}.
Le lemme montre que $K \mapsto R\Hom(E, K)$ est isomorphe à $K \mapsto K \otimes_R^\mathbf{L} N$
pour un certain $E$-module différentiel gradué $N$.
Observons que
$$
(R \otimes_R E) \otimes_E^\mathbf{L} A = R \otimes_E E \otimes_E A
$$
dans $D(A)$. D'après Algèbre différentielle graduée, lemme \ref{dga-lemma-compose-tensor-functors-general-algebra},
on en déduit que le composé de $- \otimes_R^\mathbf{L} N$ et de $- \otimes_R^\mathbf{L} A$
est de la forme $- \otimes_R M$ pour un certain $M \in D(A)$.
Pour achever la démonstration, appliquons le lemme \ref{lemma-RHom-dga}.
\end{proof}

\begin{lemma}
\label{lemma-compute-for-effective-Cartier-algebraic}
Soit $R \to A$ un homomorphisme surjectif d'anneaux dont le noyau $I$
est un $R$-module inversible. Le foncteur
$R\Hom(A, -) : D(R) \to D(A)$
est isomorphe à $K \mapsto K \otimes_R^\mathbf{L} N[-1]$,
où $N$ est l'inverse du $A$-module inversible $I \otimes_R A$.
\end{lemma}

\begin{proof}
Puisque $A$ possède la résolution projective finie
$$
0 \to I \to R \to A \to 0
$$
on voit que $A$ est un $R$-module parfait.
D'après le lemme \ref{lemma-RHom-is-tensor-special}, il suffit de prouver que $R\Hom(A, R)$
est représenté par $N[-1]$ dans $D(A)$.
Cela signifie que $R\Hom(A, R)$ possède un unique module de cohomologie non nul,
à savoir $N$ en degré $1$.
Puisque $\text{Mod}_A \to \text{Mod}_R$ est pleinement fidèle, il suffit de le prouver
après application du foncteur de restriction $i_* : D(A) \to D(R)$.
D'après le lemme \ref{lemma-RHom-ext}, on a
$$
i_*R\Hom(A, R) = R\Hom_R(A, R)
$$
À l'aide de la résolution projective finie ci-dessus, on trouve que
ce dernier est représenté par le complexe $R \to I^{\otimes -1}$ avec $R$
en degré $0$. Le morphisme $R \to I^{\otimes -1}$ est injectif
et son conoyau est $N$.
\end{proof}





\section{Changement de base pour la dualité triviale}
\label{section-base-change-trivial-duality}

\noindent
Dans cette section, considérons un carré cocartésien d'anneaux
$$
\xymatrix{
A \ar[r]_\alpha & A' \\
R \ar[u]^\varphi \ar[r]^\rho & R' \ar[u]_{\varphi'}
}
$$
Autrement dit, on a $A' = A \otimes_R R'$. Si $A$ et $R'$
sont {\bf Tor-indépendants sur} $R$,
il existe un morphisme canonique de changement de base
\begin{equation}
\label{equation-base-change}
R\Hom(A, K) \otimes_A^\mathbf{L} A'
\longrightarrow
R\Hom(A', K \otimes_R^\mathbf{L} R')
\end{equation}
dans $D(A')$, fonctoriel en $K$ dans $D(R)$.
En effet, par l'adjonction du lemme \ref{lemma-right-adjoint},
une telle flèche revient à se donner un morphisme
$$
\varphi'_*\left(R\Hom(A, K) \otimes_A^\mathbf{L} A'\right)
\longrightarrow
K \otimes_R^\mathbf{L} R'
$$
dans $D(R')$, où $\varphi'_* : D(A') \to D(R')$ est le foncteur de restriction.
On peut appliquer Compléments d'algèbre, lemme \ref{more-algebra-lemma-base-change-comparison},
au membre de gauche pour voir que cela revient à se donner un morphisme
$$
\varphi_*(R\Hom(A, K)) \otimes_R^\mathbf{L} R'
\longrightarrow
K \otimes_R^\mathbf{L} R'
$$
dans $D(R')$, où $\varphi_* : D(A) \to D(R)$ est le foncteur de restriction.
On peut choisir pour cela $can \otimes^\mathbf{L} \text{id}_{R'}$,
où $can : \varphi_*(R\Hom(A, K)) \to K$ est la counité de l'adjonction entre $R\Hom(A, -)$ et $\varphi_*$.

\begin{lemma}
\label{lemma-check-base-change-is-iso}
Dans la situation ci-dessus, le morphisme (\ref{equation-base-change})
est un isomorphisme si et seulement si le morphisme
$$
R\Hom_R(A, K) \otimes_R^\mathbf{L} R'
\longrightarrow
R\Hom_R(A, K \otimes_R^\mathbf{L} R')
$$
de Compléments d'algèbre, lemme \ref{more-algebra-lemma-internal-hom-diagonal-better}, est un isomorphisme.
\end{lemma}

\begin{proof}
Pour voir que le morphisme est un isomorphisme, il suffit de le vérifier
après application de $\varphi'_*$.
En appliquant le foncteur $\varphi'_*$ à (\ref{equation-base-change})
et en utilisant l'égalité $A' = A \otimes_R^\mathbf{L} R'$, on obtient le morphisme de changement de base
$R\Hom_R(A, K) \otimes_R^\mathbf{L} R' \to
R\Hom_{R'}(A \otimes_R^\mathbf{L} R', K \otimes_R^\mathbf{L} R')$
pour le Hom dérivé de Compléments d'algèbre, équation (\ref{more-algebra-equation-base-change-RHom}).
En explicitant les membres de gauche et de droite comme dans la démonstration
de Compléments d'algèbre, lemme \ref{more-algebra-lemma-base-change-RHom}, et en utilisant notamment
Compléments d'algèbre, lemme \ref{more-algebra-lemma-upgrade-adjoint-tensor-RHom}, on obtient le résultat voulu.
\end{proof}

\begin{lemma}
\label{lemma-flat-bc-surjection}
Soient $R \to A$ et $R \to R'$ des homomorphismes d'anneaux et $A' = A \otimes_R R'$.
Supposons que
\begin{enumerate}
\item $A$ soit pseudo-cohérent en tant que $R$-module ;
\item $R'$ soit de Tor-dimension finie en tant que $R$-module
(par exemple, $R \to R'$ est plat) ;
\item $A$ et $R'$ soient Tor-indépendants sur $R$.
\end{enumerate}
Alors (\ref{equation-base-change}) est un isomorphisme pour $K \in D^+(R)$.
\end{lemma}

\begin{proof}
Cela résulte du lemme \ref{lemma-check-base-change-is-iso} et de
Compléments d'algèbre, lemme \ref{more-algebra-lemma-internal-hom-evaluate-tensor-isomorphism} (4).
\end{proof}

\begin{lemma}
\label{lemma-bc-surjection}
Soient $R \to A$ et $R \to R'$ des homomorphismes d'anneaux et $A' = A \otimes_R R'$.
Supposons que
\begin{enumerate}
\item $A$ soit parfait en tant que $R$-module ;
\item $A$ et $R'$ soient Tor-indépendants sur $R$.
\end{enumerate}
Alors (\ref{equation-base-change}) est un isomorphisme pour tout $K \in D(R)$.
\end{lemma}

\begin{proof}
Cela résulte du lemme \ref{lemma-check-base-change-is-iso} et de
Compléments d'algèbre, lemme \ref{more-algebra-lemma-internal-hom-evaluate-tensor-isomorphism} (1).
\end{proof}







\section{Complexes dualisants}
\label{section-dualizing}

\noindent
Dans cette section, nous définissons les complexes dualisants pour les anneaux noethériens.

\begin{definition}
\label{definition-dualizing}
Soit $A$ un anneau noethérien. Un {\it complexe dualisant}
est un complexe de $A$-modules $\omega_A^\bullet$ tel que
\begin{enumerate}
\item $\omega_A^\bullet$ soit de dimension injective finie ;
\item $H^i(\omega_A^\bullet)$ soit un $A$-module de type fini pour tout $i$ ;
\item $A \to R\Hom_A(\omega_A^\bullet, \omega_A^\bullet)$
soit un quasi-isomorphisme.
\end{enumerate}
\end{definition}

\noindent
Il faut un peu de temps pour se familiariser avec cette définition.
Il peut être utile de démontrer soi-même certains des lemmes suivants
avant d'en lire les démonstrations.

\begin{lemma}
\label{lemma-finite-ext-into-bounded-injective}
Soit $A$ un anneau noethérien. Soient $K, L \in D_{\textit{Coh}}(A)$,
et supposons $L$ de dimension injective finie.
Alors $R\Hom_A(K, L)$ appartient à $D_{\textit{Coh}}(A)$.
\end{lemma}

\begin{proof}
Choisissons un entier $n$ et considérons le triangle distingué
$$
\tau_{\leq n}K \to K \to \tau_{\geq n + 1}K \to \tau_{\leq n}K[1]
$$
voir Catégories dérivées, remarque \ref{derived-remark-truncation-distinguished-triangle}.
Puisque $L$ est de dimension injective finie, $R\Hom_A(\tau_{\geq n + 1}K, L)$
a une cohomologie nulle en degrés $\geq c - n$ pour une constante $c$.
Ainsi, étant donné $i$, on voit que $\Ext^i_A(K, L) \to \Ext^i_A(\tau_{\leq n}K, L)$
est un isomorphisme pour un certain $n \gg - i$.
En appliquant Catégories dérivées de schémas, lemme \ref{perfect-lemma-coherent-internal-hom},
à $\tau_{\leq n}K$ et $L$, on en déduit que $\Ext^i_A(K, L)$
est un $A$-module de type fini pour tout $i$.
Ainsi $R\Hom_A(K, L)$ est bien un objet de $D_{\textit{Coh}}(A)$.
\end{proof}

\begin{lemma}
\label{lemma-dualizing}
Soit $A$ un anneau noethérien. Si $\omega_A^\bullet$ est un complexe dualisant,
le foncteur
$$
D : K \longmapsto R\Hom_A(K, \omega_A^\bullet)
$$
est une anti-équivalence $D_{\textit{Coh}}(A) \to D_{\textit{Coh}}(A)$
qui échange $D^+_{\textit{Coh}}(A)$ et $D^-_{\textit{Coh}}(A)$
et induit une anti-équivalence
$D^b_{\textit{Coh}}(A) \to D^b_{\textit{Coh}}(A)$.
De plus, $D \circ D$ est isomorphe au foncteur identité.
\end{lemma}

\begin{proof}
Soit $K$ un objet de $D_{\textit{Coh}}(A)$. D'après le lemme \ref{lemma-finite-ext-into-bounded-injective},
on voit que $R\Hom_A(K, \omega_A^\bullet)$ est un objet de $D_{\textit{Coh}}(A)$.
Compléments d'algèbre, lemme \ref{more-algebra-lemma-internal-hom-evaluate-isomorphism-technical},
et les hypothèses sur le complexe dualisant
donnent un isomorphisme canonique
$$
K = R\Hom_A(\omega_A^\bullet, \omega_A^\bullet) \otimes_A^\mathbf{L} K
\longrightarrow
R\Hom_A(R\Hom_A(K, \omega_A^\bullet), \omega_A^\bullet)
$$
Notre foncteur possède donc un quasi-inverse, ce qui achève la démonstration.
\end{proof}

\noindent
Soit $R$ un anneau. Rappelons qu'un objet $L$ de $D(R)$
est {\it inversible} s'il est inversible pour la structure
monoïdale symétrique sur $D(R)$ donnée par le produit tensoriel dérivé.
Dans Compléments d'algèbre, lemme \ref{more-algebra-lemma-invertible-derived}, nous avons vu que cela signifie
que $L$ est parfait, que $L = \bigoplus H^n(L)[-n]$, cette somme étant finie,
que chaque $H^n(L)$ est projectif de type fini
et qu'il existe un recouvrement ouvert $\Spec(R) = \bigcup D(f_i)$ tel que
$L \otimes_R R_{f_i} \cong R_{f_i}[-n_i]$ pour certains entiers $n_i$.

\begin{lemma}
\label{lemma-equivalence-comes-from-invertible}
Soit $A$ un anneau noethérien. Soit $F : D^b_{\textit{Coh}}(A) \to D^b_{\textit{Coh}}(A)$
une équivalence de catégories $A$-linéaire.
Alors $F(A)$ est un objet inversible de $D(A)$.
\end{lemma}

\begin{proof}
Soit $\mathfrak m \subset A$ un idéal maximal de corps résiduel $\kappa$.
Considérons l'objet $F(\kappa)$. Puisque $\kappa = \Hom_{D(A)}(\kappa, \kappa)$,
tous les groupes de cohomologie de $F(\kappa)$ sont annulés par $\mathfrak m$.
On voit aussi que
$$
\Ext^i_A(\kappa, \kappa) = \text{Ext}^i_A(F(\kappa), F(\kappa))
= \Hom_{D(A)}(F(\kappa), F(\kappa)[i])
$$
est nul pour $i < 0$. Supposons $H^a(F(\kappa)) \not = 0$ et $H^b(F(\kappa)) \not = 0$
avec $a$ minimal et $b$ maximal
(en particulier $a \leq b$). Il existe alors un morphisme non nul
$$
F(\kappa) \to H^b(F(\kappa))[-b] \to H^a(F(\kappa))[-b]
\to F(\kappa)[a - b]
$$
dans $D(A)$ (non nul parce qu'il induit un morphisme non nul en cohomologie).
Cela prouve que $b = a$. On en déduit que $F(\kappa) = \kappa[-a]$.

\medskip\noindent
Soit $G$ un quasi-inverse du foncteur $F$.
En raisonnant comme ci-dessus, on trouve un entier $b$ tel que $G(\kappa) = \kappa[-b]$.
En composant, on obtient $a + b = 0$. Soit $E$ un $A$-module de type fini
annulé par une puissance de $\mathfrak m$. Par récurrence sur la longueur de $E$,
on trouve que $G(E) = E'[-b]$ pour un certain $A$-module de type fini $E'$
annulé par une puissance de $\mathfrak m$. Alors $E[-a] = F(E')$.
Considérons ensuite les groupes
$$
\Ext^i_A(A, E') = \text{Ext}^i_A(F(A), F(E')) =
\Hom_{D(A)}(F(A), E[-a + i])
$$
Le membre de gauche est non nul si et seulement si $i = 0$,
auquel cas on obtient $E'$. En appliquant cela à $E = E' = \kappa$
et en utilisant le lemme de Nakayama, on voit que $H^j(F(A))_\mathfrak m$
est nul pour $j > a$ et engendré par $1$ élément pour $j = a$.
D'autre part, si $H^j(F(A))_\mathfrak m$ n'est pas nul pour un certain $j < a$,
il existe un morphisme $F(A) \to E[-a + i]$ pour un certain $i < 0$
et un certain $E$ (Compléments d'algèbre, lemme \ref{more-algebra-lemma-detect-cohomology}),
ce qui est contradictoire.
Ainsi $F(A)_\mathfrak m = M[-a]$ pour un certain $A_\mathfrak m$-module $M$
engendré par $1$ élément. Cependant, puisque
$$
A_\mathfrak m = \Hom_{D(A)}(A, A)_\mathfrak m =
\Hom_{D(A)}(F(A), F(A))_\mathfrak m = \Hom_{A_\mathfrak m}(M, M)
$$
on voit que $M \cong A_\mathfrak m$. On en déduit qu'il existe un élément $f \in A$
tel que $f \not \in \mathfrak m$ et tel que $F(A)_f$ soit isomorphe à $A_f[-a]$.
Cela achève la démonstration.
\end{proof}

\begin{lemma}
\label{lemma-dualizing-unique}
Soit $A$ un anneau noethérien. Si $\omega_A^\bullet$ et $(\omega'_A)^\bullet$
sont des complexes dualisants, alors $(\omega'_A)^\bullet$ est quasi-isomorphe à
$\omega_A^\bullet \otimes_A^\mathbf{L} L$
pour un certain objet inversible $L$ de $D(A)$.
\end{lemma}

\begin{proof}
D'après les lemmes \ref{lemma-dualizing} et \ref{lemma-equivalence-comes-from-invertible}, le foncteur $K \mapsto R\Hom_A(R\Hom_A(K, \omega_A^\bullet), (\omega_A')^\bullet)$
envoie $A$ sur un objet inversible $L$.
Autrement dit, il existe un isomorphisme
$$
L \longrightarrow R\Hom_A(\omega_A^\bullet, (\omega_A')^\bullet)
$$
Puisque $L$ est de Tor-dimension finie, on peut appliquer
Compléments d'algèbre, lemme \ref{more-algebra-lemma-internal-hom-evaluate-isomorphism-technical}, pour voir que
$$
R\Hom_A(\omega_A^\bullet, (\omega'_A)^\bullet) \otimes_A^\mathbf{L} K
\longrightarrow
R\Hom_A(R\Hom_A(K, \omega_A^\bullet), (\omega_A')^\bullet)
$$
est un isomorphisme pour $K$ dans $D^b_{\textit{Coh}}(A)$.
En particulier, en posant $K = \omega_A^\bullet$, on achève la démonstration.
\end{proof}

\begin{lemma}
\label{lemma-dualizing-localize}
Soit $A$ un anneau noethérien. Soit $B = S^{-1}A$ un localisé.
Si $\omega_A^\bullet$ est un complexe dualisant, alors $\omega_A^\bullet \otimes_A B$
est un complexe dualisant pour $B$.
\end{lemma}

\begin{proof}
Soit $\omega_A^\bullet \to I^\bullet$ un quasi-isomorphisme, où $I^\bullet$ est un complexe borné d'injectifs.
Alors $S^{-1}I^\bullet$ est un complexe borné de $B = S^{-1}A$-modules injectifs
(lemme \ref{lemma-localization-injective-modules}) représentant $\omega_A^\bullet \otimes_A B$.
Ainsi $\omega_A^\bullet \otimes_A B$ est de dimension injective finie.
Puisque $H^i(\omega_A^\bullet \otimes_A B) = H^i(\omega_A^\bullet) \otimes_A B$ par platitude de $A \to B$,
on voit que $\omega_A^\bullet \otimes_A B$ a des modules de cohomologie de type fini.
Enfin, le morphisme
$$
B \longrightarrow
R\Hom_A(\omega_A^\bullet \otimes_A B, \omega_A^\bullet \otimes_A B)
$$
est un quasi-isomorphisme, car la formation du Hom interne commute
au changement de base plat dans ce cas ; voir
Compléments d'algèbre, lemme \ref{more-algebra-lemma-base-change-RHom}.
\end{proof}

\begin{lemma}
\label{lemma-dualizing-glue}
Soit $A$ un anneau noethérien. Soient $f_1, \ldots, f_n \in A$ engendrant l'idéal unité.
Si $\omega_A^\bullet$ est un complexe de $A$-modules tel que $(\omega_A^\bullet)_{f_i}$
soit un complexe dualisant pour $A_{f_i}$ pour tout $i$,
alors $\omega_A^\bullet$ est un complexe dualisant pour $A$.
\end{lemma}

\begin{proof}
Considérons le complexe double
$$
\prod\nolimits_{i_0} (\omega_A^\bullet)_{f_{i_0}}
\to
\prod\nolimits_{i_0 < i_1} (\omega_A^\bullet)_{f_{i_0}f_{i_1}}
\to \ldots
$$
Le complexe total associé est quasi-isomorphe à $\omega_A^\bullet$,
par exemple d'après Descente, remarque \ref{descent-remark-standard-covering},
ou Catégories dérivées de schémas, lemme \ref{perfect-lemma-alternating-cech-complex-complex-computes-cohomology}.
Par hypothèse, les complexes $(\omega_A^\bullet)_{f_i}$ sont de dimension injective finie
en tant que complexes de $A_{f_i}$-modules.
Cela entraîne que chacun des complexes $(\omega_A^\bullet)_{f_{i_0} \ldots f_{i_p}}$, $p > 0$,
est de dimension injective finie sur $A_{f_{i_0} \ldots f_{i_p}}$ ; voir le lemme \ref{lemma-localization-injective-modules}.
Cela entraîne à son tour que chacun des complexes $(\omega_A^\bullet)_{f_{i_0} \ldots f_{i_p}}$, $p > 0$,
est de dimension injective finie sur $A$ ; voir le lemme \ref{lemma-injective-flat}.
Ainsi $\omega_A^\bullet$ est de dimension injective finie comme complexe de $A$-modules
(car il peut être représenté par un complexe muni d'une filtration finie
dont les gradués sont de dimension injective finie).
Puisque $H^n(\omega_A^\bullet)_{f_i}$ est un $A_{f_i}$-module de type fini pour chaque $i$,
on voit que $H^i(\omega_A^\bullet)$ est un $A$-module de type fini ; voir Algèbre, lemme \ref{algebra-lemma-cover}.
Enfin, le changement de base dérivé du morphisme $A \to R\Hom_A(\omega_A^\bullet, \omega_A^\bullet)$ à $A_{f_i}$
est le morphisme $A_{f_i} \to R\Hom_A((\omega_A^\bullet)_{f_i}, (\omega_A^\bullet)_{f_i})$ d'après Compléments d'algèbre, lemme \ref{more-algebra-lemma-base-change-RHom}.
On en déduit donc que
$A \to R\Hom_A(\omega_A^\bullet, \omega_A^\bullet)$
est un isomorphisme, ce qui achève la démonstration.
\end{proof}

\begin{lemma}
\label{lemma-dualizing-finite}
Soit $A \to B$ un homomorphisme fini d'anneaux noethériens.
Soit $\omega_A^\bullet$ un complexe dualisant.
Alors $R\Hom(B, \omega_A^\bullet)$ est un complexe dualisant pour $B$.
\end{lemma}

\begin{proof}
Soit $\omega_A^\bullet \to I^\bullet$ un quasi-isomorphisme, où $I^\bullet$ est un complexe borné d'injectifs.
Alors $\Hom_A(B, I^\bullet)$ est un complexe borné de $B$-modules injectifs
(lemme \ref{lemma-hom-injective}) représentant $R\Hom(B, \omega_A^\bullet)$.
Ainsi $R\Hom(B, \omega_A^\bullet)$ est de dimension injective finie.
D'après le lemme \ref{lemma-exact-support-coherent}, c'est un objet de $D_{\textit{Coh}}(B)$.
Enfin, on calcule
$$
\Hom_{D(B)}(R\Hom(B, \omega_A^\bullet), R\Hom(B, \omega_A^\bullet)) =
\Hom_{D(A)}(R\Hom(B, \omega_A^\bullet), \omega_A^\bullet) = B
$$
et, pour $n \not = 0$, on calcule
$$
\Hom_{D(B)}(R\Hom(B, \omega_A^\bullet), R\Hom(B, \omega_A^\bullet)[n]) =
\Hom_{D(A)}(R\Hom(B, \omega_A^\bullet), \omega_A^\bullet[n]) = 0
$$
ce qui prouve la dernière propriété d'un complexe dualisant.
Dans les égalités affichées, la première égalité résulte du lemme \ref{lemma-right-adjoint}
et la seconde du lemme \ref{lemma-dualizing}.
\end{proof}

\begin{lemma}
\label{lemma-dualizing-quotient}
Soit $A \to B$ un homomorphisme surjectif d'anneaux noethériens.
Soit $\omega_A^\bullet$ un complexe dualisant.
Alors $R\Hom(B, \omega_A^\bullet)$ est un complexe dualisant pour $B$.
\end{lemma}

\begin{proof}
C'est un cas particulier du lemme \ref{lemma-dualizing-finite}.
\end{proof}

\begin{lemma}
\label{lemma-dualizing-polynomial-ring}
Soit $A$ un anneau noethérien. Si $\omega_A^\bullet$ est un complexe dualisant,
alors $\omega_A^\bullet \otimes_A A[x]$ est un complexe dualisant pour $A[x]$.
\end{lemma}

\begin{proof}
Posons $B = A[x]$ et $\omega_B^\bullet = \omega_A^\bullet \otimes_A B$.
Le lemme \ref{lemma-injective-dimension-over-polynomial-ring} et Compléments d'algèbre, lemme \ref{more-algebra-lemma-finite-injective-dimension},
montrent que $\omega_B^\bullet$ est de dimension injective finie.
Puisque $H^i(\omega_B^\bullet) = H^i(\omega_A^\bullet) \otimes_A B$ par platitude de $A \to B$, on voit que $\omega_A^\bullet \otimes_A B$
a des modules de cohomologie de type fini.
Enfin, le morphisme
$$
B \longrightarrow R\Hom_B(\omega_B^\bullet, \omega_B^\bullet)
$$
est un quasi-isomorphisme, car la formation du Hom interne commute
au changement de base plat dans ce cas ; voir
Compléments d'algèbre, lemme \ref{more-algebra-lemma-base-change-RHom}.
\end{proof}

\begin{proposition}
\label{proposition-dualizing-essentially-finite-type}
Soit $A$ un anneau noethérien possédant un complexe dualisant.
Alors toute $A$-algèbre essentiellement de type fini sur $A$
possède un complexe dualisant.
\end{proposition}

\begin{proof}
Cela résulte de la combinaison des lemmes \ref{lemma-dualizing-localize},
\ref{lemma-dualizing-quotient} et \ref{lemma-dualizing-polynomial-ring}.
\end{proof}

\begin{lemma}
\label{lemma-find-function}
Soit $A$ un anneau noethérien. Soit $\omega_A^\bullet$ un complexe dualisant.
Soit $\mathfrak m \subset A$ un idéal maximal et posons $\kappa = A/\mathfrak m$. Alors
$R\Hom_A(\kappa, \omega_A^\bullet) \cong \kappa[n]$ pour un certain $n \in \mathbf{Z}$.
\end{lemma}

\begin{proof}
Cela résulte de ce que $R\Hom_A(\kappa, \omega_A^\bullet)$ est un complexe dualisant sur $\kappa$
(lemme \ref{lemma-dualizing-quotient}), de ce que les complexes dualisants sur $\kappa$
sont uniques à décalage près (lemme \ref{lemma-dualizing-unique})
et de ce que $\kappa$ est un complexe dualisant sur $\kappa$.
\end{proof}




\section{Complexes dualisants sur les anneaux locaux}
\label{section-dualizing-local}

\noindent
Dans cette section, $(A, \mathfrak m, \kappa)$ sera un anneau local noethérien
muni d'un complexe dualisant $\omega_A^\bullet$ tel que l'entier $n$
du lemme \ref{lemma-find-function} soit nul. Plus précisément, supposons $R\Hom_A(\kappa, \omega_A^\bullet) = \kappa[0]$.
Nous dirons alors que le complexe dualisant est {\it normalisé}.
Observons qu'un complexe dualisant normalisé est unique à isomorphisme près
et que tout autre complexe dualisant pour $A$ est isomorphe
à un décalé d'un complexe normalisé (lemme \ref{lemma-dualizing-unique}).

\begin{lemma}
\label{lemma-normalized-finite}
Soit $(A, \mathfrak m, \kappa) \to (B, \mathfrak m', \kappa')$ un homomorphisme local fini d'anneaux locaux noethériens.
Soit $\omega_A^\bullet$ un complexe dualisant normalisé.
Alors $\omega_B^\bullet = R\Hom(B, \omega_A^\bullet)$ est un complexe dualisant normalisé pour $B$.
\end{lemma}

\begin{proof}
D'après le lemme \ref{lemma-dualizing-finite}, le complexe $\omega_B^\bullet$ est dualisant pour $B$.
On a
$$
R\Hom_B(\kappa', \omega_B^\bullet) =
R\Hom_B(\kappa', R\Hom(B, \omega_A^\bullet)) =
R\Hom_A(\kappa', \omega_A^\bullet)
$$
d'après le lemme \ref{lemma-right-adjoint}. Puisque $\kappa'$ est isomorphe à une somme directe
finie de copies de $\kappa$ comme $A$-module, et que $\omega_A^\bullet$ est normalisé,
ce complexe n'a de cohomologie qu'en degré $0$.
Ainsi $\omega_B^\bullet$ est lui aussi un complexe dualisant normalisé.
\end{proof}

\begin{lemma}
\label{lemma-normalized-quotient}
Soit $(A, \mathfrak m, \kappa)$ un anneau local noethérien muni d'un complexe dualisant
normalisé $\omega_A^\bullet$. Soit $A \to B$ surjectif.
Alors $\omega_B^\bullet = R\Hom_A(B, \omega_A^\bullet)$ est un complexe dualisant normalisé pour $B$.
\end{lemma}

\begin{proof}
C'est un cas particulier du lemme \ref{lemma-normalized-finite}.
\end{proof}

\begin{lemma}
\label{lemma-equivalence-finite-length}
Soit $(A, \mathfrak m, \kappa)$ un anneau local noethérien. Soit $F$
une auto-équivalence $A$-linéaire de la catégorie des $A$-modules
de longueur finie. Alors $F$ est isomorphe au foncteur identité.
\end{lemma}

\begin{proof}
Puisque $\kappa$ est l'unique objet simple de la catégorie, on a $F(\kappa) \cong \kappa$.
Notre catégorie étant abélienne, $F$ est exact.
Ainsi $F(E)$ a la même longueur que $E$ pour tout module de longueur finie $E$.
Puisque $\Hom(E, \kappa) = \Hom(F(E), F(\kappa)) \cong \Hom(F(E), \kappa)$, le lemme de Nakayama montre que $E$ et $F(E)$
ont le même nombre de générateurs. Ainsi $F(A/\mathfrak m^n)$ est un $A$-module cyclique.
Choisissons un générateur $e \in F(A/\mathfrak m^n)$.
Puisque $F$ est $A$-linéaire, on a $\mathfrak m^n e = 0$.
Le morphisme $A/\mathfrak m^n \to F(A/\mathfrak m^n)$ est nécessairement un isomorphisme, les longueurs étant égales.
Choisissons un élément
$$
e \in \lim F(A/\mathfrak m^n)
$$
dont l'image soit un générateur pour tout $n$ (nous omettons un petit argument).
On obtient alors un système d'isomorphismes $A/\mathfrak m^n \to F(A/\mathfrak m^n)$
compatible avec tous les homomorphismes de $A$-modules $A/\mathfrak m^n \to A/\mathfrak m^{n'}$
(encore par $A$-linéarité de $F$).
Puisque tout module de longueur finie est conoyau d'un morphisme
entre sommes directes de modules cycliques, on obtient l'isomorphisme du lemme.
\end{proof}

\begin{lemma}
\label{lemma-dualizing-finite-length}
Soit $(A, \mathfrak m, \kappa)$ un anneau local noethérien muni d'un complexe dualisant
normalisé $\omega_A^\bullet$. Soit $E$ une enveloppe injective de $\kappa$.
Il existe alors un isomorphisme fonctoriel
$$
R\Hom_A(N, \omega_A^\bullet) = \Hom_A(N, E)[0]
$$
pour $N$ parcourant les $A$-modules de longueur finie.
\end{lemma}

\begin{proof}
Par récurrence sur la longueur de $N$, on voit que $R\Hom_A(N, \omega_A^\bullet)$
est un module de longueur finie placé en degré $0$.
Ainsi $R\Hom_A(-, \omega_A^\bullet)$ induit une anti-équivalence de la catégorie
des modules de longueur finie. Il en est de même de $\Hom_A(-, E)$
d'après la proposition \ref{proposition-matlis} ; par conséquent,
$$
N \longmapsto \Hom_A(R\Hom_A(N, \omega_A^\bullet), E)
$$
est une équivalence comme dans le lemme \ref{lemma-equivalence-finite-length}.
Elle est donc isomorphe au foncteur identité.
Puisque $\Hom_A(-, E)$ appliqué deux fois est l'identité
(proposition \ref{proposition-matlis}), on obtient l'énoncé du lemme.
\end{proof}

\begin{lemma}
\label{lemma-sitting-in-degrees}
Soit $(A, \mathfrak m, \kappa)$ un anneau local noethérien muni d'un complexe dualisant normalisé $\omega_A^\bullet$.
Soit $M$ un $A$-module de type fini et soit $d = \dim(\text{Supp}(M))$. Alors :
\begin{enumerate}
\item si $\Ext^i_A(M, \omega_A^\bullet)$ est non nul, alors $i \in \{-d, \ldots, 0\}$ ;
\item la dimension du support de $\Ext^i_A(M, \omega_A^\bullet)$ est au plus $-i$ ;
\item $\text{profondeur}(M)$ est le plus petit entier $\delta \geq 0$ tel que $\Ext^{-\delta}_A(M, \omega_A^\bullet) \not = 0$.
\end{enumerate}
\end{lemma}

\begin{proof}
Démontrons cela par récurrence sur $d$.
Si $d = 0$, cela résulte du lemme \ref{lemma-dualizing-finite-length} et de la dualité de Matlis
(proposition \ref{proposition-matlis}), qui assure que $\Hom_A(M, E)$ est non nul si $M$ est non nul.

\medskip\noindent
Supposons le résultat vrai pour les modules dont le support est de dimension
$< d$, et supposons $M$ de profondeur $> 0$.
Choisissons $f \in \mathfrak m$ qui soit un non-diviseur de zéro sur $M$
et considérons la suite exacte courte
$$
0 \to M \to M \to M/fM \to 0
$$
Puisque $\dim(\text{Supp}(M/fM)) = d - 1$ (Algèbre, lemme \ref{algebra-lemma-one-equation-module}), on peut appliquer l'hypothèse de récurrence.
En écrivant $E^i = \Ext^i_A(M, \omega_A^\bullet)$ et $F^i = \Ext^i_A(M/fM, \omega_A^\bullet)$, on obtient une suite exacte longue
$$
\ldots \to F^i \to E^i \xrightarrow{f} E^i \to F^{i + 1} \to \ldots
$$
Par récurrence, $E^i/fE^i = 0$ pour $i + 1 \not \in \{-\dim(\text{Supp}(M/fM)), \ldots, -\text{profondeur}(M/fM)\}$.
Le lemme de Nakayama (Algèbre, lemme \ref{algebra-lemma-NAK})
et Algèbre, lemme \ref{algebra-lemma-depth-drops-by-one}, donnent $E^i = 0$ pour $i \not \in \{-\dim(\text{Supp}(M)), \ldots, -\text{profondeur}(M)\}$.
De plus, dans le cas limite $i = - \text{profondeur}(M)$, on en déduit que $E^i$
est non nul puisque $F^{i + 1}$ est non nul par récurrence.
Puisque $E^i/fE^i \subset F^{i + 1}$, on obtient
$$
\dim(\text{Supp}(F^{i + 1})) \geq \dim(\text{Supp}(E^i/fE^i))
\geq \dim(\text{Supp}(E^i)) - 1
$$
(voir le lemme utilisé ci-dessus), ce qui donne aussi la majoration
de la dimension dans (2).

\medskip\noindent
Si $M$ est de profondeur $0$ et si $d > 0$,
posons $N = M[\mathfrak m^\infty]$ et $M' = M/N$ (comparer au lemme \ref{lemma-divide-by-torsion}).
Alors $M'$ est de profondeur $> 0$ et $\dim(\text{Supp}(M')) = d$.
On connaît donc le résultat pour $M'$ et, puisque
$R\Hom_A(N, \omega_A^\bullet) = \Hom_A(N, E)$
(lemme \ref{lemma-dualizing-finite-length}), la suite exacte longue de cohomologie des $\Ext$
donne le résultat pour $M$.
\end{proof}

\begin{remark}
\label{remark-vanishing-for-arbitrary-modules}
Soient $(A, \mathfrak m)$ et $\omega_A^\bullet$ comme dans le lemme \ref{lemma-sitting-in-degrees}.
D'après Compléments d'algèbre, lemme \ref{more-algebra-lemma-injective-amplitude},
on voit que $\omega_A^\bullet$ a son amplitude injective dans $[-d, 0]$,
car l'assertion (3) de ce lemme s'applique.
En particulier, pour tout $A$-module $M$ (pas nécessairement de type fini),
on a $\Ext^i_A(M, \omega_A^\bullet) = 0$ pour $i \not \in \{-d, \ldots, 0\}$.
\end{remark}

\begin{lemma}
\label{lemma-local-CM}
Soit $(A, \mathfrak m, \kappa)$ un anneau local noethérien muni d'un complexe dualisant normalisé $\omega_A^\bullet$.
Soit $M$ un $A$-module de type fini. Les conditions suivantes sont équivalentes :
\begin{enumerate}
\item $M$ est de Cohen-Macaulay ;
\item $\Ext^i_A(M, \omega_A^\bullet)$ est non nul pour au plus un $i$ ;
\item $\Ext^{-i}_A(M, \omega_A^\bullet)$ est nul pour $i \not = \dim(\text{Supp}(M))$.
\end{enumerate}
Notons $CM_d$ la catégorie des $A$-modules de Cohen-Macaulay
de type fini et de profondeur $d$.
Alors $M \mapsto \Ext^{-d}_A(M, \omega_A^\bullet)$ définit une anti-auto-équivalence de $CM_d$.
\end{lemma}

\begin{proof}
Nous utiliserons les résultats du lemme \ref{lemma-sitting-in-degrees} sans autre mention.
Fixons un module de type fini $M$. Si $M$ est de Cohen-Macaulay,
seul $\Ext^{-d}_A(M, \omega_A^\bullet)$ peut être non nul, d'où (1) $\Rightarrow$ (3).
L'implication (3) $\Rightarrow$ (2) est immédiate.
Supposons (2) et notons $N = \Ext^{-\delta}_A(M, \omega_A^\bullet)$ le groupe $\Ext$ non nul, où $\delta = \text{profondeur}(M)$.
Alors, puisque
$$
M[0] = R\Hom_A(R\Hom_A(M, \omega_A^\bullet), \omega_A^\bullet) =
R\Hom_A(N[\delta], \omega_A^\bullet)
$$
(lemme \ref{lemma-dualizing}), on en déduit que $M = \Ext_A^{-\delta}(N, \omega_A^\bullet)$.
Ainsi $\delta \geq \dim(\text{Supp}(M))$. Cependant, sachant aussi que $\delta \leq \dim(\text{Supp}(M))$
(Algèbre, lemme \ref{algebra-lemma-bound-depth}), on en déduit que $M$ est de Cohen-Macaulay.

\medskip\noindent
Pour prouver la dernière assertion, il suffit de montrer que $N = \Ext^{-d}_A(M, \omega_A^\bullet)$
appartient à $CM_d$ pour $M$ dans $CM_d$.
Nous avons vu ci-dessus que $M[0] = R\Hom_A(N[d], \omega_A^\bullet)$, ce qui prouve le résultat voulu
par l'équivalence de (1) et (3).
\end{proof}

\begin{lemma}
\label{lemma-dualizing-artinian}
Soit $(A, \mathfrak m, \kappa)$ un anneau local noethérien muni d'un complexe dualisant normalisé $\omega_A^\bullet$.
Si $\dim(A) = 0$, alors $\omega_A^\bullet \cong E[0]$,
où $E$ est une enveloppe injective du corps résiduel.
\end{lemma}

\begin{proof}
Cela résulte immédiatement du lemme \ref{lemma-dualizing-finite-length}.
\end{proof}

\begin{lemma}
\label{lemma-divide-by-finite-length-ideal}
Soit $(A, \mathfrak m, \kappa)$ un anneau local noethérien muni d'un complexe dualisant normalisé.
Soit $I \subset \mathfrak m$ un idéal de longueur finie. Posons $B = A/I$.
Il existe alors un triangle distingué
$$
\omega_B^\bullet \to \omega_A^\bullet \to \Hom_A(I, E)[0] \to
\omega_B^\bullet[1]
$$
dans $D(A)$, où $E$ est une enveloppe injective de $\kappa$
et $\omega_B^\bullet$ un complexe dualisant normalisé pour $B$.
\end{lemma}

\begin{proof}
Utiliser la suite exacte courte $0 \to I \to A \to B \to 0$
et les lemmes \ref{lemma-dualizing-finite-length} et \ref{lemma-normalized-quotient}.
\end{proof}

\begin{lemma}
\label{lemma-divide-by-nonzerodivisor}
Soit $(A, \mathfrak m, \kappa)$ un anneau local noethérien muni d'un complexe dualisant normalisé $\omega_A^\bullet$.
Soit $f \in \mathfrak m$ un non-diviseur de zéro. Posons $B = A/(f)$.
Il existe alors un triangle distingué
$$
\omega_B^\bullet \to \omega_A^\bullet \to \omega_A^\bullet \to
\omega_B^\bullet[1]
$$
dans $D(A)$, où $\omega_B^\bullet$ est un complexe dualisant normalisé pour $B$.
\end{lemma}

\begin{proof}
Utiliser la suite exacte courte $0 \to A \to A \to B \to 0$ et le lemme \ref{lemma-normalized-quotient}.
\end{proof}

\begin{lemma}
\label{lemma-nonvanishing-generically-local}
Soit $(A, \mathfrak m, \kappa)$ un anneau local noethérien muni d'un complexe dualisant normalisé $\omega_A^\bullet$.
Soit $\mathfrak p$ un idéal premier minimal de $A$ tel que $\dim(A/\mathfrak p) = e$.
Alors $H^i(\omega_A^\bullet)_\mathfrak p$ est non nul si et seulement si $i = -e$.
\end{lemma}

\begin{proof}
Puisque $A_\mathfrak p$ est de dimension zéro, il existe un entier $n > 0$
tel que $\mathfrak p^nA_\mathfrak p$ soit nul. Posons $B = A/\mathfrak p^n$ et $\omega_B^\bullet = R\Hom_A(B, \omega_A^\bullet)$.
Puisque $B_\mathfrak p = A_\mathfrak p$, on voit que
$$
(\omega_B^\bullet)_\mathfrak p =
R\Hom_A(B, \omega_A^\bullet) \otimes_A^\mathbf{L} A_\mathfrak p =
R\Hom_{A_\mathfrak p}(B_\mathfrak p, (\omega_A^\bullet)_\mathfrak p) =
(\omega_A^\bullet)_\mathfrak p
$$
La deuxième égalité résulte de Compléments d'algèbre, lemme \ref{more-algebra-lemma-base-change-RHom}.
D'après le lemme \ref{lemma-normalized-quotient}, on peut remplacer $A$ par $B$.
On a alors $\dim(A) = e$. Ainsi $H^i(\omega_A^\bullet)_\mathfrak p$ ne peut être non nul que si $i = -e$
d'après le lemme \ref{lemma-sitting-in-degrees} (1) et (2).
D'autre part, puisque $(\omega_A^\bullet)_\mathfrak p$ est un complexe dualisant
pour l'anneau non nul $A_\mathfrak p$ (lemme \ref{lemma-dualizing-localize}),
le module restant doit être non nul.
\end{proof}





\section{Complexes dualisants et fonctions de dimension}
\label{section-dimension-function}

\noindent
Nos résultats dans le cas local ont la conséquence suivante :
un anneau noethérien possédant un complexe dualisant est
universellement caténaire et de dimension finie.

\begin{lemma}
\label{lemma-nonvanishing-generically}
Soit $A$ un anneau noethérien. Soit $\mathfrak p$ un idéal premier minimal de $A$.
Alors $H^i(\omega_A^\bullet)_\mathfrak p$ est non nul pour exactement un $i$.
\end{lemma}

\begin{proof}
Le complexe $\omega_A^\bullet \otimes_A A_\mathfrak p$ est un complexe dualisant pour $A_\mathfrak p$
(lemme \ref{lemma-dualizing-localize}). La dimension de $A_\mathfrak p$ est zéro puisque $\mathfrak p$ est minimal.
Le résultat découle donc du lemme \ref{lemma-dualizing-artinian}.
\end{proof}

\noindent
Soient $A$ un anneau noethérien et $\omega_A^\bullet$ un complexe dualisant.
Le lemme \ref{lemma-find-function} permet de définir une fonction
$$
\delta = \delta_{\omega_A^\bullet} : \Spec(A) \longrightarrow \mathbf{Z}
$$
en associant à $\mathfrak p$ l'entier du lemme \ref{lemma-find-function}
pour le complexe dualisant $(\omega_A^\bullet)_\mathfrak p$ sur $A_\mathfrak p$
(lemme \ref{lemma-dualizing-localize}) et le corps résiduel $\kappa(\mathfrak p)$.
Plus précisément, $\delta(\mathfrak p)$ est l'unique entier tel que
$$
(\omega_A^\bullet)_\mathfrak p[-\delta(\mathfrak p)]
$$
soit un complexe dualisant normalisé sur l'anneau local noethérien $A_\mathfrak p$.

\begin{lemma}
\label{lemma-quotient-function}
Soient $A$ un anneau noethérien et $\omega_A^\bullet$ un complexe dualisant.
Soit $A \to B$ un homomorphisme surjectif d'anneaux et soit $\omega_B^\bullet = R\Hom(B, \omega_A^\bullet)$
le complexe dualisant pour $B$ du lemme \ref{lemma-dualizing-quotient}. Alors on a
$$
\delta_{\omega_B^\bullet} = \delta_{\omega_A^\bullet}|_{\Spec(B)}
$$
\end{lemma}

\begin{proof}
Cela résulte de la définition des fonctions et du lemme \ref{lemma-normalized-quotient}.
\end{proof}

\begin{lemma}
\label{lemma-dimension-function}
Soient $A$ un anneau noethérien et $\omega_A^\bullet$ un complexe dualisant.
La fonction $\delta = \delta_{\omega_A^\bullet}$ définie ci-dessus est une fonction de dimension
(Topologie, définition \ref{topology-definition-dimension-function}).
\end{lemma}

\begin{proof}
Soit $\mathfrak p \subset \mathfrak q$ une spécialisation immédiate. Il faut montrer que $\delta(\mathfrak p) = \delta(\mathfrak q) + 1$.
On peut remplacer $A$ par $A/\mathfrak p$, le complexe $\omega_A^\bullet$ par $\omega_{A/\mathfrak p}^\bullet = R\Hom(A/\mathfrak p, \omega_A^\bullet)$,
l'idéal premier $\mathfrak p$ par $(0)$ et l'idéal premier $\mathfrak q$ par $\mathfrak q/\mathfrak p$ ;
voir le lemme \ref{lemma-quotient-function}. On peut donc supposer que $A$ est intègre,
que $\mathfrak p = (0)$ et que $\mathfrak q$ est un idéal premier de hauteur $1$.

\medskip\noindent
Alors $H^i(\omega_A^\bullet)_{(0)}$ est non nul pour exactement un $i$,
disons $i_0$, d'après le lemme \ref{lemma-nonvanishing-generically}.
En fait, $i_0 = -\delta((0))$ parce que $(\omega_A^\bullet)_{(0)}[-\delta((0))]$ est un complexe dualisant normalisé
sur le corps $A_{(0)}$.

\medskip\noindent
D'autre part, $(\omega_A^\bullet)_\mathfrak q[-\delta(\mathfrak q)]$ est un complexe dualisant normalisé pour $A_\mathfrak q$.
D'après le lemme \ref{lemma-nonvanishing-generically-local}, on voit que
$$
H^e((\omega_A^\bullet)_\mathfrak q[-\delta(\mathfrak q)])_{(0)} =
H^{e - \delta(\mathfrak q)}(\omega_A^\bullet)_{(0)}
$$
n'est non nul que pour $e = -\dim(A_\mathfrak q) = -1$.
On en déduit
$$
-\delta((0)) = -1 - \delta(\mathfrak q)
$$
comme voulu.
\end{proof}

\begin{lemma}
\label{lemma-universally-catenary}
Soit $A$ un anneau noethérien possédant un complexe dualisant.
Alors $A$ est universellement caténaire et de dimension finie.
\end{lemma}

\begin{proof}
Puisque $\Spec(A)$ possède une fonction de dimension d'après le lemme \ref{lemma-dimension-function},
il est caténaire ; voir Topologie, lemme \ref{topology-lemma-dimension-function-catenary}.
Ainsi $A$ est caténaire ; voir Algèbre, lemme \ref{algebra-lemma-catenary}.
La proposition \ref{proposition-dualizing-essentially-finite-type} entraîne alors que $A$ est universellement caténaire.

\medskip\noindent
Puisque tout complexe dualisant $\omega_A^\bullet$ appartient à $D^b_{\textit{Coh}}(A)$,
les valeurs de la fonction $\delta_{\omega_A^\bullet}$ aux idéaux premiers minimaux
sont bornées d'après le lemme \ref{lemma-nonvanishing-generically}.
D'autre part, pour un idéal maximal $\mathfrak m$ de corps résiduel $\kappa$,
l'entier $i = -\delta(\mathfrak m)$ est l'unique entier tel que $\Ext_A^i(\kappa, \omega_A^\bullet)$ soit non nul
(lemme \ref{lemma-find-function}).
Puisque $\omega_A^\bullet$ est de dimension injective finie, ces valeurs sont aussi bornées.
La dimension de $A$ étant la valeur maximale de $\delta(\mathfrak p) - \delta(\mathfrak m)$,
où $\mathfrak p \subset \mathfrak m$ parcourt les couples formés d'un idéal premier minimal
et d'un idéal maximal, on voit que la dimension de $\Spec(A)$ est bornée.
\end{proof}

\begin{lemma}
\label{lemma-depth-dualizing-module}
Soit $(A, \mathfrak m, \kappa)$ un anneau local noethérien muni d'un complexe dualisant normalisé $\omega_A^\bullet$.
Soient $d = \dim(A)$ et $\omega_A = H^{-d}(\omega_A^\bullet)$. Alors :
\begin{enumerate}
\item le support de $\omega_A$ est la réunion des composantes irréductibles
de $\Spec(A)$ de dimension $d$ ;
\item $\omega_A$ satisfait à $(S_2)$, voir Algèbre, définition \ref{algebra-definition-conditions}.
\end{enumerate}
\end{lemma}

\begin{proof}
Nous utiliserons le lemme \ref{lemma-sitting-in-degrees} sans autre mention.
D'après le lemme \ref{lemma-nonvanishing-generically-local}, le support de $\omega_A$ contient les composantes
irréductibles de dimension $d$. Soit $\mathfrak p \subset A$ un idéal premier.
D'après le lemme \ref{lemma-dimension-function}, le complexe $(\omega_A^\bullet)_{\mathfrak p}[-\dim(A/\mathfrak p)]$ est un complexe dualisant
normalisé pour $A_\mathfrak p$. Ainsi, si $\dim(A/\mathfrak p) + \dim(A_\mathfrak p) < d$, alors $(\omega_A)_\mathfrak p = 0$.
Cela prouve que le support de $\omega_A$ est la réunion des composantes
irréductibles de dimension $d$, car le complémentaire de cette réunion
est exactement formé des idéaux premiers $\mathfrak p$ de $A$ pour lesquels
$\dim(A/\mathfrak p) + \dim(A_\mathfrak p) < d$, puisque $A$ est caténaire (lemme \ref{lemma-universally-catenary}).
D'autre part, si $\dim(A/\mathfrak p) + \dim(A_\mathfrak p) = d$, alors
$$
(\omega_A)_\mathfrak p =
H^{-\dim(A_\mathfrak p)}\left(
(\omega_A^\bullet)_{\mathfrak p}[-\dim(A/\mathfrak p)] \right)
$$
Pour prouver que $\omega_A$ satisfait à $(S_2)$, il suffit donc de montrer
que la profondeur de $\omega_A$ est au moins $\min(\dim(A), 2)$.
Procédons par récurrence sur $\dim(A)$. Le cas $\dim(A) = 0$ est trivial.

\medskip\noindent
Supposons $\text{profondeur}(A) > 0$. Choisissons un non-diviseur de zéro $f \in \mathfrak m$
et posons $B = A/fA$. Alors $\dim(B) = \dim(A) - 1$ et l'on peut appliquer
l'hypothèse de récurrence à $B$.
D'après le lemme \ref{lemma-divide-by-nonzerodivisor}, la multiplication par $f$ est injective sur $\omega_A$
et l'on obtient $\omega_A/f\omega_A \subset \omega_B$. La profondeur de $\omega_A$ est donc au moins $1$.
Si $\dim(A) > 1$, alors $\dim(B) > 0$ et $\omega_B$ est de profondeur $ > 0$.
Ainsi $\omega_A$ est de profondeur $> 1$, ce qui conclut dans ce cas.

\medskip\noindent
Supposons $\dim(A) > 0$ et $\text{profondeur}(A) = 0$. Soit $I = A[\mathfrak m^\infty]$ et posons $B = A/I$.
Alors $B$ est de profondeur $\geq 1$ et $\omega_A = \omega_B$ d'après le lemme \ref{lemma-divide-by-finite-length-ideal}.
Puisque le résultat a été démontré ci-dessus pour $\omega_B$,
la démonstration est achevée.
\end{proof}





\section{Le théorème de dualité locale}
\label{section-local-duality}

\noindent
Le résultat principal de cette section est dû à Grothendieck.

\begin{lemma}
\label{lemma-local-cohomology-of-dualizing}
Soit $(A, \mathfrak m, \kappa)$ un anneau local noethérien.
Soit $\omega_A^\bullet$ un complexe dualisant normalisé.
Soit $Z = V(\mathfrak m) \subset \Spec(A)$.
Alors $E = R^0\Gamma_Z(\omega_A^\bullet)$ est une enveloppe injective de $\kappa$ et $R\Gamma_Z(\omega_A^\bullet) = E[0]$.
\end{lemma}

\begin{proof}
D'après le lemme \ref{lemma-local-cohomology-noetherian}, on a $R\Gamma_{\mathfrak m} = R\Gamma_Z$. Ainsi
$$
R\Gamma_Z(\omega_A^\bullet) =
R\Gamma_{\mathfrak m}(\omega_A^\bullet) =
\text{hocolim}\ R\Hom_A(A/\mathfrak m^n, \omega_A^\bullet)
$$
d'après le lemme \ref{lemma-local-cohomology-ext}. Soit $E'$ une enveloppe injective du corps résiduel.
D'après le lemme \ref{lemma-dualizing-finite-length}, on peut trouver des isomorphismes
$$
R\Hom_A(A/\mathfrak m^n, \omega_A^\bullet) \cong \Hom_A(A/\mathfrak m^n, E')[0]
$$
compatibles avec les morphismes de transition. Puisque
$E' = \bigcup E'[\mathfrak m^n] = \colim \Hom_A(A/\mathfrak m^n, E')$
d'après le lemme \ref{lemma-union-artinian}, on en déduit que $E \cong E'$
et que tous les autres groupes de cohomologie du complexe $R\Gamma_Z(\omega_A^\bullet)$ sont nuls.
\end{proof}

\begin{remark}
\label{remark-specific-injective-hull}
Soit $(A, \mathfrak m, \kappa)$ un anneau local noethérien muni d'un complexe dualisant normalisé $\omega_A^\bullet$.
D'après le lemme \ref{lemma-local-cohomology-of-dualizing} ci-dessus, $R\Gamma_Z(\omega_A^\bullet)$ est une enveloppe injective
du corps résiduel placée en degré $0$.
Cela donne en fait une « construction » ou une « réalisation »
de l'enveloppe injective un peu plus canonique que le choix
d'une enveloppe injective quelconque.
En effet, un complexe dualisant normalisé est unique à isomorphisme près,
et son groupe d'automorphismes est le groupe des unités de $A$,
tandis qu'une enveloppe injective de $\kappa$ est unique à isomorphisme près,
avec pour groupe d'automorphismes le groupe des unités du complété
$A^\wedge$ de $A$ relativement à $\mathfrak m$.
\end{remark}

\noindent
Voici le résultat principal de cette section.

\begin{theorem}
\label{theorem-local-duality}
Soit $(A, \mathfrak m, \kappa)$ un anneau local noethérien.
Soit $\omega_A^\bullet$ un complexe dualisant normalisé.
Soit $E$ une enveloppe injective du corps résiduel.
Soit $Z = V(\mathfrak m) \subset \Spec(A)$.
Notons ${}^\wedge$ la complétion dérivée relativement à $\mathfrak m$.
Alors
$$
R\Hom_A(K, \omega_A^\bullet)^\wedge \cong R\Hom_A(R\Gamma_Z(K), E[0])
$$
pour $K$ dans $D(A)$.
\end{theorem}

\begin{proof}
Observons que $E[0] \cong R\Gamma_Z(\omega_A^\bullet)$ d'après le lemme \ref{lemma-local-cohomology-of-dualizing}.
D'après Compléments d'algèbre, lemme \ref{more-algebra-lemma-completion-RHom},
la complétion dans le membre de gauche passe à l'intérieur.
Il faut donc prouver
$$
R\Hom_A(K^\wedge, (\omega_A^\bullet)^\wedge)
=
R\Hom_A(R\Gamma_Z(K), R\Gamma_Z(\omega_A^\bullet))
$$
Cela résulte de l'équivalence entre $D_{comp}(A, \mathfrak m)$ et $D_{\mathfrak m^\infty\text{-torsion}}(A)$
donnée dans la proposition \ref{proposition-torsion-complete}.
Plus précisément, c'est un cas particulier du lemme \ref{lemma-compare-RHom}.
\end{proof}

\noindent
Voici un cas particulier du théorème précédent.

\begin{lemma}
\label{lemma-special-case-local-duality}
Soit $(A, \mathfrak m, \kappa)$ un anneau local noethérien.
Soit $\omega_A^\bullet$ un complexe dualisant normalisé.
Soit $E$ une enveloppe injective du corps résiduel.
Soit $K \in D_{\textit{Coh}}(A)$. Alors
$$
\Ext^{-i}_A(K, \omega_A^\bullet)^\wedge =
\Hom_A(H^i_{\mathfrak m}(K), E)
$$
où ${}^\wedge$ désigne la complétion $\mathfrak m$-adique.
\end{lemma}

\begin{proof}
D'après le lemme \ref{lemma-dualizing}, on voit que $R\Hom_A(K, \omega_A^\bullet)$ est un objet de $D_{\textit{Coh}}(A)$.
Il en résulte que les modules de cohomologie du complété dérivé de $R\Hom_A(K, \omega_A^\bullet)$
sont les complétés usuels $\Ext^i_A(K, \omega_A^\bullet)^\wedge$ d'après
Compléments d'algèbre, lemme \ref{more-algebra-lemma-derived-completion-pseudo-coherent}.
D'autre part, on a $R\Gamma_{\mathfrak m} = R\Gamma_Z$ pour $Z = V(\mathfrak m)$ d'après le lemme \ref{lemma-local-cohomology-noetherian}.
De plus, le foncteur $\Hom_A(-, E)$ est exact, donc se factorise par la cohomologie.
Le lemme est ainsi une conséquence du théorème \ref{theorem-local-duality}.
\end{proof}





\section{Modules dualisants}
\label{section-dualizing-module}

\noindent
Si $(A, \mathfrak m, \kappa)$ est un anneau local noethérien et $\omega_A^\bullet$ un complexe dualisant normalisé,
on dit que le module $\omega_A = H^{-\dim(A)}(\omega_A^\bullet)$, décrit dans le lemme \ref{lemma-depth-dualizing-module},
est un {\it module dualisant} pour $A$.
Ce module est un module canonique de $A$.
Il semble généralement admis de définir un {\it module canonique}
pour un anneau local noethérien $(A, \mathfrak m, \kappa)$ comme un $A$-module de type fini
$K$ tel que
$$
\Hom_A(K, E) \cong H^{\dim(A)}_\mathfrak m(A)
$$
où $E$ est une enveloppe injective du corps résiduel.
Un module dualisant est canonique parce que
$$
\Hom_A(H^{\dim(A)}_\mathfrak m(A), E) = (\omega_A)^\wedge
$$
d'après le lemme \ref{lemma-special-case-local-duality} ; en appliquant $\Hom_A(-, E)$, on obtient donc
\begin{align*}
\Hom_A(\omega_A, E)
& =
\Hom_A((\omega_A)^\wedge, E) \\
& =
\Hom_A(\Hom_A(H^{\dim(A)}_\mathfrak m(A), E), E) \\
& = H^{\dim(A)}_\mathfrak m(A)
\end{align*}
La première égalité vient de ce que $E$ est de torsion par les puissances
de $\mathfrak m$, la deuxième de ce qui précède et la troisième de la dualité de Matlis
(proposition \ref{proposition-matlis}).
L'intérêt de la définition du module canonique donnée ci-dessus
est qu'elle a un sens même si $A$ ne possède pas de complexe dualisant.




\section{Anneaux de Cohen-Macaulay}
\label{section-CM}

\noindent
Les modules et anneaux de Cohen-Macaulay ont été étudiés dans
Algèbre, sections \ref{algebra-section-CM} et \ref{algebra-section-CM-ring}.

\begin{lemma}
\label{lemma-depth-in-terms-dualizing-complex}
Soit $(A, \mathfrak m, \kappa)$ un anneau local noethérien muni d'un complexe dualisant normalisé $\omega_A^\bullet$.
Alors $\text{profondeur}(A)$ est égale au plus petit entier $\delta \geq 0$ tel que $H^{-\delta}(\omega_A^\bullet) \not = 0$.
\end{lemma}

\begin{proof}
Cela résulte immédiatement du lemme \ref{lemma-sitting-in-degrees}.
Voici deux autres façons de le voir.

\medskip\noindent
Première variante. Le lemme de Nakayama montre que $\delta$
est le plus petit entier tel que $\Hom_A(H^{-\delta}(\omega_A^\bullet), \kappa) \not = 0$.
Autrement dit, c'est le plus petit entier tel que $\Ext_A^{-\delta}(\omega_A^\bullet, \kappa)$ soit non nul.
Le lemme \ref{lemma-dualizing} et le fait que $\omega_A^\bullet$ soit normalisé montrent qu'il s'agit
du plus petit entier tel que $\Ext_A^\delta(\kappa, A)$ soit non nul.
C'est la profondeur de $A$ d'après Algèbre, lemme \ref{algebra-lemma-depth-ext}.

\medskip\noindent
Deuxième variante. Le théorème de dualité locale
(sous la forme du lemme \ref{lemma-special-case-local-duality}) montre que $\delta$ est le plus petit entier
tel que $H^\delta_\mathfrak m(A)$ soit non nul. C'est la profondeur de $A$
d'après le lemme \ref{lemma-depth}.
\end{proof}

\begin{lemma}
\label{lemma-apply-CM}
Soit $(A, \mathfrak m, \kappa)$ un anneau local noethérien muni d'un complexe dualisant normalisé $\omega_A^\bullet$
et d'un module dualisant $\omega_A = H^{-\dim(A)}(\omega_A^\bullet)$. Les conditions suivantes sont équivalentes :
\begin{enumerate}
\item $A$ est de Cohen-Macaulay ;
\item $\omega_A^\bullet$ est concentré en un seul degré ;
\item $\omega_A^\bullet = \omega_A[\dim(A)]$.
\end{enumerate}
Dans ce cas, $\omega_A$ est un module de Cohen-Macaulay maximal.
\end{lemma}

\begin{proof}
Cela résulte immédiatement du lemme \ref{lemma-local-CM}.
\end{proof}

\begin{lemma}
\label{lemma-has-dualizing-module-CM}
Soit $A$ un anneau noethérien. S'il existe un $A$-module de type fini $\omega_A$
tel que $\omega_A[0]$ soit un complexe dualisant, alors $A$ est de Cohen-Macaulay.
\end{lemma}

\begin{proof}
On peut remplacer $A$ par son localisé en un idéal premier
(lemme \ref{lemma-dualizing-localize} et Algèbre, définition \ref{algebra-definition-ring-CM}).
Dans ce cas, le résultat découle immédiatement du lemme \ref{lemma-apply-CM}.
\end{proof}

\begin{lemma}
\label{lemma-CM-open}
Soit $A$ un anneau noethérien muni d'un complexe dualisant $\omega_A^\bullet$.
Soit $M$ un $A$-module de type fini. Alors
$$
U = \{\mathfrak p \in \Spec(A) \mid M_\mathfrak p\text{ est de Cohen-Macaulay}\}
$$
est un ouvert de $\Spec(A)$ dont l'intersection avec $\text{Supp}(M)$ est dense.
\end{lemma}

\begin{proof}
Si $\mathfrak p$ est un point générique de $\text{Supp}(M)$, alors $\text{profondeur}(M_\mathfrak p) = \dim(M_\mathfrak p) = 0$,
donc $\mathfrak p \in U$. Cela prouve la densité.
Si $\mathfrak p \in U$, on voit que
$$
R\Hom_A(M, \omega_A^\bullet)_\mathfrak p =
R\Hom_{A_\mathfrak p}(M_\mathfrak p, (\omega_A^\bullet)_\mathfrak p)
$$
a un unique module de cohomologie non nul, disons en degré $i_0$,
d'après le lemme \ref{lemma-local-CM}.
Puisque $R\Hom_A(M, \omega_A^\bullet)$ n'a qu'un nombre fini de modules de cohomologie non nuls $H^i$,
et que chacun est un $A$-module de type fini,
il existe $f \in A$ tel que $f \not \in \mathfrak p$ et $(H^i)_f = 0$ pour $i \not = i_0$.
Alors $R\Hom_A(M, \omega_A^\bullet)_f$ a un unique module de cohomologie non nul ;
en reprenant les arguments précédents en sens inverse, on trouve $D(f) \subset U$.
\end{proof}

\begin{lemma}
\label{lemma-CM}
Soit $A$ un anneau noethérien. Si $A$ possède un complexe dualisant $\omega_A^\bullet$, alors
$\{\mathfrak p \in \Spec(A) \mid A_\mathfrak p\text{ est de Cohen-Macaulay}\}$
est un ouvert dense de $\Spec(A)$.
\end{lemma}

\begin{proof}
Conséquence immédiate du lemme \ref{lemma-CM-open} et des définitions.
\end{proof}






\section{Anneaux de Gorenstein}
\label{section-gorenstein}

\noindent
Jusqu'ici, le seul complexe dualisant explicite que nous ayons rencontré
est $\kappa$ sur $\kappa$ pour un corps $\kappa$ ; voir la démonstration du lemme \ref{lemma-find-function}.
D'après la proposition \ref{proposition-dualizing-essentially-finite-type}, cela signifie que toute algèbre de type fini
sur un corps possède un complexe dualisant.
Cependant, il existe des anneaux noethériens (locaux) sans complexe dualisant.
En effet, nous avons vu qu'un anneau possédant un complexe dualisant
est universellement caténaire (lemme \ref{lemma-universally-catenary}),
mais il existe des exemples d'anneaux locaux noethériens non caténaires ;
voir Exemples, section \ref{examples-section-non-catenary-Noetherian-local}.

\medskip\noindent
Néanmoins, beaucoup d'anneaux de la géométrie algébrique possèdent
des complexes dualisants simplement parce qu'ils sont quotients
d'anneaux de Gorenstein. Cette condition est en fait nécessaire et suffisante :
un anneau noethérien possède un complexe dualisant si et seulement s'il est
quotient d'un anneau de Gorenstein de dimension finie.
C'est la conjecture de Sharp (\cite{Sharp}), que l'on trouve dans la littérature
sous la forme \cite[corollaire 1.4]{Kawasaki}. Revenons à notre sujet et définissons les anneaux de Gorenstein.

\begin{definition}
\label{definition-gorenstein}
Anneaux de Gorenstein.
\begin{enumerate}
\item Soit $A$ un anneau local noethérien. On dit que $A$ est
{\it de Gorenstein} si $A[0]$ est un complexe dualisant pour $A$.
\item Soit $A$ un anneau noethérien. On dit que $A$ est
{\it de Gorenstein} si $A_\mathfrak p$ est de Gorenstein
pour tout idéal premier $\mathfrak p$ de $A$.
\end{enumerate}
\end{definition}

\noindent
Cette définition a un sens, car, si $A[0]$ est un complexe dualisant
pour $A$, alors $S^{-1}A[0]$ est un complexe dualisant pour $S^{-1}A$
d'après le lemme \ref{lemma-dualizing-localize}.
Un anneau noethérien de dimension finie est de Gorenstein
s'il est de dimension injective finie comme module sur lui-même ;
voir le lemme \ref{lemma-characterize-gorenstein}.

\begin{lemma}
\label{lemma-gorenstein-CM}
Un anneau de Gorenstein est de Cohen-Macaulay.
\end{lemma}

\begin{proof}
Cela résulte du lemme \ref{lemma-apply-CM}.
\end{proof}

\noindent
Un anneau régulier est un exemple d'anneau de Gorenstein.

\begin{lemma}
\label{lemma-regular-gorenstein}
Un anneau local régulier est de Gorenstein.
Un anneau régulier est de Gorenstein.
\end{lemma}

\begin{proof}
Soit $A$ un anneau régulier de dimension finie $d$.
Alors $A$ est de dimension globale finie $d$ ; voir Algèbre, lemme \ref{algebra-lemma-finite-gl-dim-finite-dim-regular}.
Ainsi $\Ext^{d + 1}_A(M, A) = 0$ pour tout $A$-module $M$ ; voir Algèbre, lemme \ref{algebra-lemma-projective-dimension-ext}.
Par conséquent, $A$ est de dimension injective finie comme $A$-module
d'après Compléments d'algèbre, lemme \ref{more-algebra-lemma-injective-amplitude}.
Il en résulte que $A[0]$ est un complexe dualisant, donc $A$
est de Gorenstein d'après la remarque qui suit la définition.
\end{proof}

\begin{lemma}
\label{lemma-gorenstein}
Soit $A$ un anneau noethérien.
\begin{enumerate}
\item Si $A$ possède un complexe dualisant $\omega_A^\bullet$, alors
\begin{enumerate}
\item $A$ est de Gorenstein $\Leftrightarrow$ $\omega_A^\bullet$ est un objet inversible de $D(A)$ ;
\item $A_\mathfrak p$ est de Gorenstein $\Leftrightarrow$
$(\omega_A^\bullet)_\mathfrak p$ est un objet inversible de $D(A_\mathfrak p)$ ;
\item $\{\mathfrak p \in \Spec(A) \mid A_\mathfrak p\text{ est de Gorenstein}\}$ est un ouvert.
\end{enumerate}
\item Si $A$ est de Gorenstein, alors $A$ possède un complexe dualisant
si et seulement si $A[0]$ est un complexe dualisant.
\end{enumerate}
\end{lemma}

\begin{proof}
Pour les objets inversibles de $D(A)$, voir Compléments d'algèbre,
lemme \ref{more-algebra-lemma-invertible-derived} et la discussion de la section \ref{section-dualizing}.

\medskip\noindent
D'après le lemme \ref{lemma-dualizing-localize}, pour tout $\mathfrak p$, le complexe $(\omega_A^\bullet)_\mathfrak p$
est un complexe dualisant sur $A_\mathfrak p$.
La définition et l'unicité des complexes dualisants (lemme \ref{lemma-dualizing-unique})
donnent (1)(b).

\medskip\noindent
Pour voir (1)(c), supposons $A_\mathfrak p$ de Gorenstein.
Soit $n_x$ l'unique entier tel que $H^{n_{x}}((\omega_A^\bullet)_\mathfrak p)$ soit non nul et isomorphe à $A_\mathfrak p$.
Puisque $\omega_A^\bullet$ appartient à $D^b_{\textit{Coh}}(A)$, il n'y a qu'un nombre fini
de $A$-modules de type fini non nuls $H^i(\omega_A^\bullet)$.
Il existe donc $f \in A$ tel que $f \not \in \mathfrak p$ et tel que seul $H^{n_x}((\omega_A^\bullet)_f)$
soit non nul et engendré par $1$ élément sur $A_f$.
Les complexes dualisants étant fidèles par définition, on en déduit que $A_f \cong H^{n_x}((\omega_A^\bullet)_f)$.
Ainsi $A_\mathfrak q$ est de Gorenstein pour tout $\mathfrak q \in D(f)$.
Cela prouve que l'ensemble de (1)(c) est ouvert.

\medskip\noindent
Preuve de (1)(a). L'implication $\Leftarrow$ résulte de (1)(b).
L'implication $\Rightarrow$ découle du paragraphe précédent,
où l'on a montré que, si $A_\mathfrak p$ est de Gorenstein,
il existe $f \in A$ tel que $f \not \in \mathfrak p$ et tel que le complexe $(\omega_A^\bullet)_f$
n'ait qu'un module de cohomologie non nul, qui est inversible.

\medskip\noindent
Si $A[0]$ est un complexe dualisant, alors $A$ est de Gorenstein d'après (1).
Réciproquement, (1) montre que $\omega_A^\bullet$ est localement isomorphe à un décalé de $A$.
Puisque la propriété d'être un complexe dualisant est locale
(lemme \ref{lemma-dualizing-glue}), le résultat est clair.
\end{proof}

\begin{lemma}
\label{lemma-gorenstein-ext}
Soit $(A, \mathfrak m, \kappa)$ un anneau local noethérien.
Alors $A$ est de Gorenstein si et seulement si $\Ext^i_A(\kappa, A)$ est nul pour $i \gg 0$.
\end{lemma}

\begin{proof}
Observons que $A[0]$ est un complexe dualisant pour $A$ si et seulement si
$A$ est de dimension injective finie comme $A$-module
(cela résulte immédiatement de la définition \ref{definition-dualizing}).
Le lemme découle donc de Compléments d'algèbre, lemme \ref{more-algebra-lemma-finite-injective-dimension-Noetherian-local}.
\end{proof}

\begin{lemma}
\label{lemma-characterize-gorenstein}
Soit $A$ un anneau noethérien. Les conditions suivantes sont équivalentes :
\begin{enumerate}
\item $A$ est de dimension injective finie comme module sur lui-même ;
\item $A$ est de Gorenstein et de dimension finie.
\end{enumerate}
Dans ce cas, $A[0]$ est un complexe dualisant.
\end{lemma}

\begin{proof}
Si $A$ est de dimension injective finie comme module sur $A$,
alors $A[0]$ est un complexe dualisant sur $A$
(voir la définition \ref{definition-dualizing}). On en déduit que $A$ est de dimension finie
d'après le lemme \ref{lemma-universally-catenary} et de Gorenstein d'après le lemme \ref{lemma-gorenstein}.
Supposons $A$ de Gorenstein et $\dim(A) \leq d < \infty$.
Soit $I \subset A$ un idéal quelconque. Affirmons que $\Ext^i_A(A/I, A) = 0$ pour $i > d$ ;
cela prouve que $A$ est de dimension injective finie comme module
sur lui-même d'après Compléments d'algèbre, lemme \ref{more-algebra-lemma-injective-amplitude}.
La formation des groupes $\Ext$ commute à la localisation
(voir par exemple Compléments d'algèbre, lemme \ref{more-algebra-lemma-base-change-RHom}).
Pour prouver l'annulation, on peut donc supposer $A$ local de dimension $e < d$.
Puisque $A$ est de Gorenstein, $\omega_A^\bullet = A[e]$ est un complexe dualisant normalisé.
L'annulation annoncée découle alors, par exemple, du lemme \ref{lemma-sitting-in-degrees}.
\end{proof}

\begin{lemma}
\label{lemma-gorenstein-divide-by-nonzerodivisor}
Soit $(A, \mathfrak m, \kappa)$ un anneau local noethérien. Soit $f \in \mathfrak m$ un non-diviseur de zéro.
Posons $B = A/(f)$. Alors $A$ est de Gorenstein si et seulement si $B$ l'est.
\end{lemma}

\begin{proof}
Si $A$ est de Gorenstein, alors $B$ l'est d'après le lemme \ref{lemma-divide-by-nonzerodivisor}.
Réciproquement, supposons $B$ de Gorenstein.
Alors $\Ext^i_B(\kappa, B)$ est nul pour $i \gg 0$ (lemme \ref{lemma-gorenstein-ext}).
Rappelons que $R\Hom(B, -) : D(A) \to D(B)$ est adjoint à droite de la restriction (lemme \ref{lemma-right-adjoint}).
Ainsi
$$
R\Hom_A(\kappa, A) = R\Hom_B(\kappa, R\Hom(B, A)) =
R\Hom_B(\kappa, B[1])
$$
La dernière égalité s'obtient par calcul direct ou par le lemme \ref{lemma-compute-for-effective-Cartier-algebraic}.
On voit donc que $\Ext^i_A(\kappa, A)$ est nul pour $i \gg 0$,
et $A$ est de Gorenstein (lemme \ref{lemma-gorenstein-ext}).
\end{proof}

\begin{lemma}
\label{lemma-gorenstein-lci}
Si $A \to B$ est un homomorphisme d'anneaux localement d'intersection complète
et si $A$ est un anneau noethérien de Gorenstein,
alors $B$ est un anneau de Gorenstein.
\end{lemma}

\begin{proof}
D'après Compléments d'algèbre, définition \ref{more-algebra-definition-local-complete-intersection},
on peut écrire $B = A[x_1, \ldots, x_n]/I$, où $I$ est un idéal Koszul-régulier.
Observons qu'un anneau de polynômes sur un anneau de Gorenstein $A$
est de Gorenstein : se ramener au cas où $A$ est local,
puis utiliser les lemmes \ref{lemma-dualizing-polynomial-ring} et \ref{lemma-gorenstein}.
Un idéal Koszul-régulier est, par définition, localement engendré
par une suite Koszul-régulière ; voir Compléments d'algèbre, section \ref{more-algebra-section-ideals}.
En examinant les anneaux locaux de $A[x_1, \ldots, x_n]$, on voit qu'il suffit de montrer ceci :
si $R$ est un anneau local noethérien de Gorenstein et si $f_1, \ldots, f_c \in \mathfrak m_R$
est une suite Koszul-régulière, alors $R/(f_1, \ldots, f_c)$ est de Gorenstein.
Cela résulte du lemme \ref{lemma-gorenstein-divide-by-nonzerodivisor} et du fait qu'une suite Koszul-régulière
dans $R$ est simplement une suite régulière
(Compléments d'algèbre, lemme \ref{more-algebra-lemma-noetherian-finite-all-equivalent}).
\end{proof}

\begin{lemma}
\label{lemma-flat-under-gorenstein}
Soit $A \to B$ un homomorphisme local plat d'anneaux locaux noethériens.
Les conditions suivantes sont équivalentes :
\begin{enumerate}
\item $B$ est de Gorenstein ;
\item $A$ et $B/\mathfrak m_A B$ sont de Gorenstein.
\end{enumerate}
\end{lemma}

\begin{proof}
Nous utiliserons ci-dessous sans autre mention le fait qu'un anneau local
de Gorenstein est de dimension injective finie, ainsi que le lemme \ref{lemma-gorenstein-ext}.
D'après Compléments d'algèbre, lemme \ref{more-algebra-lemma-pseudo-coherence-and-base-change-ext}, on a
$$
\Ext^i_A(\kappa_A, A) \otimes_A B =
\Ext^i_B(B/\mathfrak m_A B, B)
$$
pour tout $i$.

\medskip\noindent
Supposons (2). Puisque $R\Hom(B/\mathfrak m_A B, -) : D(B) \to D(B/\mathfrak m_A B)$ est adjoint à droite de la restriction
(lemme \ref{lemma-right-adjoint}), on obtient
$$
R\Hom_B(\kappa_B, B) =
R\Hom_{B/\mathfrak m_A B}(\kappa_B, R\Hom(B/\mathfrak m_A B, B))
$$
Les modules de cohomologie de $R\Hom(B/\mathfrak m_A B, B)$ sont les modules $\Ext^i_B(B/\mathfrak m_A B, B) =
\Ext^i_A(\kappa_A, A) \otimes_A B$.
Puisque $A$ est de Gorenstein, seul un nombre fini d'entre eux est non nul
et chacun est isomorphe à une somme directe de copies de $B/\mathfrak m_A B$.
Puisque $B/\mathfrak m_A B$ est de Gorenstein, on en déduit que $R\Hom_B(B/\mathfrak m_B, B)$
n'a qu'un nombre fini de modules de cohomologie non nuls.
Ainsi $B$ est de Gorenstein.

\medskip\noindent
Supposons (1). Puisque $B$ est de dimension injective finie,
$\Ext^i_B(B/\mathfrak m_A B, B)$ est $0$ pour $i \gg 0$.
Puisque $A \to B$ est fidèlement plat, on en déduit que $\Ext^i_A(\kappa_A, A)$ est $0$ pour $i \gg 0$.
Ainsi $A$ est de Gorenstein. Cela entraîne que $\Ext^i_A(\kappa_A, A)$ est non nul
pour exactement un $i$, à savoir $i = \dim(A)$, et que $\Ext^{\dim(A)}_A(\kappa_A, A) \cong \kappa_A$
(voir les lemmes \ref{lemma-normalized-finite}, \ref{lemma-apply-CM} et \ref{lemma-gorenstein-CM}).
On voit donc que $\Ext^i_B(B/\mathfrak m_A B, B)$ est nul sauf pour un $i$,
à savoir $i = \dim(A)$, et que $\Ext^{\dim(A)}_B(B/\mathfrak m_A B, B) \cong B/\mathfrak m_A B$.
Ainsi $B/\mathfrak m_A B$ est de Gorenstein d'après le lemme \ref{lemma-normalized-finite}.
\end{proof}

\begin{lemma}
\label{lemma-tor-injective-hull}
Soit $(A, \mathfrak m, \kappa)$ un anneau local noethérien de Gorenstein de dimension $d$.
Soit $E$ l'enveloppe injective de $\kappa$.
Alors $\text{Tor}_i^A(E, \kappa)$ est nul pour $i \not = d$ et $\text{Tor}_d^A(E, \kappa) = \kappa$.
\end{lemma}

\begin{proof}
Puisque $A$ est de Gorenstein, $\omega_A^\bullet = A[d]$ est un complexe dualisant normalisé pour $A$.
De plus, $E$ est l'unique module de cohomologie non nul de $R\Gamma_\mathfrak m(\omega_A^\bullet)$,
placé en degré $0$ ; voir le lemme \ref{lemma-local-cohomology-of-dualizing}.
D'après le lemme \ref{lemma-torsion-tensor-product}, on a
$$
E \otimes_A^\mathbf{L} \kappa =
R\Gamma_\mathfrak m(\omega_A^\bullet) \otimes_A^\mathbf{L} \kappa =
R\Gamma_\mathfrak m(\omega_A^\bullet \otimes_A^\mathbf{L} \kappa) =
R\Gamma_\mathfrak m(\kappa[d]) = \kappa[d]
$$
d'où le lemme.
\end{proof}




\section{L'abondance des complexes dualisants}
\label{section-ubiquity-dualizing}

\noindent
Beaucoup d'anneaux noethériens possèdent des complexes dualisants.

\begin{lemma}
\label{lemma-flat-unramified}
Soit $A \to B$ un homomorphisme local d'anneaux locaux noethériens.
Soit $\omega_A^\bullet$ un complexe dualisant normalisé.
Si $A \to B$ est plat et si $\mathfrak m_A B = \mathfrak m_B$,
alors $\omega_A^\bullet \otimes_A B$ est un complexe dualisant normalisé pour $B$.
\end{lemma}

\begin{proof}
Il est clair que $\omega_A^\bullet \otimes_A B$ appartient à $D^b_{\textit{Coh}}(B)$.
Soient $\kappa_A$ et $\kappa_B$ les corps résiduels de $A$ et de $B$.
D'après Compléments d'algèbre, lemme \ref{more-algebra-lemma-base-change-RHom}, on voit que
$$
R\Hom_B(\kappa_B, \omega_A^\bullet \otimes_A B) =
R\Hom_A(\kappa_A, \omega_A^\bullet) \otimes_A B =
\kappa_A[0] \otimes_A B = \kappa_B[0]
$$
Ainsi $\omega_A^\bullet \otimes_A B$ est de dimension injective finie d'après
Compléments d'algèbre, lemme \ref{more-algebra-lemma-finite-injective-dimension-Noetherian-local}.
Enfin, les mêmes arguments donnent
$$
R\Hom_B(\omega_A^\bullet \otimes_A B, \omega_A^\bullet \otimes_A B) =
R\Hom_A(\omega_A^\bullet, \omega_A^\bullet) \otimes_A B = A \otimes_A B = B
$$
comme voulu.
\end{proof}

\begin{lemma}
\label{lemma-flat-iso-mod-I}
Soit $A \to B$ un homomorphisme plat d'anneaux noethériens.
Soit $I \subset A$ un idéal tel que $A/I = B/IB$ et tel que $IB$
soit contenu dans le radical de Jacobson de $B$.
Soit $\omega_A^\bullet$ un complexe dualisant.
Alors $\omega_A^\bullet \otimes_A B$ est un complexe dualisant pour $B$.
\end{lemma}

\begin{proof}
Il est clair que $\omega_A^\bullet \otimes_A B$ appartient à $D^b_{\textit{Coh}}(B)$.
D'après Compléments d'algèbre, lemme \ref{more-algebra-lemma-base-change-RHom}, on voit que
$$
R\Hom_B(K \otimes_A B, \omega_A^\bullet \otimes_A B) =
R\Hom_A(K, \omega_A^\bullet) \otimes_A B
$$
pour tout $K \in D^b_{\textit{Coh}}(A)$. Pour tout idéal $IB \subset J \subset B$,
il existe un unique idéal $I \subset J' \subset A$ tel que $A/J' \otimes_A B = B/J$.
Ainsi $\omega_A^\bullet \otimes_A B$ est de dimension injective finie
d'après Compléments d'algèbre, lemme \ref{more-algebra-lemma-finite-injective-dimension-Noetherian-radical}.
Enfin, on a aussi
$$
R\Hom_B(\omega_A^\bullet \otimes_A B, \omega_A^\bullet \otimes_A B) =
R\Hom_A(\omega_A^\bullet, \omega_A^\bullet) \otimes_A B = A \otimes_A B = B
$$
comme voulu.
\end{proof}

\begin{lemma}
\label{lemma-completion-henselization-dualizing}
Soient $A$ un anneau noethérien et $I \subset A$ un idéal.
Soit $\omega_A^\bullet$ un complexe dualisant.
\begin{enumerate}
\item $\omega_A^\bullet \otimes_A A^h$ est un complexe dualisant sur l'hensélisé $(A^h, I^h)$ du couple $(A, I)$ ;
\item $\omega_A^\bullet \otimes_A A^\wedge$ est un complexe dualisant sur le complété $I$-adique $A^\wedge$ ;
\item si $A$ est local, alors $\omega_A^\bullet \otimes_A A^h$,
resp.\ $\omega_A^\bullet \otimes_A A^{sh}$, est un complexe dualisant
sur l'hensélisé, resp.\ l'hensélisé strict de $A$.
\end{enumerate}
\end{lemma}

\begin{proof}
Cela résulte immédiatement des lemmes \ref{lemma-flat-unramified} et \ref{lemma-flat-iso-mod-I}.
Voir Compléments d'algèbre, sections \ref{more-algebra-section-henselian-pairs}, \ref{more-algebra-section-permanence-completion} et \ref{more-algebra-section-permanence-henselization},
ainsi que Algèbre, sections \ref{algebra-section-completion} et \ref{algebra-section-completion-noetherian},
pour les complétions et les hensélisations.
\end{proof}

\begin{lemma}
\label{lemma-ubiquity-dualizing}
Les anneaux des types suivants possèdent un complexe dualisant :
\begin{enumerate}
\item les corps ;
\item les anneaux locaux noethériens complets ;
\item $\mathbf{Z}$ ;
\item les anneaux de Dedekind ;
\item tout anneau obtenu à partir d'un des anneaux précédents
en prenant une algèbre essentiellement de type fini, un complété
pour la topologie définie par un idéal, un hensélisé ou un hensélisé strict.
\end{enumerate}
\end{lemma}

\begin{proof}
L'assertion (5) résulte de la proposition \ref{proposition-dualizing-essentially-finite-type} et du lemme \ref{lemma-completion-henselization-dualizing}.
D'après le lemme \ref{lemma-regular-gorenstein}, un anneau local régulier possède un complexe dualisant.
Un anneau local noethérien complet est quotient d'un anneau local régulier
par le théorème de structure de Cohen (Algèbre, théorème \ref{algebra-theorem-cohen-structure-theorem}).
Soit $A$ un anneau de Dedekind. Tout idéal $I$ est alors un
$A$-module projectif de type fini (cela résulte d'Algèbre, lemme \ref{algebra-lemma-finite-projective},
et du fait que les anneaux locaux de $A$ sont des anneaux
de valuation discrète, donc principaux).
Ainsi tout $A$-module est de dimension injective finie au plus $1$
d'après Compléments d'algèbre, lemme \ref{more-algebra-lemma-injective-amplitude}.
Il en résulte aisément que $A[0]$ est un complexe dualisant.
\end{proof}





\section{Fibres formelles}
\label{section-formal-fibres}

\noindent
Cette section poursuit Compléments d'algèbre, section \ref{more-algebra-section-properties-formal-fibres}.
Nous y avons établi une théorie assez satisfaisante des anneaux noethériens
$A$ dont les anneaux locaux ont des fibres formelles de Cohen-Macaulay.
Plus précisément, nous avons prouvé que
(1) il suffit de vérifier que les fibres formelles des localisés
aux idéaux maximaux sont de Cohen-Macaulay,
(2) la propriété se transmet aux anneaux de type fini sur $A$,
(3) les fibres de $A \to A^\wedge$ sont de Cohen-Macaulay pour tout complété $A^\wedge$ de $A$,
et (4) la propriété se transmet aux hensélisés de $A$.
Voir Compléments d'algèbre, lemme \ref{more-algebra-lemma-check-P-ring-maximal-ideals}, proposition \ref{more-algebra-proposition-finite-type-over-P-ring},
lemme \ref{more-algebra-lemma-map-P-ring-to-completion-P} et lemme \ref{more-algebra-lemma-henselization-pair-P-ring}.
Il en va de même pour les anneaux noethériens dont les anneaux locaux
ont des fibres formelles géométriquement réduites, géométriquement normales,
$(S_n)$ ou géométriquement $(R_n)$.
Nous verrons ici que les mêmes résultats valent pour les anneaux noethériens
dont les anneaux locaux ont des fibres formelles de Gorenstein
ou localement d'intersection complète.
Cela concerne ce chapitre, car un anneau noethérien possédant
un complexe dualisant en est un exemple.

\begin{lemma}
\label{lemma-formal-fibres-gorenstein}
Les propriétés (A), (B), (C), (D) et (E) de
Compléments d'algèbre, section \ref{more-algebra-section-properties-formal-fibres},
sont satisfaites pour $P(k \to R) =$« $R$ est un anneau de Gorenstein ».
\end{lemma}

\begin{proof}
Puisque le résultat est déjà connu avec « Cohen-Macaulay » à la place
de « Gorenstein », on peut à chaque étape supposer l'anneau de Cohen-Macaulay.
Cela n'aide guère à la démonstration, mais peut être utile psychologiquement.

\medskip\noindent
Assertion (A). Soit $K/k$ une extension de corps de type fini.
Soit $R$ une $k$-algèbre de Gorenstein.
On peut trouver une intersection complète globale
$A = k[x_1, \ldots, x_n]/(f_1, \ldots, f_c)$
sur $k$ telle que $K$ soit isomorphe au corps des fractions de $A$ ;
voir Algèbre, lemme \ref{algebra-lemma-colimit-syntomic}.
Alors $R \to R \otimes_k A$ est une intersection complète globale relative.
Ainsi $R \otimes_k A$ est de Gorenstein d'après le lemme \ref{lemma-gorenstein-lci}.
Il en est de même de $R \otimes_k K$, qui en est un localisé.

\medskip\noindent
Preuve de (B). C'est clair, car un anneau est de Gorenstein
si et seulement si tous ses anneaux locaux le sont.

\medskip\noindent
Assertion (C). Soient $A \to B \to C$ des homomorphismes plats d'anneaux noethériens.
Supposons que les fibres de $A \to B$ soient de Gorenstein et que $B \to C$ soit régulier.
Il faut montrer que les fibres de $A \to C$ sont de Gorenstein.
On peut clairement supposer que $A = k$ est un corps,
puis que $B \to C$ est un homomorphisme local régulier d'anneaux locaux noethériens.
Alors $B$ est de Gorenstein et $C/\mathfrak m_B C$ est régulier,
donc en particulier de Gorenstein (lemme \ref{lemma-regular-gorenstein}).
Ainsi $C$ est de Gorenstein d'après le lemme \ref{lemma-flat-under-gorenstein}.

\medskip\noindent
Assertion (D). Cela résulte du lemme \ref{lemma-flat-under-gorenstein}.
L'assertion (E) est immédiate, car la condition ne fait pas intervenir le corps de base.
\end{proof}

\begin{lemma}
\label{lemma-dualizing-gorenstein-formal-fibres}
Soit $A$ un anneau local noethérien. Si $A$ possède un complexe dualisant,
les fibres formelles de $A$ sont de Gorenstein.
\end{lemma}

\begin{proof}
Soit $\mathfrak p$ un idéal premier de $A$.
La fibre formelle de $A$ en $\mathfrak p$ est isomorphe à la fibre formelle
de $A/\mathfrak p$ en $(0)$. Le quotient $A/\mathfrak p$ possède un complexe dualisant
(lemme \ref{lemma-dualizing-quotient}). Il suffit donc de vérifier l'énoncé lorsque $A$
est un anneau local intègre et $\mathfrak p = (0)$.
Soit $\omega_A^\bullet$ un complexe dualisant pour $A$.
Alors $\omega_A^\bullet \otimes_A A^\wedge$ est un complexe dualisant pour le complété $A^\wedge$
(lemme \ref{lemma-flat-unramified}).
De plus, $\omega_A^\bullet \otimes_A K$ est un complexe dualisant pour le corps des fractions $K$
de $A$ (lemme \ref{lemma-dualizing-localize}).
Ainsi $\omega_A^\bullet \otimes_A K$ est isomorphe à $K[n]$ pour un certain $n \in \mathbf{Z}$.
De même, on obtient pour la fibre formelle $A^\wedge \otimes_A K$ le complexe dualisant
$$
\omega_A^\bullet \otimes_A A^\wedge \otimes_{A^\wedge} (A^\wedge \otimes_A K) =
(\omega_A^\bullet \otimes_A K) \otimes_K (A^\wedge \otimes_A K) \cong
(A^\wedge \otimes_A K)[n]
$$
comme voulu.
\end{proof}

\noindent
Voici la vérification annoncée dans Algèbre à puissances divisées,
remarque \ref{dpa-remark-no-good-ci-map}.

\begin{lemma}
\label{lemma-formal-fibres-lci}
Les propriétés (A), (B), (C), (D) et (E) de
Compléments d'algèbre, section \ref{more-algebra-section-properties-formal-fibres},
sont satisfaites pour $P(k \to R) =$« $R$ est localement une intersection complète ».
Voir Algèbre à puissances divisées, définition \ref{dpa-definition-lci}.
\end{lemma}

\begin{proof}
Assertion (A). Soit $K/k$ une extension de corps de type fini.
Soit $R$ une $k$-algèbre localement d'intersection complète.
On peut trouver une intersection complète globale
$A = k[x_1, \ldots, x_n]/(f_1, \ldots, f_c)$
sur $k$ telle que $K$ soit isomorphe au corps des fractions de $A$ ;
voir Algèbre, lemme \ref{algebra-lemma-colimit-syntomic}.
Alors $R \to R \otimes_k A$ est une intersection complète globale relative.
Il en résulte que $R \otimes_k A$ est localement d'intersection complète
d'après Algèbre à puissances divisées, lemme \ref{dpa-lemma-avramov}.

\medskip\noindent
Preuve de (B). C'est clair, car un anneau est localement d'intersection
complète si et seulement si tous ses anneaux locaux sont des intersections complètes.

\medskip\noindent
Assertion (C). Soient $A \to B \to C$ des homomorphismes plats d'anneaux noethériens.
Supposons que les fibres de $A \to B$ soient localement d'intersection complète
et que $B \to C$ soit régulier. Il faut montrer que les fibres de $A \to C$
sont localement d'intersection complète.
On peut clairement supposer que $A = k$ est un corps,
puis que $B \to C$ est un homomorphisme local régulier d'anneaux locaux noethériens.
Alors $B$ est une intersection complète et $C/\mathfrak m_B C$ est régulier,
donc en particulier une intersection complète par définition.
Ainsi $C$ est une intersection complète d'après
Algèbre à puissances divisées, lemme \ref{dpa-lemma-avramov}.

\medskip\noindent
Assertion (D). Les mêmes arguments que pour (C) s'appliquent,
en utilisant l'autre implication d'Algèbre à puissances divisées,
lemme \ref{dpa-lemma-avramov}.
L'assertion (E) est immédiate, car la condition ne fait pas intervenir le corps de base.
\end{proof}






\section{Construction algébrique de l'image inverse extraordinaire}
\label{section-relative-dualizing-complex-algebraic}

\noindent
Pour un homomorphisme de type fini $R \to A$ d'anneaux noethériens,
nous construirons un foncteur $\varphi^! : D(R) \to D(A)$ bien défini à isomorphisme non unique près.
Comme nous le verrons dans Dualité pour les schémas, remarque \ref{duality-remark-local-calculation-shriek},
il coïncide à isomorphisme près avec les foncteurs image inverse
extraordinaire de la théorie de la dualité pour les schémas.
Pour motiver cette construction, signalons deux propriétés supplémentaires :
\begin{enumerate}
\item $\varphi^!$ envoie un complexe dualisant pour $R$, s'il en existe,
sur un complexe dualisant pour $A$ ;
\item $\omega_{A/R}^\bullet = \varphi^!(R)$ est une sorte de complexe dualisant relatif :
il appartient à $D^b_{\textit{Coh}}(A)$ et sa restriction aux fibres est un complexe
dualisant si $R \to A$ est plat.
\end{enumerate}
Ces énoncés sont les lemmes \ref{lemma-shriek-dualizing-algebraic} et \ref{lemma-relative-dualizing-algebraic}.

\medskip\noindent
Soit $\varphi : R \to A$ un homomorphisme de type fini d'anneaux noethériens.
Définissons un foncteur $\varphi^! : D(R) \to D(A)$ de la manière suivante :
\begin{enumerate}
\item si $\varphi : R \to A$ est surjectif, posons $\varphi^!(K) = R\Hom(A, K)$.
Nous utilisons ici le foncteur $R\Hom(A, -) : D(R) \to D(A)$ de la section \ref{section-trivial} ;
\item en général, choisissons une surjection $\psi : P \to A$ avec $P = R[x_1, \ldots, x_n]$
et posons $\varphi^!(K) = \psi^!(K \otimes_R^\mathbf{L} P)[n]$.
Nous utilisons ici le foncteur $- \otimes_R^\mathbf{L} P : D(R) \to D(P)$
de Compléments d'algèbre, section \ref{more-algebra-section-derived-base-change}.
\end{enumerate}
Noter le décalage $[n]$ par le nombre de variables de l'anneau de polynômes.
Cette construction n'est {\bf pas} canonique,
et le foncteur $\varphi^!$ ne sera bien défini qu'à isomorphisme
de foncteurs non unique près\footnote{On peut rendre la construction canonique :
utiliser $\Omega^n_{P/R}[n]$ au lieu de $P[n]$ dans la construction,
puis dans le lemme \ref{lemma-well-defined}.
Cela rend toutefois le contenu de cette section beaucoup plus compliqué.}.

\begin{lemma}
\label{lemma-well-defined}
Soit $\varphi : R \to A$ un homomorphisme de type fini d'anneaux noethériens.
Le foncteur $\varphi^!$ est bien défini à isomorphisme près.
\end{lemma}

\begin{proof}
Supposons que $\psi_1 : P_1 = R[x_1, \ldots, x_n] \to A$ et $\psi_2 : P_2 = R[y_1, \ldots, y_m] \to A$ soient deux surjections
d'anneaux de polynômes sur $A$. On obtient un diagramme commutatif
$$
\xymatrix{
R[x_1, \ldots, x_n, y_1, \ldots, y_m]
\ar[d]^{x_i \mapsto g_i} \ar[rr]_-{y_j \mapsto f_j} & &
R[x_1, \ldots, x_n] \ar[d] \\
R[y_1, \ldots, y_m] \ar[rr] & & A
}
$$
où $f_j$ et $g_i$ sont choisis de façon que $\psi_1(f_j) = \psi_2(y_j)$ et $\psi_2(g_i) = \psi_1(x_i)$.
Par symétrie, il suffit de prouver que les foncteurs définis à l'aide de
$P \to A$ et de $P[y_1, \ldots, y_m] \to A$ sont isomorphes.
Par récurrence, on peut supposer $m = 1$.
On se ramène ainsi au cas traité dans le paragraphe suivant.

\medskip\noindent
Ici $\psi : P \to A$ est donné et $\chi : P[y] \to A$ induit $\psi$ sur $P$.
Écrivons $Q = P[y]$. Choisissons $g \in P$ tel que $\psi(g) = \chi(y)$.
Notons $\pi : Q \to P$ l'homomorphisme de $P$-algèbres tel que $\pi(y) = g$.
Alors $\chi = \psi \circ \pi$, donc $\chi^! = \psi^! \circ \pi^!$, puisque les deux sont adjoints
au foncteur de restriction $D(A) \to D(Q)$ d'après la section \ref{section-trivial}.
Ainsi
$$
\chi^!\left(K \otimes_R^\mathbf{L} Q\right)[n + 1] =
\psi^!\left(\pi^!\left(K \otimes_R^\mathbf{L} Q\right)[1]\right)[n]
$$
Il suffit donc de montrer que $\pi^!(K \otimes_R^\mathbf{L} Q[1]) = K \otimes_R^\mathbf{L} P$.
Il suffit par conséquent de montrer que le foncteur $\pi^!(-) : D(Q) \to D(P)$
est isomorphe à $K \mapsto K \otimes_Q^\mathbf{L} P[-1]$, ce qui résulte du lemme \ref{lemma-compute-for-effective-Cartier-algebraic}.
\end{proof}

\begin{lemma}
\label{lemma-shriek-boundedness}
Soit $\varphi : R \to A$ un homomorphisme de type fini d'anneaux noethériens.
\begin{enumerate}
\item $\varphi^!$ envoie $D^+(R)$ dans $D^+(A)$ et $D^+_{\textit{Coh}}(R)$ dans $D^+_{\textit{Coh}}(A)$.
\item Si $\varphi$ est parfait, alors $\varphi^!$ envoie
$D^-(R)$ dans $D^-(A)$, $D^-_{\textit{Coh}}(R)$ dans $D^-_{\textit{Coh}}(A)$ et $D^b_{\textit{Coh}}(R)$ dans $D^b_{\textit{Coh}}(A)$.
\end{enumerate}
\end{lemma}

\begin{proof}
Choisissons une factorisation $R \to P \to A$ comme dans la définition de $\varphi^!$.
Le foncteur $- \otimes_R^\mathbf{L} : D(R) \to D(P)$ préserve les sous-catégories $D^+, D^+_{\textit{Coh}}, D^-, D^-_{\textit{Coh}}, D^b_{\textit{Coh}}$.
Le foncteur $R\Hom(A, -) : D(P) \to D(A)$ préserve $D^+$ et $D^+_{\textit{Coh}}$ d'après le lemme \ref{lemma-exact-support-coherent}.
Si $R \to A$ est parfait, alors $A$ est parfait comme $P$-module ;
voir Compléments d'algèbre, lemme \ref{more-algebra-lemma-perfect-ring-map}.
Rappelons que la restriction de $R\Hom(A, K)$ à $D(P)$ est $R\Hom_P(A, K)$.
D'après Compléments d'algèbre, lemme \ref{more-algebra-lemma-dual-perfect-complex}, on a $R\Hom_P(A, K) = E \otimes_P^\mathbf{L} K$
pour un certain objet parfait $E \in D(P)$.
Puisque $E$ peut être représenté par un complexe fini de $P$-modules
projectifs de type fini, il est clair que $R\Hom_P(A, K)$ appartient à $D^-(P), D^-_{\textit{Coh}}(P), D^b_{\textit{Coh}}(P)$
dès que $K$ y appartient.
Puisque le foncteur de restriction $D(A) \to D(P)$ reflète l'appartenance
à ces sous-catégories, la démonstration est achevée.
\end{proof}

\begin{lemma}
\label{lemma-shriek-dualizing-algebraic}
Soit $\varphi$ un homomorphisme de type fini d'anneaux noethériens.
Si $\omega_R^\bullet$ est un complexe dualisant pour $R$,
alors $\varphi^!(\omega_R^\bullet)$ est un complexe dualisant pour $A$.
\end{lemma}

\begin{proof}
Cela résulte des lemmes \ref{lemma-dualizing-polynomial-ring} et \ref{lemma-dualizing-quotient}.
\end{proof}

\begin{lemma}
\label{lemma-flat-bc}
Soit $R \to R'$ un homomorphisme plat d'anneaux noethériens.
Soit $\varphi : R \to A$ un homomorphisme de type fini.
Soit $\varphi' : R' \to A' = A \otimes_R R'$ l'homomorphisme induit par $\varphi$.
On a alors des morphismes fonctoriels
$$
\varphi^!(K) \otimes_A^\mathbf{L} A' \longrightarrow
(\varphi')^!(K \otimes_R^\mathbf{L} R')
$$
pour $K$ dans $D(R)$, qui sont des isomorphismes pour $K \in D^+(R)$.
\end{lemma}

\begin{proof}
Choisissons une factorisation $R \to P \to A$ où $P$ est un anneau de polynômes sur $R$.
Elle donne par changement de base une factorisation correspondante $R' \to P' \to A'$.
Puisque l'on a $(K \otimes_R^\mathbf{L} P) \otimes_P^\mathbf{L} P' =
(K \otimes_R^\mathbf{L} R') \otimes_{R'}^\mathbf{L} P'$ d'après Compléments d'algèbre, lemme \ref{more-algebra-lemma-double-base-change},
il suffit de construire des morphismes
$$
R\Hom(A, K \otimes_R^\mathbf{L} P[n]) \otimes_A^\mathbf{L} A'
\longrightarrow
R\Hom(A', (K \otimes_R^\mathbf{L} P[n]) \otimes_P^\mathbf{L} P')
$$
fonctoriels en $K$. Pour cela, utilisons le morphisme (\ref{equation-base-change})
construit dans la section \ref{section-base-change-trivial-duality} pour $P, A, P', A'$.
Ce morphisme est un isomorphisme pour $K \in D^+(R)$ d'après le lemme \ref{lemma-flat-bc-surjection}.
\end{proof}

\begin{lemma}
\label{lemma-bc}
Soit $R \to R'$ un homomorphisme d'anneaux noethériens.
Soit $\varphi : R \to A$ un homomorphisme parfait
(Compléments d'algèbre, définition \ref{more-algebra-definition-pseudo-coherent-perfect})
tel que $R'$ et $A$ soient Tor-indépendants sur $R$.
Soit $\varphi' : R' \to A' = A \otimes_R R'$ l'homomorphisme induit par $\varphi$.
On a alors un isomorphisme fonctoriel
$$
\varphi^!(K) \otimes_A^\mathbf{L} A' =
(\varphi')^!(K \otimes_R^\mathbf{L} R')
$$
pour $K$ dans $D(R)$.
\end{lemma}

\begin{proof}
On peut choisir une factorisation $R \to P \to A$ où $P$ est un anneau de polynômes
sur $R$ tel que $A$ soit un $P$-module parfait ;
voir Compléments d'algèbre, lemme \ref{more-algebra-lemma-perfect-ring-map}.
Elle donne par changement de base une factorisation correspondante $R' \to P' \to A'$.
Puisque l'on a $(K \otimes_R^\mathbf{L} P) \otimes_P^\mathbf{L} P' =
(K \otimes_R^\mathbf{L} R') \otimes_{R'}^\mathbf{L} P'$ d'après Compléments d'algèbre, lemme \ref{more-algebra-lemma-double-base-change},
il suffit de construire des morphismes
$$
R\Hom(A, K \otimes_R^\mathbf{L} P[n]) \otimes_A^\mathbf{L} A'
\longrightarrow
R\Hom(A', (K \otimes_R^\mathbf{L} P[n]) \otimes_P^\mathbf{L} P')
$$
fonctoriels en $K$. On a
$$
A \otimes_P^\mathbf{L} P' = A \otimes_R^\mathbf{L} R' = A'
$$
La première égalité résulte de Compléments d'algèbre, lemme \ref{more-algebra-lemma-base-change-comparison},
appliqué à $R, R', P, P'$. La seconde vient de ce que $A$ et $R'$
sont Tor-indépendants sur $R$.
Ainsi $A$ et $P'$ sont Tor-indépendants sur $P$,
et l'on peut utiliser le morphisme (\ref{equation-base-change}) construit dans la section \ref{section-base-change-trivial-duality}
pour $P, A, P', A'$ afin d'obtenir la flèche cherchée.
D'après le lemme \ref{lemma-bc-surjection}, il suffit pour conclure de prouver que $A$
est un $P$-module parfait, ce qui a été vu ci-dessus.
\end{proof}

\begin{lemma}
\label{lemma-bc-flat}
Soit $R \to R'$ un homomorphisme d'anneaux noethériens.
Soit $\varphi : R \to A$ plat et de type fini.
Soit $\varphi' : R' \to A' = A \otimes_R R'$ l'homomorphisme induit par $\varphi$.
On a alors un isomorphisme fonctoriel
$$
\varphi^!(K) \otimes_A^\mathbf{L} A' =
(\varphi')^!(K \otimes_R^\mathbf{L} R')
$$
pour $K$ dans $D(R)$.
\end{lemma}

\begin{proof}
C'est un cas particulier du lemme \ref{lemma-bc} d'après
Compléments d'algèbre, lemme \ref{more-algebra-lemma-flat-finite-presentation-perfect}.
\end{proof}

\begin{lemma}
\label{lemma-composition-shriek-algebraic}
Soient $A \xrightarrow{a} B \xrightarrow{b} C$ des homomorphismes de type fini d'anneaux noethériens.
Il existe alors une transformation de foncteurs $b^! \circ a^! \to (b \circ a)^!$
qui est un isomorphisme sur $D^+(A)$.
\end{lemma}

\begin{proof}
Choisissons un anneau de polynômes $P = A[x_1, \ldots, x_n]$ sur $A$
et une surjection $P \to B$. Choisissons des éléments $c_1, \ldots, c_m \in C$
engendrant $C$ sur $B$. Posons $Q = P[y_1, \ldots, y_m]$ et notons $Q' = Q \otimes_P B = B[y_1, \ldots, y_m]$.
Soit $\chi : Q' \to C$ la surjection envoyant $y_j$ sur $c_j$.
On a le diagramme
$$
\xymatrix{
& Q \ar[r]_{\psi'} & Q' \ar[r]_\chi & C \\
A \ar[r] & P \ar[r]^\psi \ar[u] & B \ar[u]
}
$$
D'après le lemme \ref{lemma-flat-bc-surjection}, pour $M \in D(P)$, on a une flèche $\psi^!(M) \otimes_B^\mathbf{L} Q' \to (\psi')^!(M \otimes_P^\mathbf{L} Q)$
qui est un isomorphisme dès que $M$ est borné inférieurement.
De plus, $\chi^! \circ (\psi')^! = (\chi \circ \psi')^!$ puisque les deux foncteurs sont adjoints
au foncteur de restriction $D(C) \to D(Q)$ d'après la section \ref{section-trivial}.
On obtient alors
\begin{align*}
b^!(a^!(K))
& =
\chi^!(\psi^!(K \otimes_A^\mathbf{L} P)[n] \otimes_B^\mathbf{L} Q)[m] \\
& \to
\chi^!((\psi')^!(K \otimes_A^\mathbf{L} P \otimes_P^\mathbf{L} Q))[n + m] \\
& =
(\chi \circ \psi')^!(K\otimes_A^\mathbf{L} Q)[n + m] \\
& =
(b \circ a)^!(K)
\end{align*}
en utilisant, outre ce qui précède, Compléments d'algèbre, lemme \ref{more-algebra-lemma-double-base-change}.
\end{proof}

\begin{lemma}
\label{lemma-upper-shriek-finite}
Soit $\varphi : R \to A$ un homomorphisme fini d'anneaux noethériens.
Alors $\varphi^!$ est isomorphe au foncteur $R\Hom(A, -) : D(R) \to D(A)$ de la section \ref{section-trivial}.
\end{lemma}

\begin{proof}
Supposons $A$ engendré par $n > 1$ éléments sur $R$.
On peut factoriser $R \to A$ en deux homomorphismes finis d'anneaux,
avec à chaque étape un nombre de générateurs $< n$.
Les lemmes \ref{lemma-composition-shriek-algebraic} et \ref{lemma-composition-right-adjoints} montrent donc qu'il suffit de prouver
le lemme lorsque $A$ est engendré par un seul élément sur $R$.
Puisque $A$ est fini sur $R$, $A$ est quotient de $B = R[x]/(f)$,
où $f$ est un polynôme unitaire en $x$ (Algèbre, lemme \ref{algebra-lemma-finite-is-integral}).
En utilisant encore les lemmes de composition et l'accord des foncteurs
pour les surjections par définition, on se ramène à $R \to B = R[x]/(f)$.
Dans ce cas, le foncteur $\varphi^!$ est isomorphe à $K \mapsto K \otimes_R^\mathbf{L} B$ :
cela se prouve en utilisant le lemme \ref{lemma-compute-for-effective-Cartier-algebraic} pour l'homomorphisme $R[x] \to B$
(noter que les décalages dans la définition de $\varphi^!$ et dans le lemme ont somme nulle).
Pour le foncteur $R\Hom(B, -) : D(R) \to D(B)$, le lemme \ref{lemma-RHom-is-tensor-special} permet de se ramener
à montrer que $\Hom_R(B, R) \cong B$ comme $B$-modules.
Supposons $f$ de degré $d$.
Alors une $R$-base de $B$ est donnée par $1, x, \ldots, x^{d - 1}$.
Soit $\delta_i : B \to R$, pour $i = 0, \ldots, d - 1$, l'application $R$-linéaire
qui extrait le coefficient de $x^i$ dans cette base.
Alors $\delta_0, \ldots, \delta_{d - 1}$ est une base de $\Hom_R(B, R)$.
Enfin, pour $0 \leq i \leq d - 1$, un calcul donne
$$
x^i \delta_{d - 1} =
\delta_{d - 1 - i} + b_1 \delta_{d - i} + \ldots + b_i \delta_{d - 1}
$$
pour certains $c_1, \ldots, c_d \in R$\footnote{Si $f = x^d + a_1 x^{d - 1} + \ldots + a_d$, alors
$c_1 = -a_1$, $c_2 = a_1^2 - a_2$, $c_3 = -a_1^3 + 2a_1a_2 -a_3$, etc.}.
Ainsi $\Hom_R(B, R)$ est un $B$-module monogène de générateur $\delta_{d - 1}$.
En examinant les rangs, on en déduit que c'est un module libre
de rang $1$ sur $B$.
\end{proof}

\begin{lemma}
\label{lemma-upper-shriek-localize}
Soit $R$ un anneau noethérien et soit $f \in R$.
Si $\varphi$ désigne l'homomorphisme $R \to R_f$, alors $\varphi^!$ est isomorphe à $- \otimes_R^\mathbf{L} R_f$.
Plus généralement, si $\varphi : R \to R'$ est un homomorphisme tel que
$\Spec(R') \to \Spec(R)$ soit une immersion ouverte, alors $\varphi^!$ est isomorphe à $- \otimes_R^\mathbf{L} R'$.
\end{lemma}

\begin{proof}
Choisissons la présentation $R \to R[x] \to R[x]/(fx - 1) = R_f$ et observons que $fx - 1$
est un non-diviseur de zéro dans $R[x]$.
On peut donc appliquer le lemme \ref{lemma-compute-for-effective-Cartier-algebraic} pour calculer le foncteur $\varphi^!$.
Nous omettons les détails ; noter que les décalages dans la définition
de $\varphi^!$ et dans le lemme ont somme nulle.

\medskip\noindent
Dans le cas général, noter que $R' \otimes_R R' = R'$.
Le résultat découle donc des résultats de changement de base ci-dessus :
on peut utiliser soit le lemme \ref{lemma-flat-bc}, soit le lemme \ref{lemma-bc}.
\end{proof}

\begin{lemma}
\label{lemma-upper-shriek-is-tensor-functor}
Soit $\varphi : R \to A$ un homomorphisme parfait d'anneaux noethériens
(par exemple, $\varphi$ est plat de type fini).
Alors $\varphi^!(K) = K \otimes_R^\mathbf{L} \varphi^!(R)$ pour $K \in D(R)$.
\end{lemma}

\begin{proof}
(L'assertion entre parenthèses résulte de Compléments d'algèbre, lemme \ref{more-algebra-lemma-flat-finite-presentation-perfect}.)
On peut choisir une factorisation $R \to P \to A$ où $P$ est un anneau de polynômes
en $n$ variables sur $R$ ; alors $A$ est un $P$-module parfait,
voir Compléments d'algèbre, lemme \ref{more-algebra-lemma-perfect-ring-map}.
Rappelons que $\varphi^!(K) = R\Hom(A, K \otimes_R^\mathbf{L} P[n])$.
Le résultat découle donc du lemme \ref{lemma-RHom-is-tensor-special}
et de Compléments d'algèbre, lemme \ref{more-algebra-lemma-double-base-change}.
\end{proof}

\begin{lemma}
\label{lemma-relative-dualizing-if-have-omega}
Soit $\varphi : A \to B$ un homomorphisme de type fini d'anneaux noethériens.
Soit $\omega_A^\bullet$ un complexe dualisant pour $A$. Posons $\omega_B^\bullet = \varphi^!(\omega_A^\bullet)$.
Notons $D_A(K) = R\Hom_A(K, \omega_A^\bullet)$ pour $K \in D_{\textit{Coh}}(A)$ et $D_B(L) = R\Hom_B(L, \omega_B^\bullet)$ pour $L \in D_{\textit{Coh}}(B)$.
Il existe alors un isomorphisme fonctoriel
$$
\varphi^!(K) = D_B(D_A(K) \otimes_A^\mathbf{L} B)
$$
pour $K \in D_{\textit{Coh}}(A)$.
\end{lemma}

\begin{proof}
Observons que $\omega_B^\bullet$ est un complexe dualisant pour $B$ d'après le lemme \ref{lemma-shriek-dualizing-algebraic}.
Soient $A \to B \to C$ des homomorphismes de type fini d'anneaux noethériens.
Si le lemme vaut pour $A \to B$ et $B \to C$, il vaut pour $A \to C$.
Cela résulte du lemme \ref{lemma-composition-shriek-algebraic} et du fait que $D_B \circ D_B \cong \text{id}$ d'après le lemme \ref{lemma-dualizing}.
Il suffit donc de prouver le lemme lorsque $A \to B$ est une surjection
et lorsque $B$ est un anneau de polynômes sur $A$.

\medskip\noindent
Supposons $B = A[x_1, \ldots, x_n]$. Puisque $D_A \circ D_A \cong \text{id}$, il suffit de prouver
$D_B(K \otimes_A B) \cong D_A(K) \otimes_A B[n]$ pour $K$ dans $D_{\textit{Coh}}(A)$.
Choisissons un complexe borné $I^\bullet$ d'injectifs représentant $\omega_A^\bullet$.
Choisissons un quasi-isomorphisme $I^\bullet \otimes_A B \to J^\bullet$ où $J^\bullet$
est un complexe borné de $B$-modules.
Étant donné un complexe $K^\bullet$ de $A$-modules,
considérons le morphisme évident de complexes
$$
\Hom^\bullet(K^\bullet, I^\bullet) \otimes_A B[n]
\longrightarrow
\Hom^\bullet(K^\bullet \otimes_A B, J^\bullet[n])
$$
Le membre de gauche représente $D_A(K) \otimes_A B[n]$ et celui de droite représente $D_B(K \otimes_A B)$.
Il suffit donc de prouver que ce morphisme est un quasi-isomorphisme
si les modules de cohomologie de $K^\bullet$ sont des $A$-modules de type fini.
Observons que la cohomologie du complexe en degré $r$,
de chaque côté, ne dépend que d'un nombre fini des $K^i$.
On peut donc remplacer $K^\bullet$ par une troncature, c'est-à-dire supposer
que $K^\bullet$ représente un objet de $D^-_{\textit{Coh}}(A)$.
Alors $K^\bullet$ est quasi-isomorphe à un complexe borné supérieurement
de $A$-modules libres de type fini.
On peut donc supposer que $K^\bullet$ est un complexe borné supérieurement
de $A$-modules libres de type fini.
Il est alors facile de voir que le morphisme affiché est
un isomorphisme de complexes, ce qui conclut dans ce cas.

\medskip\noindent
Supposons $A \to B$ surjectif. Notons $i_* : D(B) \to D(A)$ le foncteur de restriction
et rappelons que $\varphi^!(-) = R\Hom(A, -)$ est adjoint à droite de $i_*$ (lemme \ref{lemma-right-adjoint}).
Pour $F \in D(B)$, on a
\begin{align*}
\Hom_B(F, D_B(D_A(K) \otimes_A^\mathbf{L} B))
& =
\Hom_B((D_A(K) \otimes_A^\mathbf{L} B) \otimes_B^\mathbf{L} F,
\omega_B^\bullet) \\
& =
\Hom_A(D_A(K) \otimes_A^\mathbf{L} i_*F, \omega_A^\bullet) \\
& =
\Hom_A(i_*F, D_A(D_A(K))) \\
& =
\Hom_A(i_*F, K) \\
& =
\Hom_B(F, \varphi^!(K))
\end{align*}
La première égalité résulte de Compléments d'algèbre, lemme \ref{more-algebra-lemma-internal-hom},
et de la définition de $D_B$.
La deuxième résulte de l'adjonction mentionnée ci-dessus et de l'égalité
$i_*((D_A(K) \otimes_A^\mathbf{L} B) \otimes_B^\mathbf{L} F) =
D_A(K) \otimes_A^\mathbf{L} i_*F$
(Compléments d'algèbre, lemme \ref{more-algebra-lemma-derived-base-change}).
La troisième résulte de Compléments d'algèbre, lemme \ref{more-algebra-lemma-internal-hom}.
La quatrième vient de ce que $D_A \circ D_A = \text{id}$, et la dernière à nouveau de l'adjonction.
Le résultat découle donc du lemme de Yoneda.
\end{proof}






\section{Complexes dualisants relatifs dans le cas noethérien}
\label{section-relative-dualizing-complexes-Noetherian}

\noindent
Soit $\varphi : R \to A$ un homomorphisme de type fini d'anneaux noethériens.
On définit le {\it complexe dualisant relatif de $A$ sur $R$}
comme l'objet
$$
\omega_{A/R}^\bullet = \varphi^!(R)
$$
de $D(A)$. Ici $\varphi^!$ est celui de la section \ref{section-relative-dualizing-complex-algebraic}.
Les résultats de cette section montrent que $\omega_{A/R}^\bullet$
est bien défini à isomorphisme non unique près.

\begin{lemma}
\label{lemma-base-change-relative-algebraic}
Soit $R \to R'$ un homomorphisme d'anneaux noethériens.
Soit $R \to A$ de type fini. Posons $A' = A \otimes_R R'$. Si
\begin{enumerate}
\item $R \to R'$ est plat, ou
\item $R \to A$ est plat, ou
\item $R \to A$ est parfait et $R'$ et $A$ sont Tor-indépendants sur $R$,
\end{enumerate}
alors il existe un isomorphisme
$\omega_{A/R}^\bullet \otimes_A^\mathbf{L} A' \to \omega^\bullet_{A'/R'}$
dans $D(A')$.
\end{lemma}

\begin{proof}
Cela résulte des lemmes \ref{lemma-flat-bc}, \ref{lemma-bc-flat} et \ref{lemma-bc} et des définitions.
\end{proof}

\begin{lemma}
\label{lemma-relative-dualizing-algebraic}
Soit $\varphi : R \to A$ un homomorphisme de type fini d'anneaux noethériens.
\begin{enumerate}
\item Si $\varphi$ est parfait, alors
\begin{enumerate}
\item $\omega_{A/R}^\bullet$ appartient à $D^b_{\textit{Coh}}(A)$ ;
\item $\omega_{A/R}^\bullet$ est de Tor-dimension finie sur $R$ ;
\item $A \to R\Hom_A(\omega_{A/R}^\bullet, \omega_{A/R}^\bullet)$ est un isomorphisme.
\end{enumerate}
\item Si $\varphi$ est plat, alors (1)(a), (1)(b) et (1)(c) sont satisfaites, et
\begin{enumerate}
\item $\omega_{A/R}^\bullet$ est $R$-parfait
(Compléments d'algèbre, définition \ref{more-algebra-definition-relatively-perfect}) ;
\item pour tout homomorphisme $R \to k$ vers un corps, le changement de base
$\omega_{A/R}^\bullet \otimes_A^\mathbf{L} (A \otimes_R k)$
est un complexe dualisant pour $A \otimes_R k$.
\end{enumerate}
\end{enumerate}
\end{lemma}

\begin{proof}
Supposons $R \to A$ parfait. Choisissons $R \to P \to A$ comme dans la définition de $\varphi^!$.
Alors $A$ est parfait comme $P$-module
(Compléments d'algèbre, lemme \ref{more-algebra-lemma-perfect-ring-map}).
Cela montre que $\omega_{A/R}^\bullet$ appartient à $D^b_{\textit{Coh}}(A)$ d'après le lemme \ref{lemma-shriek-boundedness}, d'où (1)(a).
Pour montrer que $\omega_{A/R}^\bullet$ est de Tor-dimension finie comme complexe de $R$-modules,
observons que $\omega_{A/R}^\bullet = \varphi^!(R) = R\Hom(A, P)[n]$ a pour image $R\Hom_P(A, P)[n]$ dans $D(P)$.
Cet objet est parfait dans $D(P)$ (Compléments d'algèbre, lemme \ref{more-algebra-lemma-dual-perfect-complex}),
donc de Tor-dimension finie dans $D(R)$ puisque $R \to P$ est plat.
Cela prouve (1)(b). L'objet $R\Hom_A(\omega_{A/R}^\bullet, \omega_{A/R}^\bullet)$ de $D(A)$ a pour image dans $D(P)$
\begin{align*}
R\Hom_P(\omega_{A/R}^\bullet, R\Hom(A, P)[n])
& =
R\Hom_P(R\Hom_P(A, P)[n], P)[n] \\
& =
R\Hom_P(R\Hom_P(A, P), P)
\end{align*}
Cet objet est égal à $A$ d'après Compléments d'algèbre, lemme \ref{more-algebra-lemma-dual-perfect-complex},
déjà utilisé. Cela prouve (1)(c).

\medskip\noindent
Supposons $\varphi$ plat. Alors $R \to A$ est un homomorphisme parfait
(Compléments d'algèbre, lemme \ref{more-algebra-lemma-flat-finite-presentation-perfect}),
et (1)(a), (1)(b) et (1)(c) sont satisfaites.
Bien entendu, $\omega_{A/R}^\bullet$ est alors $R$-parfait d'après (1)(a), (1)(b) et les définitions.
Soit $R \to k$ comme dans (2)(b).
D'après le lemme \ref{lemma-base-change-relative-algebraic}, il existe un isomorphisme
$$
\omega_{A/R}^\bullet \otimes_A^\mathbf{L} (A \otimes_R k) \cong
\omega^\bullet_{A \otimes_R k/k}
$$
et le membre de droite est un complexe dualisant d'après le lemme \ref{lemma-shriek-dualizing-algebraic}.
Cela achève la démonstration.
\end{proof}

\begin{lemma}
\label{lemma-base-change-dualizing-over-field}
Soit $K/k$ une extension de corps. Soit $A$ une $k$-algèbre de type fini.
Soit $A_K = A \otimes_k K$. Si $\omega_A^\bullet$ est un complexe dualisant pour $A$,
alors $\omega_A^\bullet \otimes_A A_K$ est un complexe dualisant pour $A_K$.
\end{lemma}

\begin{proof}
Par unicité des complexes dualisants, peu importe celui que l'on choisit
pour $A$ ; nous omettons la preuve détaillée.
Notons $\varphi : k \to A$ la structure d'algèbre.
On peut prendre $\omega_A^\bullet = \varphi^!(k[0])$ d'après le lemme \ref{lemma-shriek-dualizing-algebraic}.
On conclut par le lemme \ref{lemma-relative-dualizing-algebraic}.
\end{proof}

\begin{lemma}
\label{lemma-lci-shriek}
Soit $\varphi : R \to A$ un homomorphisme d'anneaux noethériens localement d'intersection complète.
Alors $\omega_{A/R}^\bullet$ est un objet inversible de $D(A)$ et $\varphi^!(K) = K \otimes_R^\mathbf{L} \omega_{A/R}^\bullet$ pour tout $K \in D(R)$.
\end{lemma}

\begin{proof}
Rappelons qu'un homomorphisme localement d'intersection complète
est parfait d'après Compléments d'algèbre, lemme \ref{more-algebra-lemma-lci-perfect}.
La dernière assertion résulte donc du lemme \ref{lemma-upper-shriek-is-tensor-functor}.
D'après Compléments d'algèbre, définition \ref{more-algebra-definition-local-complete-intersection},
on peut écrire $A = R[x_1, \ldots, x_n]/I$, où $I$ est un idéal Koszul-régulier.
La construction de $\varphi^!$ dans la section \ref{section-relative-dualizing-complex-algebraic}
montre qu'il suffit de prouver le lemme lorsque $A = R/I$,
où $I \subset R$ est un idéal Koszul-régulier.
Vérifier que $\omega_{A/R}^\bullet$ est inversible dans $D(A)$ est local sur $\Spec(A)$
d'après Compléments d'algèbre, lemme \ref{more-algebra-lemma-invertible-derived}.
De plus, la formation de $\omega_{A/R}^\bullet$ commute à la localisation sur $R$
d'après le lemme \ref{lemma-flat-bc}.
En combinant Compléments d'algèbre, définition \ref{more-algebra-definition-regular-ideal} et lemme \ref{more-algebra-lemma-noetherian-finite-all-equivalent},
avec Algèbre, lemme \ref{algebra-lemma-regular-sequence-in-neighbourhood}, on trouve $g_1, \ldots, g_r \in R$
engendrant l'idéal unité dans $A$ tels que $I_{g_j} \subset R_{g_j}$
soit engendré par une suite régulière.
On peut donc supposer $A = R/(f_1, \ldots, f_c)$, où $f_1, \ldots, f_c$ est une suite régulière dans $R$.
Considérons alors les homomorphismes d'anneaux
$$
R \to R/(f_1) \to R/(f_1, f_2) \to \ldots \to R/(f_1, \ldots, f_c) = A
$$
et utilisons le lemme \ref{lemma-composition-shriek-algebraic} (ainsi que la dernière assertion déjà prouvée)
pour nous ramener à chaque étape.
Enfin, lorsque $A = R/(f)$ pour un non-diviseur de zéro $f$,
le lemme est vrai puisque $\varphi^!(R) = R\Hom(A, R)$ est inversible d'après le lemme \ref{lemma-compute-for-effective-Cartier-algebraic}.
\end{proof}

\begin{lemma}
\label{lemma-gorenstein-shriek}
Soit $\varphi : R \to A$ un homomorphisme plat de type fini d'anneaux noethériens.
Les conditions suivantes sont équivalentes :
\begin{enumerate}
\item les fibres $A \otimes_R \kappa(\mathfrak p)$ sont de Gorenstein pour tous les idéaux premiers $\mathfrak p \subset R$ ;
\item $\omega_{A/R}^\bullet$ est un objet inversible de $D(A)$,
voir Compléments d'algèbre, lemme \ref{more-algebra-lemma-invertible-derived}.
\end{enumerate}
\end{lemma}

\begin{proof}
Si (2) est satisfaite, les anneaux des fibres $A \otimes_R \kappa(\mathfrak p)$
ont des complexes dualisants inversibles, donc sont de Gorenstein.
Voir les lemmes \ref{lemma-relative-dualizing-algebraic} et \ref{lemma-gorenstein}.

\medskip\noindent
Réciproquement, supposons (1).
Observons que $\omega_{A/R}^\bullet$ appartient à $D^b_{\textit{Coh}}(A)$ d'après le lemme \ref{lemma-shriek-boundedness}
(puisque les homomorphismes plats de type fini d'anneaux noethériens
sont parfaits, voir Compléments d'algèbre, lemme \ref{more-algebra-lemma-flat-finite-presentation-perfect}).
Prenons un idéal premier $\mathfrak q \subset A$ au-dessus de $\mathfrak p \subset R$. Alors
$$
\omega_{A/R}^\bullet \otimes_A^\mathbf{L} \kappa(\mathfrak q) =
\omega_{A/R}^\bullet \otimes_A^\mathbf{L}
(A \otimes_R \kappa(\mathfrak p))
\otimes_{(A \otimes_R \kappa(\mathfrak p))}^\mathbf{L}
\kappa(\mathfrak q)
$$
Les lemmes \ref{lemma-relative-dualizing-algebraic} et \ref{lemma-gorenstein} et l'hypothèse (1) montrent que ce complexe
a $1$ groupe de cohomologie non nul,
qui est un espace vectoriel de dimension $1$ sur $\kappa(\mathfrak q)$.
D'après Compléments d'algèbre, lemme \ref{more-algebra-lemma-lift-bounded-pseudo-coherent-to-perfect},
on en déduit que $(\omega_{A/R}^\bullet)_f$ est un objet inversible de $D(A_f)$
pour un certain $f \in A$ tel que $f \not \in \mathfrak q$.
Cela prouve (2).
\end{proof}

\noindent
Le lemme suivant permet de comprendre comment les fonctions de dimension
se transforment lorsqu'on passe à une algèbre de type fini sur un anneau noethérien.

\begin{lemma}
\label{lemma-shriek-normalized}
Soit $\varphi : R \to A$ un homomorphisme de type fini d'anneaux noethériens.
Supposons $R$ local et soit $\mathfrak m \subset A$ un idéal maximal au-dessus
de l'idéal maximal de $R$.
Si $\omega_R^\bullet$ est un complexe dualisant normalisé pour $R$,
alors $\varphi^!(\omega_R^\bullet)_\mathfrak m$ est un complexe dualisant normalisé pour $A_\mathfrak m$.
\end{lemma}

\begin{proof}
Nous savons déjà que $\varphi^!(\omega_R^\bullet)$ est un complexe dualisant pour $A$ ; voir le lemme \ref{lemma-shriek-dualizing-algebraic}.
Choisissons une factorisation $R \to P \to A$ avec $P = R[x_1, \ldots, x_n]$ comme dans la construction de $\varphi^!$.
Si l'on prouve le lemme pour $R \to P$ et l'idéal maximal $\mathfrak m'$ de $P$
correspondant à $\mathfrak m$, on obtient le résultat pour $R \to A$
en appliquant le lemme \ref{lemma-normalized-quotient} à $P_{\mathfrak m'} \to A_\mathfrak m$,
ou le lemme \ref{lemma-quotient-function} à $P \to A$.
Dans le cas $A = R[x_1, \ldots, x_n]$, on a $\dim(A_\mathfrak m) = \dim(R) + n$, par exemple d'après Algèbre, lemme \ref{algebra-lemma-dimension-base-fibre-equals-total}
(combiné avec Algèbre, lemme \ref{algebra-lemma-dim-affine-space} pour calculer la dimension de la fibre).
La normalisation de $\omega_R^\bullet$ signifie que $i = -\dim(R)$ est le plus petit indice
tel que $H^i(\omega_R^\bullet)$ soit non nul (cela résulte des lemmes \ref{lemma-sitting-in-degrees} et \ref{lemma-nonvanishing-generically-local}).
Alors $\varphi^!(\omega_R^\bullet)_\mathfrak m =
\omega_R^\bullet \otimes_R A_\mathfrak m[n]$ a son premier module de cohomologie non nul en degré $-\dim(R) - n$
et est donc le complexe dualisant normalisé pour $A_\mathfrak m$.
\end{proof}

\begin{lemma}
\label{lemma-relative-dualizing-trivial-vanishing}
Soit $R \to A$ un homomorphisme de type fini d'anneaux noethériens.
Soit $\mathfrak q \subset A$ un idéal premier au-dessus de $\mathfrak p \subset R$. Alors
$$
H^i(\omega_{A/R}^\bullet)_\mathfrak q \not = 0
\Rightarrow - d \leq i
$$
où $d$ est la dimension de la fibre de $\Spec(A) \to \Spec(R)$ au-dessus de $\mathfrak p$ au point $\mathfrak q$.
\end{lemma}

\begin{proof}
Choisissons une factorisation $R \to P \to A$ avec $P = R[x_1, \ldots, x_n]$ comme dans la section \ref{section-relative-dualizing-complex-algebraic},
de sorte que $\omega_{A/R}^\bullet = R\Hom(A, P)[n]$. Il faut montrer que $R\Hom(A, P)_\mathfrak q$
a une cohomologie nulle en degrés $< n - d$.
D'après le lemme \ref{lemma-RHom-ext}, cela revient à montrer que $\Ext_P^i(P/I, P)_{\mathfrak r} = 0$ pour $i < n - d$,
où $\mathfrak r \subset P$ est l'idéal premier correspondant à $\mathfrak q$
et $I$ le noyau de $P \to A$.
On peut récrire le groupe considéré sous la forme $\Ext_{P_\mathfrak r}^i(P_\mathfrak r/IP_\mathfrak r, P_\mathfrak r)$
d'après Compléments d'algèbre, lemme \ref{more-algebra-lemma-pseudo-coherence-and-base-change-ext}.
Il faut donc montrer
$$
\text{profondeur}_{IP_\mathfrak r}(P_\mathfrak r) \geq n - d
$$
d'après le lemme \ref{lemma-depth}. Le lemme \ref{lemma-depth-flat-CM} donne
$$
\text{profondeur}_{IP_\mathfrak r}(P_\mathfrak r) \geq
\dim((P \otimes_R \kappa(\mathfrak p))_\mathfrak r) -
\dim((P/I \otimes_R \kappa(\mathfrak p))_\mathfrak r)
$$
Les membres de droite des deux inégalités coïncident
d'après Algèbre, lemme \ref{algebra-lemma-codimension}.
\end{proof}

\begin{lemma}
\label{lemma-relative-dualizing-flat-vanishing}
Soit $R \to A$ un homomorphisme plat de type fini d'anneaux noethériens.
Soit $\mathfrak q \subset A$ un idéal premier au-dessus de $\mathfrak p \subset R$. Alors
$$
H^i(\omega_{A/R}^\bullet)_\mathfrak q \not = 0
\Rightarrow - d \leq i \leq 0
$$
où $d$ est la dimension de la fibre de $\Spec(A) \to \Spec(R)$ au-dessus de $\mathfrak p$ au point $\mathfrak q$.
Si toutes les fibres de $\Spec(A) \to \Spec(R)$ sont de dimension $\leq d$,
alors $\omega_{A/R}^\bullet$ a son amplitude de Tor dans $[-d, 0]$
en tant que complexe de $R$-modules.
\end{lemma}

\begin{proof}
La borne inférieure a été établie dans le lemme \ref{lemma-relative-dualizing-trivial-vanishing}.
Choisissons une factorisation $R \to P \to A$ avec $P = R[x_1, \ldots, x_n]$ comme dans la section \ref{section-relative-dualizing-complex-algebraic},
de sorte que $\omega_{A/R}^\bullet = R\Hom(A, P)[n]$.
La borne supérieure signifie que $\Ext^i_P(A, P)$ est nul pour $i > n$.
Cela résulte de Compléments d'algèbre, lemme \ref{more-algebra-lemma-perfect-over-polynomial-ring},
qui montre que $A$ est un $P$-module parfait
d'amplitude de Tor dans $[-n, 0]$.

\medskip\noindent
Preuve de la dernière assertion. Soit $R \to R'$ un homomorphisme d'anneaux noethériens.
Posons $A' = A \otimes_R R'$. Alors
$$
\omega_{A'/R'}^\bullet =
\omega_{A/R}^\bullet \otimes_A^\mathbf{L} A' =
\omega_{A/R}^\bullet \otimes_R^\mathbf{L} R'
$$
Le premier isomorphisme résulte du lemme \ref{lemma-base-change-relative-algebraic} et le second,
dans $D(R')$, de Compléments d'algèbre, lemme \ref{more-algebra-lemma-base-change-comparison}.
D'après la première partie de la démonstration
(noter que les fibres de $\Spec(A') \to \Spec(R')$ sont de dimension $\leq d$),
on en déduit que $\omega_{A/R}^\bullet \otimes_R^\mathbf{L} R'$ n'a de cohomologie qu'en degrés $[-d, 0]$.
En prenant pour $R' = R \oplus M$ l'épaississement de carré nul de $R$
par un $R$-module de type fini $M$, on voit que $R\Hom(A, P) \otimes_R^\mathbf{L} M$
n'a de cohomologie que dans l'intervalle $[-d, 0]$
pour tout $R$-module de type fini $M$.
Puisque tout $R$-module est limite inductive filtrante de $R$-modules
de type fini et que les produits tensoriels commutent aux limites inductives,
on conclut.
\end{proof}

\begin{lemma}
\label{lemma-relative-dualizing-CM-vanishing}
Soit $R \to A$ un homomorphisme de type fini d'anneaux noethériens.
Soit $\mathfrak p \subset R$ un idéal premier. Supposons que
\begin{enumerate}
\item $R_\mathfrak p$ soit de Cohen-Macaulay ;
\item pour tout idéal premier minimal $\mathfrak q \subset A$, on ait $\text{trdeg}_{\kappa(R \cap \mathfrak q)} \kappa(\mathfrak q) \leq r$.
\end{enumerate}
Alors
$$
H^i(\omega_{A/R}^\bullet)_\mathfrak p \not = 0 \Rightarrow - r \leq i
$$
et $H^{-r}(\omega_{A/R}^\bullet)_\mathfrak p$ satisfait à $(S_2)$ en tant que $A_\mathfrak p$-module.
\end{lemma}

\begin{proof}
On peut remplacer $R$ par $R_\mathfrak p$ d'après le lemme \ref{lemma-base-change-relative-algebraic}.
On peut donc supposer que $R$ est un anneau local de Cohen-Macaulay,
et il faut établir les assertions du lemme pour les $A$-modules $H^i(\omega_{A/R}^\bullet)$.

\medskip\noindent
Soit $R^\wedge$ le complété de $R$.
L'homomorphisme $R \to R^\wedge$ est plat et $R^\wedge$ est de Cohen-Macaulay
(Compléments d'algèbre, lemme \ref{more-algebra-lemma-completion-CM}).
Observons que les idéaux premiers minimaux de $A \otimes_R R^\wedge$
sont au-dessus d'idéaux premiers minimaux de $A$
par platitude de $A \to A \otimes_R R^\wedge$ (et par le théorème de descente pour les morphismes plats,
voir Algèbre, lemme \ref{algebra-lemma-flat-going-down}).
La condition (2) vaut donc pour l'homomorphisme de type fini $R^\wedge \to A \otimes_R R^\wedge$
d'après Morphismes, lemme \ref{morphisms-lemma-dimension-fibre-after-base-change}.
En appliquant encore le lemme \ref{lemma-base-change-relative-algebraic}, il suffit de prouver le lemme pour $R^\wedge \to A \otimes_R R^\wedge$.
De cette manière, en utilisant le lemme \ref{lemma-ubiquity-dualizing},
on peut supposer que $R$ est un anneau local noethérien de Cohen-Macaulay
possédant un complexe dualisant $\omega_R^\bullet$.

\medskip\noindent
Soit $\mathfrak m \subset A$ un idéal maximal.
Il suffit de montrer les assertions du lemme pour $H^i(\omega_{A/R}^\bullet)_\mathfrak m$.
Si $\mathfrak m$ n'est pas au-dessus de l'idéal maximal de $R$,
on remplace $R$ par un localisé pour se ramener à ce cas
(nous omettons un petit détail).

\medskip\noindent
On peut supposer $\omega_R^\bullet$ normalisé. En posant $d = \dim(R)$,
on voit que $\omega_R^\bullet = \omega_R[d]$ pour un certain $R$-module $\omega_R$ ; voir le lemme \ref{lemma-apply-CM}.
Posons $\omega_A^\bullet = \varphi^!(\omega_R^\bullet)$. D'après le lemme \ref{lemma-relative-dualizing-if-have-omega}, on a
$$
\omega_{A/R}^\bullet =
R\Hom_A(\omega_R[d] \otimes_R^\mathbf{L} A, \omega_A^\bullet)
$$
La formule de dimension donne $\dim(A_\mathfrak m) \leq d + r$ ; voir Morphismes, lemme \ref{morphisms-lemma-dimension-formula-general},
et utiliser le fait que $\kappa(\mathfrak m)$ est fini sur le corps résiduel de $R$
par le Nullstellensatz de Hilbert.
D'après le lemme \ref{lemma-shriek-normalized}, $(\omega_A^\bullet)_\mathfrak m$ est un complexe dualisant normalisé pour $A_\mathfrak m$.
Ainsi $H^i((\omega_A^\bullet)_\mathfrak m)$ n'est non nul que pour $-d - r \leq i \leq 0$ ; voir le lemme \ref{lemma-sitting-in-degrees}.
Puisque $\omega_R[d] \otimes_R^\mathbf{L} A$ est placé en degrés $\leq -d$, on obtient l'annulation voulue.
Enfin, on voit aussi que
$$
H^{-r}(\omega_{A/R}^\bullet)_\mathfrak m =
\Hom_A(\omega_R \otimes_R A, H^{-d - r}(\omega_A^\bullet))_\mathfrak m
$$
Puisque $H^{-d - r}(\omega_A^\bullet)_\mathfrak m$ satisfait à $(S_2)$ d'après le lemme \ref{lemma-depth-dualizing-module},
la dernière assertion résulte de Compléments d'algèbre, lemme \ref{more-algebra-lemma-hom-into-S2}.
\end{proof}



\section{Compléments sur les complexes dualisants}
\label{section-more-dualizing}

\noindent
Quelques lemmes qui ne trouvent pas naturellement leur place ailleurs.

\begin{lemma}
\label{lemma-descent}
Soit $A \to B$ un homomorphisme fidèlement plat d'anneaux noethériens.
Si $K \in D(A)$ et si $K \otimes_A^\mathbf{L} B$ est un complexe dualisant pour $B$,
alors $K$ est un complexe dualisant pour $A$.
\end{lemma}

\begin{proof}
Puisque $A \to B$ est plat, on a $H^i(K) \otimes_A B = H^i(K \otimes_A^\mathbf{L} B)$.
Puisque $K \otimes_A^\mathbf{L} B$ appartient à $D^b_{\textit{Coh}}(B)$, on trouve d'abord que $K$ appartient à $D^b(A)$,
puis que $H^i(K)$ est un $A$-module de type fini d'après Algèbre, lemme \ref{algebra-lemma-descend-properties-modules}.
Soit $M$ un $A$-module de type fini. Alors
$$
R\Hom_A(M, K) \otimes_A B = R\Hom_B(M \otimes_A B, K \otimes_A^\mathbf{L} B)
$$
d'après Compléments d'algèbre, lemme \ref{more-algebra-lemma-base-change-RHom}.
Puisque $K \otimes_A^\mathbf{L} B$ est de dimension injective finie,
disons d'amplitude injective dans $[a, b]$,
le membre de droite a une cohomologie nulle en degrés $> b$.
Puisque $A \to B$ est fidèlement plat, $R\Hom_A(M, K)$ a une cohomologie nulle en degrés $> b$.
Ainsi $K$ est de dimension injective finie
d'après Compléments d'algèbre, lemme \ref{more-algebra-lemma-injective-amplitude}.
Pour conclure, il faut montrer que le morphisme $A \to R\Hom_A(K, K)$ est un isomorphisme.
Utilisons encore Compléments d'algèbre, lemme \ref{more-algebra-lemma-base-change-RHom},
et le fait que
$B \to R\Hom_B(K \otimes_A^\mathbf{L} B, K \otimes_A^\mathbf{L} B)$
est un isomorphisme.
\end{proof}

\begin{lemma}
\label{lemma-descent-ascent}
Soit $\varphi : A \to B$ un homomorphisme d'anneaux noethériens. Supposons que
\begin{enumerate}
\item $A \to B$ soit syntomique et induise une application surjective sur les spectres, ou
\item $A \to B$ soit fidèlement plat et localement d'intersection complète, ou
\item $A \to B$ soit fidèlement plat de type fini, à fibres de Gorenstein.
\end{enumerate}
Alors $K \in D(A)$ est un complexe dualisant pour $A$ si et seulement si
$K \otimes_A^\mathbf{L} B$ est un complexe dualisant pour $B$.
\end{lemma}

\begin{proof}
Observons que $A \to B$ satisfait à (1) si et seulement si $A \to B$ satisfait à (2)
d'après Compléments d'algèbre, lemme \ref{more-algebra-lemma-syntomic-lci}.
Dans les cas (2) et (3), le complexe dualisant relatif $\varphi^!(A) = \omega_{B/A}^\bullet$
est un objet inversible de $D(B)$ ; voir les lemmes \ref{lemma-lci-shriek} et \ref{lemma-gorenstein-shriek}.
De plus, on a $\varphi^!(K) = K \otimes_A^\mathbf{L} \omega_{B/A}^\bullet$ dans les deux cas ; voir le lemme \ref{lemma-upper-shriek-is-tensor-functor} pour le cas (3).
Ainsi $\varphi^!(K)$ est identique à $K \otimes_A^\mathbf{L} B$ à tensorisation près
par un objet inversible de $D(B)$.
Par conséquent, $\varphi^!(K)$ est un complexe dualisant pour $B$
si et seulement si $K \otimes_A^\mathbf{L} B$ l'est
(la propriété d'être dualisant étant locale et invariante par décalage).
Si $K$ est dualisant pour $A$, alors $K \otimes_A^\mathbf{L} B$ est donc dualisant
pour $B$ d'après le lemme \ref{lemma-shriek-dualizing-algebraic}.
Pour la descente de cette propriété, voir le lemme \ref{lemma-descent}.
\end{proof}

\begin{lemma}
\label{lemma-injective-hull-goes-up}
Soit $(A, \mathfrak m, \kappa) \to (B, \mathfrak n, l)$ un homomorphisme local plat d'anneaux noethériens tel que $\mathfrak n = \mathfrak m B$.
Si $E$ est l'enveloppe injective de $\kappa$,
alors $E \otimes_A B$ est l'enveloppe injective de $l$.
\end{lemma}

\begin{proof}
Écrivons $E = \bigcup E_n$ comme dans le lemme \ref{lemma-union-artinian}.
Il suffit de montrer que $E_n \otimes_{A/\mathfrak m^n} B/\mathfrak n^n$ est l'enveloppe injective de $l$ sur $B/\mathfrak n$.
On se ramène ainsi au cas où $A$ et $B$ sont locaux artiniens.
Observons que $\text{longueur}_A(A) = \text{longueur}_B(B)$ et $\text{longueur}_A(E) = \text{longueur}_B(E \otimes_A B)$ d'après Algèbre, lemme \ref{algebra-lemma-pullback-module}.
D'après le lemme \ref{lemma-finite}, on a $\text{longueur}_A(E) = \text{longueur}_A(A)$ et $\text{longueur}_B(E') = \text{longueur}_B(B)$,
où $E'$ est l'enveloppe injective de $l$ sur $B$.
On en déduit que $\text{longueur}_B(E') = \text{longueur}_B(E \otimes_A B)$. Observons que
$$
\dim_l((E \otimes_A B)[\mathfrak n]) =
\dim_l(E[\mathfrak m] \otimes_A B) =
\dim_\kappa(E[\mathfrak m]) = 1
$$
où l'on a utilisé la platitude de $A \to B$ et l'égalité $\mathfrak n = \mathfrak mB$.
Il existe donc un homomorphisme injectif de $B$-modules $E \otimes_A B \to E'$
d'après le lemme \ref{lemma-torsion-submodule-sum-injective-hulls}.
C'est un isomorphisme par l'égalité des longueurs établie ci-dessus.
\end{proof}

\begin{lemma}
\label{lemma-injective-goes-up}
Soit $\varphi : A \to B$ un homomorphisme plat d'anneaux noethériens tel que,
pour tout idéal premier $\mathfrak q \subset B$, on ait $\mathfrak p B_\mathfrak q = \mathfrak qB_\mathfrak q$,
où $\mathfrak p = \varphi^{-1}(\mathfrak q)$, par exemple si $\varphi$ est \'etale.
Si $I$ est un $A$-module injectif, alors $I \otimes_A B$ est un $B$-module injectif.
\end{lemma}

\begin{proof}
Les morphismes \'etales satisfont à l'hypothèse d'après Algèbre, lemme \ref{algebra-lemma-etale-at-prime}.
Le lemme \ref{lemma-sum-injective-modules} et la proposition \ref{proposition-structure-injectives-noetherian} permettent de supposer que $I$
est l'enveloppe injective de $\kappa(\mathfrak p)$ pour un idéal premier $\mathfrak p \subset A$.
Alors $I$ est un module sur $A_\mathfrak p$.
Il suffit de prouver que $I \otimes_A B = I \otimes_{A_\mathfrak p} B_\mathfrak p$ est injectif comme $B_\mathfrak p$-module ; voir le lemme \ref{lemma-injective-flat}.
On peut donc supposer $(A, \mathfrak m, \kappa)$ local noethérien
et $I = E$ égal à l'enveloppe injective du corps résiduel $\kappa$.
Notre hypothèse entraîne que l'anneau noethérien $B/\mathfrak m B$
est produit de corps (nous omettons les détails).
Il n'y a donc qu'un nombre fini d'idéaux premiers $\mathfrak m_1, \ldots, \mathfrak m_n$ dans $B$
au-dessus de $\mathfrak m$, et ils sont tous maximaux.
Écrivons $E = \bigcup E_n$ comme dans le lemme \ref{lemma-union-artinian}.
Alors $E \otimes_A B = \bigcup E_n \otimes_A B$, et $E_n \otimes_A B$ est un $B$-module de type fini
de support $\{\mathfrak m_1, \ldots, \mathfrak m_n\}$, donc se décompose en produit des localisés en les $\mathfrak m_i$.
Ainsi $E \otimes_A B = \prod (E \otimes_A B)_{\mathfrak m_i}$.
Puisque $(E \otimes_A B)_{\mathfrak m_i} = E \otimes_A B_{\mathfrak m_i}$ est l'enveloppe injective du corps résiduel de $\mathfrak m_i$
d'après le lemme \ref{lemma-injective-hull-goes-up}, on conclut.
\end{proof}






\section{Complexes dualisants relatifs}
\label{section-relative-dualizing-complexes}

\noindent
Pour un homomorphisme de type fini $\varphi : R \to A$ d'anneaux noethériens,
on dispose du complexe dualisant relatif $\omega_{A/R}^\bullet = \varphi^!(R)$ étudié dans la section \ref{section-relative-dualizing-complexes-Noetherian}.
Si $R$ n'est pas noethérien, un complexe construit de la même manière
n'aura généralement pas de bonnes propriétés.
Nous donnons ici une définition du complexe dualisant relatif
pour les homomorphismes plats de présentation finie $R \to A$
d'anneaux non noethériens. Cette définition est choisie pour se globaliser
aux morphismes plats de présentation finie de schémas ;
voir Dualité pour les schémas, section \ref{duality-section-relative-dualizing-complexes}.
Nous montrerons que les complexes dualisants relatifs existent
dans le cadre de la définition, sont uniques à isomorphisme non canonique près,
et que l'on retrouve dans le cas noethérien le complexe de la section \ref{section-relative-dualizing-complexes-Noetherian}.

\medskip\noindent
Le lecteur qui s'en tient au cas noethérien peut sans risque sauter cette section !

\begin{definition}
\label{definition-relative-dualizing-complex}
Soit $R \to A$ un homomorphisme plat de présentation finie.
Un {\it complexe dualisant relatif} est un objet $K \in D(A)$ tel que
\begin{enumerate}
\item $K$ soit $R$-parfait (Compléments d'algèbre, définition \ref{more-algebra-definition-relatively-perfect}) ;
\item $R\Hom_{A \otimes_R A}(A, K \otimes_A^\mathbf{L} (A \otimes_R A))$ soit isomorphe à $A$.
\end{enumerate}
\end{definition}

\noindent
Pour comprendre cette définition, il peut être nécessaire de lire
et de comprendre certains des lemmes suivants.
Les lemmes \ref{lemma-relative-dualizing-noetherian} et \ref{lemma-uniqueness-relative-dualizing} montrent qu'elle est compatible
avec la définition de la section \ref{section-relative-dualizing-complexes-Noetherian}.

\begin{lemma}
\label{lemma-uniqueness-relative-dualizing}
Soit $R \to A$ un homomorphisme plat de présentation finie.
Deux complexes dualisants relatifs pour $R \to A$ sont isomorphes.
\end{lemma}

\begin{proof}
Soient $K$ et $L$ deux complexes dualisants relatifs pour $R \to A$.
Notons $K_1 = K \otimes_A^\mathbf{L} (A \otimes_R A)$ et $L_2 = (A \otimes_R A) \otimes_A^\mathbf{L} L$ les changements de base dérivés
par les première et seconde coprojections $A \to A \otimes_R A$.
Par symétrie, l'hypothèse sur $L_2$ entraîne que $R\Hom_{A \otimes_R A}(A, L_2)$ est isomorphe à $A$.
En appliquant deux fois Compléments d'algèbre, lemme \ref{more-algebra-lemma-internal-hom-evaluate-tensor-isomorphism} (3), on a
$$
A \otimes_{A \otimes_R A}^\mathbf{L} L_2 \cong
R\Hom_{A \otimes_R A}(A, K_1 \otimes_{A \otimes_R A}^\mathbf{L} L_2) \cong
A \otimes_{A \otimes_R A}^\mathbf{L} K_1
$$
L'application du foncteur de restriction $D(A \otimes_R A) \to D(A)$
pour l'une ou l'autre coprojection donne le résultat voulu.
\end{proof}

\begin{lemma}
\label{lemma-relative-dualizing-noetherian}
Soit $\varphi : R \to A$ un homomorphisme plat de type fini d'anneaux noethériens.
Alors le complexe dualisant relatif $\omega_{A/R}^\bullet = \varphi^!(R)$ de la section \ref{section-relative-dualizing-complexes-Noetherian}
est un complexe dualisant relatif au sens de la définition \ref{definition-relative-dualizing-complex}.
\end{lemma}

\begin{proof}
Le lemme \ref{lemma-relative-dualizing-algebraic} montre que $\varphi^!(R)$ est $R$-parfait.
Notons $\delta : A \otimes_R A \to A$ l'homomorphisme de multiplication et $p_1, p_2 : A \to A \otimes_R A$ les coprojections.
Alors
$$
\varphi^!(R) \otimes_A^\mathbf{L} (A \otimes_R A) =
\varphi^!(R) \otimes_{A, p_1}^\mathbf{L} (A \otimes_R A) =
p_2^!(A)
$$
d'après le lemme \ref{lemma-flat-bc}. Rappelons que
$
R\Hom_{A \otimes_R A}(A, \varphi^!(R) \otimes_A^\mathbf{L} (A \otimes_R A))
$
est l'image de $\delta^!(\varphi^!(R) \otimes_A^\mathbf{L} (A \otimes_R A))$ par le foncteur de restriction $\delta_* : D(A) \to D(A \otimes_R A)$.
Utiliser la définition de $\delta^!$ dans la section \ref{section-relative-dualizing-complex-algebraic}
et le lemme \ref{lemma-RHom-ext}.
Puisque $\delta^!(p_2^!(A)) \cong A$ d'après le lemme \ref{lemma-composition-shriek-algebraic}, on conclut.
\end{proof}

\begin{lemma}
\label{lemma-base-change-relative-dualizing}
Soit $R \to A$ un homomorphisme plat de présentation finie. Alors :
\begin{enumerate}
\item il existe un complexe dualisant relatif $K$ dans $D(A)$ ;
\item pour tout homomorphisme $R \to R'$, si l'on pose $A' = A \otimes_R R'$
et $K' = K \otimes_A^\mathbf{L} A'$, alors $K'$ est un complexe dualisant relatif pour $R' \to A'$.
\end{enumerate}
De plus, si
$$
\xi : A \longrightarrow K \otimes_A^\mathbf{L} (A \otimes_R A)
$$
est un générateur du module cyclique
$\Hom_{D(A \otimes_R A)}(A, K \otimes_A^\mathbf{L} (A \otimes_R A))$
alors, dans (2), le changement de base dérivé de $\xi$ par
$A \otimes_R A \to A' \otimes_{R'} A'$ est un générateur du module cyclique
$\Hom_{D(A' \otimes_{R'} A')}(A',
K' \otimes_{A'}^\mathbf{L} (A' \otimes_{R'} A'))$
\end{lemma}

\begin{proof}
Ramenons-nous d'abord au cas noethérien.
D'après Algèbre, lemme \ref{algebra-lemma-flat-finite-presentation-limit-flat}, il existe une sous-$\mathbf{Z}$-algèbre de type fini
$R_0 \subset R$ et un homomorphisme plat de type fini $R_0 \to A_0$ tels que $A = A_0 \otimes_{R_0} R$.
D'après le lemme \ref{lemma-relative-dualizing-noetherian}, il existe un complexe dualisant relatif $K_0 \in D(A_0)$.
Si l'on montre (2) pour $K_0$, alors $K_0 \otimes_{A_0}^\mathbf{L} A$ est un complexe dualisant pour $R \to A$
et satisfait lui aussi à (2) par transitivité du changement de base dérivé.
L'unicité des complexes dualisants relatifs (lemme \ref{lemma-uniqueness-relative-dualizing})
montre alors que cela vaut pour tout complexe dualisant relatif.

\medskip\noindent
Supposons $R$ noethérien et soit $K$ un complexe dualisant relatif pour $R \to A$.
Étant donné un homomorphisme $R \to R'$, posons $A' = A \otimes_R R'$ et $K' = K \otimes_A^\mathbf{L} A'$.
Pour conclure, il faut montrer que $K'$ est un complexe dualisant relatif pour $R' \to A'$.
D'après Compléments d'algèbre, lemme \ref{more-algebra-lemma-base-change-relatively-perfect},
on voit que $K'$ est $R'$-parfait dans tous les cas.
D'après les lemmes \ref{lemma-base-change-relative-algebraic} et \ref{lemma-relative-dualizing-noetherian}, si $R'$ est noethérien,
alors $K'$ est un complexe dualisant relatif pour $R' \to A'$ dans les deux sens.
La transitivité du produit tensoriel dérivé donne $K \otimes_A^\mathbf{L} (A \otimes_R A)
\otimes_{A \otimes_R A}^\mathbf{L} (A' \otimes_{R'} A') =
K' \otimes_{A'}^\mathbf{L} (A' \otimes_{R'} A')$.
La platitude de $R \to A$ assure que $A \otimes_{A \otimes_R A}^\mathbf{L} (A' \otimes_{R'} A') = A'$ ;
en effet, $A \otimes_R A$ et $R'$ sont Tor-indépendants sur $R$,
ce qui permet d'appliquer Compléments d'algèbre, lemme \ref{more-algebra-lemma-base-change-comparison}.
Enfin, $A$ est pseudo-cohérent comme $A \otimes_R A$-module
d'après Compléments d'algèbre, lemme \ref{more-algebra-lemma-more-relative-pseudo-coherent-is-moot}.
Toutes les hypothèses de Compléments d'algèbre, lemme \ref{more-algebra-lemma-compute-RHom-relatively-perfect},
sont donc vérifiées.
On obtient un complexe borné inférieurement $E^\bullet$ de modules
plats sur $R$ et de présentation finie sur $A \otimes_R A$
tel que $E^\bullet \otimes_R R'$ représente $R\Hom_{A' \otimes_{R'} A'}(A',
K' \otimes_{A'}^\mathbf{L} (A' \otimes_{R'} A'))$,
ces identifications étant compatibles au changement de base dérivé.
Soit $n \in \mathbf{Z}$ tel que $n \not = 0$. Définissons $Q^n$ par la suite
$$
E^{n - 1} \to E^n \to Q^n \to 0
$$
Puisque $\kappa(\mathfrak p)$ est un anneau noethérien, on sait que $H^n(E^\bullet \otimes_R \kappa(\mathfrak p)) = 0$,
voir les remarques ci-dessus. Une chasse au diagramme montre alors que
$$
Q^n \otimes_R \kappa(\mathfrak p) \to E^{n + 1} \otimes_R \kappa(\mathfrak p)
$$
est injectif. Ainsi, pour un idéal premier $\mathfrak q$ de $A \otimes_R A$
au-dessus de $\mathfrak p$, le module $Q^n_\mathfrak q$ est $R_\mathfrak p$-plat
et $Q^n_\mathfrak p \to E^{n + 1}_\mathfrak q$ est $R_\mathfrak p$-universellement injectif ; voir Algèbre, lemme \ref{algebra-lemma-mod-injective}.
Puisque cela vaut pour tous les idéaux premiers,
$Q^n$ est $R$-plat et $Q^n \to E^{n + 1}$ est $R$-universellement injectif.
En particulier, $H^n(E^\bullet \otimes_R R') = 0$ pour tout homomorphisme $R \to R'$.
Soit $Z^0 = \Ker(E^0 \to E^1)$. La suite exacte $0 \to Z^0 \to E^0 \to E^1 \to Q^1 \to 0$ montre que $Z^0$ est $R$-plat et que
$Z^0 \otimes_R R' = \Ker(E^0 \otimes_R R' \to E^1 \otimes_R R')$
pour tout $R \to R'$. La suite exacte courte
$0 \to Q^{-1} \to Z^0 \to H^0(E^\bullet) \to 0$
donne alors
$$
H^0(E^\bullet \otimes_R R') = H^0(E^\bullet) \otimes_R R'
= A \otimes_R R' = A'
$$
comme voulu. Cette égalité donne aussi la dernière assertion du lemme.
\end{proof}

\begin{lemma}
\label{lemma-relative-dualizing-RHom}
Soit $R \to A$ un homomorphisme plat de présentation finie.
Soit $K$ un complexe dualisant relatif.
Alors $A \to R\Hom_A(K, K)$ est un isomorphisme.
\end{lemma}

\begin{proof}
D'après Algèbre, lemme \ref{algebra-lemma-flat-finite-presentation-limit-flat}, il existe une sous-$\mathbf{Z}$-algèbre de type fini
$R_0 \subset R$ et un homomorphisme plat de type fini $R_0 \to A_0$ tels que $A = A_0 \otimes_{R_0} R$.
D'après les lemmes \ref{lemma-uniqueness-relative-dualizing}, \ref{lemma-relative-dualizing-noetherian} et \ref{lemma-base-change-relative-dualizing},
il existe un complexe dualisant relatif $K_0 \in D(A_0)$
dont le changement de base dérivé est $K$.
On se ramène ainsi à la situation du paragraphe suivant.

\medskip\noindent
Supposons $R$ noethérien et soit $K$ un complexe dualisant relatif pour $R \to A$.
Étant donné un homomorphisme $R \to R'$, posons $A' = A \otimes_R R'$ et $K' = K \otimes_A^\mathbf{L} A'$.
Pour conclure, montrons que $R\Hom_{A'}(K', K') = A'$.
D'après le lemme \ref{lemma-relative-dualizing-algebraic}, c'est vrai dès que $R'$ est noethérien.
Puisqu'un $R'$ quelconque est limite inductive filtrante
de $R$-algèbres noethériennes, le résultat découle de
Compléments d'algèbre, lemme \ref{more-algebra-lemma-colimit-relatively-perfect}.
\end{proof}

\begin{lemma}
\label{lemma-relative-dualizing-composition}
Soient $R \to A \to B$ des homomorphismes plats de présentation finie.
Soient $K_{A/R}$ et $K_{B/A}$ des complexes dualisants relatifs pour $R \to A$ et $A \to B$.
Alors $K = K_{A/R} \otimes_A^\mathbf{L} K_{B/A}$ est un complexe dualisant relatif pour $R \to B$.
\end{lemma}

\begin{proof}
Utilisons la réduction au cas noethérien.
D'après Algèbre, lemme \ref{algebra-lemma-flat-finite-presentation-limit-flat}, il existe une sous-$\mathbf{Z}$-algèbre de type fini
$R_0 \subset R$ et un homomorphisme plat de type fini $R_0 \to A_0$ tels que $A = A_0 \otimes_{R_0} R$.
Quitte à agrandir $R_0$ et à remplacer $A_0$ en conséquence,
on peut supposer qu'il existe un homomorphisme plat de type fini $A_0 \to B_0$
tel que $B = B_0 \otimes_{R_0} R$ (utiliser le même lemme).
Si l'on prouve le lemme pour $R_0 \to A_0 \to B_0$, il en résulte
par les lemmes \ref{lemma-uniqueness-relative-dualizing}, \ref{lemma-relative-dualizing-noetherian} et \ref{lemma-base-change-relative-dualizing}.
On se ramène ainsi à la situation du paragraphe suivant.

\medskip\noindent
Supposons $R$ noethérien et notons $\varphi : R \to A$ et $\psi : A \to B$
les homomorphismes donnés. Alors $K_{A/R} \cong \varphi^!(R)$ et $K_{B/A} \cong \psi^!(A)$,
voir les références ci-dessus.
On a alors
$$
K = K_{A/R} \otimes_A^\mathbf{L} K_{B/A} \cong
\varphi^!(R) \otimes_A^\mathbf{L} \psi^!(A) \cong
\psi^!(\varphi^!(R)) \cong (\psi \circ \varphi)^!(R)
$$
d'après les lemmes \ref{lemma-upper-shriek-is-tensor-functor} et \ref{lemma-composition-shriek-algebraic}.
Ainsi $K$ est un complexe dualisant relatif pour $R \to B$.
\end{proof}





\begin{multicols}{2}[\section{Autres chapitres}]
\noindent
Préliminaires
\begin{enumerate}
\item \hyperref[introduction-section-phantom]{Introduction}
\item \hyperref[conventions-section-phantom]{Conventions}
\item \hyperref[sets-section-phantom]{Théorie des ensembles}
\item \hyperref[categories-section-phantom]{Catégories}
\item \hyperref[topology-section-phantom]{Topologie}
\item \hyperref[sheaves-section-phantom]{Faisceaux sur les espaces}
\item \hyperref[sites-section-phantom]{Sites et faisceaux}
\item \hyperref[stacks-section-phantom]{Champs}
\item \hyperref[fields-section-phantom]{Corps}
\item \hyperref[algebra-section-phantom]{Algèbre commutative}
\item \hyperref[brauer-section-phantom]{Groupes de Brauer}
\item \hyperref[homology-section-phantom]{Algèbre homologique}
\item \hyperref[derived-section-phantom]{Catégories dérivées}
\item \hyperref[simplicial-section-phantom]{Méthodes simpliciales}
\item \hyperref[more-algebra-section-phantom]{Compléments d'algèbre}
\item \hyperref[smoothing-section-phantom]{Lissage des morphismes d'anneaux}
\item \hyperref[modules-section-phantom]{Faisceaux de modules}
\item \hyperref[sites-modules-section-phantom]{Modules sur les sites}
\item \hyperref[injectives-section-phantom]{Objets injectifs}
\item \hyperref[cohomology-section-phantom]{Cohomologie des faisceaux}
\item \hyperref[sites-cohomology-section-phantom]{Cohomologie sur les sites}
\item \hyperref[dga-section-phantom]{Algèbre différentielle graduée}
\item \hyperref[dpa-section-phantom]{Algèbre à puissances divisées}
\item \hyperref[sdga-section-phantom]{Faisceaux différentiels gradués}
\item \hyperref[hypercovering-section-phantom]{Hyperrecouvrements}
\end{enumerate}
Schémas
\begin{enumerate}
\setcounter{enumi}{25}
\item \hyperref[schemes-section-phantom]{Schémas}
\item \hyperref[constructions-section-phantom]{Constructions de schémas}
\item \hyperref[properties-section-phantom]{Propriétés des schémas}
\item \hyperref[morphisms-section-phantom]{Morphismes de schémas}
\item \hyperref[coherent-section-phantom]{Cohomologie des schémas}
\item \hyperref[divisors-section-phantom]{Diviseurs}
\item \hyperref[limits-section-phantom]{Limites de schémas}
\item \hyperref[varieties-section-phantom]{Variétés}
\item \hyperref[topologies-section-phantom]{Topologies sur les schémas}
\item \hyperref[descent-section-phantom]{Descente}
\item \hyperref[perfect-section-phantom]{Catégories dérivées de schémas}
\item \hyperref[more-morphisms-section-phantom]{Compléments sur les morphismes}
\item \hyperref[flat-section-phantom]{Compléments sur la platitude}
\item \hyperref[groupoids-section-phantom]{Schémas en groupoïdes}
\item \hyperref[more-groupoids-section-phantom]{Compléments sur les schémas en groupoïdes}
\item \hyperref[etale-section-phantom]{Morphismes étales de schémas}
\end{enumerate}
Thèmes de théorie des schémas
\begin{enumerate}
\setcounter{enumi}{41}
\item \hyperref[chow-section-phantom]{Homologie de Chow}
\item \hyperref[intersection-section-phantom]{Théorie de l'intersection}
\item \hyperref[pic-section-phantom]{Schémas de Picard des courbes}
\item \hyperref[weil-section-phantom]{Théories cohomologiques de Weil}
\item \hyperref[adequate-section-phantom]{Modules adéquats}
\item \hyperref[dualizing-section-phantom]{Complexes dualisants}
\item \hyperref[duality-section-phantom]{Dualité pour les schémas}
\item \hyperref[discriminant-section-phantom]{Discriminants et différentes}
\item \hyperref[derham-section-phantom]{Cohomologie de de Rham}
\item \hyperref[local-cohomology-section-phantom]{Cohomologie locale}
\item \hyperref[algebraization-section-phantom]{Géométrie algébrique et formelle}
\item \hyperref[curves-section-phantom]{Courbes algébriques}
\item \hyperref[resolve-section-phantom]{Résolution des surfaces}
\item \hyperref[models-section-phantom]{Réduction semi-stable}
\item \hyperref[functors-section-phantom]{Foncteurs et morphismes}
\item \hyperref[equiv-section-phantom]{Catégories dérivées de variétés}
\item \hyperref[pione-section-phantom]{Groupes fondamentaux de schémas}
\item \hyperref[etale-cohomology-section-phantom]{Cohomologie étale}
\item \hyperref[crystalline-section-phantom]{Cohomologie cristalline}
\item \hyperref[proetale-section-phantom]{Cohomologie pro-étale}
\item \hyperref[relative-cycles-section-phantom]{Cycles relatifs}
\item \hyperref[more-etale-section-phantom]{Compléments de cohomologie étale}
\item \hyperref[trace-section-phantom]{La formule des traces}
\end{enumerate}
Espaces algébriques
\begin{enumerate}
\setcounter{enumi}{64}
\item \hyperref[spaces-section-phantom]{Espaces algébriques}
\item \hyperref[spaces-properties-section-phantom]{Propriétés des espaces algébriques}
\item \hyperref[spaces-morphisms-section-phantom]{Morphismes d'espaces algébriques}
\item \hyperref[decent-spaces-section-phantom]{Espaces algébriques décents}
\item \hyperref[spaces-cohomology-section-phantom]{Cohomologie des espaces algébriques}
\item \hyperref[spaces-limits-section-phantom]{Limites d'espaces algébriques}
\item \hyperref[spaces-divisors-section-phantom]{Diviseurs sur les espaces algébriques}
\item \hyperref[spaces-over-fields-section-phantom]{Espaces algébriques sur des corps}
\item \hyperref[spaces-topologies-section-phantom]{Topologies sur les espaces algébriques}
\item \hyperref[spaces-descent-section-phantom]{Descente et espaces algébriques}
\item \hyperref[spaces-perfect-section-phantom]{Catégories dérivées d'espaces}
\item \hyperref[spaces-more-morphisms-section-phantom]{Compléments sur les morphismes d'espaces}
\item \hyperref[spaces-flat-section-phantom]{Platitude sur les espaces algébriques}
\item \hyperref[spaces-groupoids-section-phantom]{Groupoïdes dans les espaces algébriques}
\item \hyperref[spaces-more-groupoids-section-phantom]{Compléments sur les groupoïdes dans les espaces}
\item \hyperref[bootstrap-section-phantom]{Amorçage}
\item \hyperref[spaces-pushouts-section-phantom]{Sommes amalgamées d'espaces algébriques}
\end{enumerate}
Thèmes de géométrie
\begin{enumerate}
\setcounter{enumi}{81}
\item \hyperref[spaces-chow-section-phantom]{Groupes de Chow des espaces}
\item \hyperref[groupoids-quotients-section-phantom]{Quotients de groupoïdes}
\item \hyperref[spaces-more-cohomology-section-phantom]{Compléments sur la cohomologie des espaces}
\item \hyperref[spaces-simplicial-section-phantom]{Espaces simpliciaux}
\item \hyperref[spaces-duality-section-phantom]{Dualité pour les espaces}
\item \hyperref[formal-spaces-section-phantom]{Espaces algébriques formels}
\item \hyperref[restricted-section-phantom]{Algébrisation des espaces formels}
\item \hyperref[spaces-resolve-section-phantom]{Résolution des surfaces revisitée}
\end{enumerate}
Théorie des déformations
\begin{enumerate}
\setcounter{enumi}{89}
\item \hyperref[formal-defos-section-phantom]{Théorie formelle des déformations}
\item \hyperref[defos-section-phantom]{Théorie des déformations}
\item \hyperref[cotangent-section-phantom]{Le complexe cotangent}
\item \hyperref[examples-defos-section-phantom]{Problèmes de déformation}
\end{enumerate}
Champs algébriques
\begin{enumerate}
\setcounter{enumi}{93}
\item \hyperref[algebraic-section-phantom]{Champs algébriques}
\item \hyperref[examples-stacks-section-phantom]{Exemples de champs}
\item \hyperref[stacks-sheaves-section-phantom]{Faisceaux sur les champs algébriques}
\item \hyperref[criteria-section-phantom]{Critères de représentabilité}
\item \hyperref[artin-section-phantom]{Axiomes d'Artin}
\item \hyperref[quot-section-phantom]{Espaces de Quot et de Hilbert}
\item \hyperref[stacks-properties-section-phantom]{Propriétés des champs algébriques}
\item \hyperref[stacks-morphisms-section-phantom]{Morphismes de champs algébriques}
\item \hyperref[stacks-limits-section-phantom]{Limites de champs algébriques}
\item \hyperref[stacks-cohomology-section-phantom]{Cohomologie des champs algébriques}
\item \hyperref[stacks-perfect-section-phantom]{Catégories dérivées de champs}
\item \hyperref[stacks-introduction-section-phantom]{Introduction aux champs algébriques}
\item \hyperref[stacks-more-morphisms-section-phantom]{Compléments sur les morphismes de champs}
\item \hyperref[stacks-geometry-section-phantom]{La géométrie des champs}
\end{enumerate}
Thèmes de théorie des espaces de modules
\begin{enumerate}
\setcounter{enumi}{107}
\item \hyperref[moduli-section-phantom]{Champs de modules}
\item \hyperref[moduli-curves-section-phantom]{Espaces de modules de courbes}
\end{enumerate}
Divers
\begin{enumerate}
\setcounter{enumi}{109}
\item \hyperref[examples-section-phantom]{Exemples}
\item \hyperref[exercises-section-phantom]{Exercices}
\item \hyperref[guide-section-phantom]{Guide bibliographique}
\item \hyperref[desirables-section-phantom]{Desiderata}
\item \hyperref[coding-section-phantom]{Style de programmation}
\item \hyperref[obsolete-section-phantom]{Obsolète}
\item \hyperref[fdl-section-phantom]{Licence de documentation libre GNU}
\item \hyperref[index-section-phantom]{Index généré automatiquement}
\end{enumerate}
\end{multicols}


\bibliography{my}
\bibliographystyle{amsalpha}

\end{document}
